% Qu'est-ce que la philosophie? — Philosophie 1re Tech — Cours signé OnBuch+
% Programme officiel MINESEC. Compiler avec : tectonic N15-qu-est-ce-que-philosophie.tex
\newcommand{\DOCMATIERE}{Philosophie}
\newcommand{\DOCNIVEAU}{1re Tech}
\newcommand{\DOCMODULE}{Philosophie – classe de Première technique}
\newcommand{\DOCLECON}{N15}
\newcommand{\DOCTITRE}{Qu'est-ce que la philosophie?}
% OnBuch+ — préambule « cours signés » ENRICHI (compilé avec tectonic / XeLaTeX)
% Système visuel « Pop » d'OnBuch+ : fond crème (jamais blanc), bordures encre
% épaisses, ombres « dures » décalées, titres Archivo Black en capitales.
% Version MATHÉMATIQUES : graphes (pgfplots/tikz), boîtes pédagogiques
% (définition, propriété, méthode, attention, le savais-tu, exercice résolu,
% exercice, TP, bilan), page de garde et signature OnBuch+ ancrée en bas.
%
% Macros attendues dans main.tex AVANT \input{preamble} :
%   \DOCMATIERE  \DOCNIVEAU  \DOCMODULE  \DOCLECON (numéro)  \DOCTITRE
\documentclass[11pt]{article}

\usepackage{fontspec}
\usepackage[french]{babel}
\usepackage[a4paper,margin=2cm,top=3.4cm,bottom=2.8cm,headheight=1.4cm,footskip=1.3cm]{geometry}
\usepackage{amsmath,amssymb,amsfonts}
\usepackage{mathtools} % \xmapsto, \coloneqq... souvent utilisés par les modèles
\usepackage{cancel} % simplifications barrées (bilans de forces)
% \abs{z} (module, valeur absolue) et \norm{u} (norme), utilisés par les modèles
\providecommand{\abs}{}\let\abs\relax\DeclarePairedDelimiter{\abs}{\lvert}{\rvert}
\providecommand{\norm}{}\let\norm\relax\DeclarePairedDelimiter{\norm}{\lVert}{\rVert}
\usepackage[table]{xcolor} % [table] : fournit aussi \rowcolors (alternance de lignes)
\usepackage{pagecolor}
\usepackage{tikz}
\usetikzlibrary{arrows.meta,positioning,calc,shapes.geometric,shapes.misc,shapes.arrows,decorations.pathmorphing,decorations.pathreplacing,decorations.markings,patterns,shadows,mindmap,trees,fit,backgrounds,matrix,intersections,angles,quotes}
\usepackage{pgfplots}
\usepackage[version=4]{mhchem} % équations nucléaires \ce{^{235}_{92}U}
\usepackage[european,siunitx]{circuitikz} % schémas électriques
\pgfplotsset{compat=1.18}
\usepgfplotslibrary{fillbetween,groupplots} % aires entre courbes (calcul intégral)
\usepackage{enumitem}
\usepackage{fancyhdr}
% [most] charge toutes les bibliothèques tcolorbox (listings, minted, raster,
% documentation...) et gonfle inutilement la mémoire TeX sur un document long
% (~300 boîtes) : on ne charge que ce qui est réellement utilisé.
\usepackage[skins,breakable]{tcolorbox}
\usepackage{graphicx}
\usepackage{colortbl}
\usepackage{booktabs}
\usepackage{tabularx}
\usepackage{multirow}
\usepackage{diagbox} % case d'en-tête diagonale des tableaux à double entrée
% Colonnes extensibles centrée / gauche / droite (C, L, R), souvent utilisées par les modèles
\newcolumntype{C}{>{\centering\arraybackslash}X}
\newcolumntype{L}{>{\raggedright\arraybackslash}X}
\newcolumntype{R}{>{\raggedleft\arraybackslash}X}
\usepackage{array}
\usepackage{multicol}
\usepackage{siunitx}
% parse-numbers=false : accepte aussi \SI{2,5 \times 10^{-3}}{...}
\sisetup{locale=FR,output-decimal-marker={,},group-minimum-digits=4,parse-numbers=false}
\DeclareSIUnit\ohm{\ensuremath{\symup{\Omega}}} % U+2126 absent de la police
\DeclareSIUnit\division{div} % réglages d'oscilloscope (ms/div, V/div)
\DeclareSIUnit\tr{tr} % tours (vitesse de rotation, ex. \SI{1500}{\tr\per\minute})
\DeclareSIUnit\molar{M} % concentration molaire (ex. \SI{3}{\molar}, \SI{1}{\micro\molar})
\DeclareSIUnit\an{an} % année (le modèle confond parfois avec \year, un compteur TeX) — ex. \SI{5}{\kilo\meter\per\an}
\DeclareSIUnit\FCFA{FCFA} % franc CFA (ex. \SI{500}{\FCFA\per\kilo\gram})
\DeclareSIUnit\degreeBrix{\textdegree Brix} % teneur en sucre d'un sirop/jus (réfractométrie)
\DeclareSIUnit\month{mois} % le modèle utilise parfois \month, un compteur TeX, comme unité de durée
\DeclareSIUnit\microliter{\micro\liter} % le modèle écrit parfois \microliter en un seul token
\DeclareSIUnit\million{millions} % dénombrement (ex. \SI{15}{\million\per\mL} spermatozoïdes)
\usepackage{qrcode}
% \degree : commande fréquemment utilisée par les modèles pour un angle ou une
% température en dehors de \SI{}{} (non fournie par les paquets chargés).
\providecommand{\degree}{^{\circ}}
% \qed (fin de démonstration), utilisé par les modèles sans amsthm
\providecommand{\qed}{\hfill$\square$}
% \barre : double barre des valeurs interdites dans un tableau de variations (tkz-tab)
\providecommand{\barre}{\ensuremath{\|}}
% environnement proof (amsthm non chargé) utilisé par les modèles
\ifdefined\proof\else\newenvironment{proof}[1][Démonstration]{\par\noindent\textit{#1.}\ }{\qed\par}\fi
\usepackage{lastpage}
\usepackage{hyperref}
\hypersetup{hidelinks}

% ============================================================
% Palette « Pop » OnBuch+ — mêmes hex que la classe P (plus_ui.dart)
% ============================================================
\definecolor{poporange}{HTML}{F59321}
\definecolor{popdark}{HTML}{E07A0C}
\definecolor{popcream}{HTML}{FFF6E8}
\definecolor{popink}{HTML}{1C1714}
\definecolor{poppurple}{HTML}{7A5AE0}
\definecolor{popgreen}{HTML}{2E9E6A}
\definecolor{poppink}{HTML}{E0457B}
\definecolor{popblue}{HTML}{2F7DE1}
\definecolor{popgold}{HTML}{C98A12}
\definecolor{popmuted}{HTML}{8A7F76}
\definecolor{popline}{HTML}{EADFD2}
\definecolor{popgray}{HTML}{8A7F76}
\colorlet{popgrey}{popgray}
% Alias tolérants (le modèle nomme parfois les couleurs différemment)
\colorlet{popwhite}{white}
\colorlet{popblack}{popink}
\colorlet{popred}{poppink}
\colorlet{popyellow}{popgold}
% Teintes claires (fonds de tableaux/encadrés) — le modèle nomme parfois
% les couleurs avec un suffixe L (ex. popblueL) sans les avoir définies.
\colorlet{poporangeL}{poporange!20}
\colorlet{popdarkL}{popdark!20}
\colorlet{poppurpleL}{poppurple!20}
\colorlet{popgreenL}{popgreen!20}
\colorlet{poppinkL}{poppink!20}
\colorlet{popblueL}{popblue!20}
\colorlet{popgoldL}{popgold!20}
\colorlet{popredL}{popred!20}
\colorlet{popgrayL}{popgray!20}
% Teintes claires pour fonds de tableaux / zones
\colorlet{popblueL}{popblue!10!white}
\colorlet{poporangeL}{poporange!14!white}
\colorlet{popgreenL}{popgreen!12!white}
\colorlet{poppurpleL}{poppurple!12!white}
\colorlet{poppinkL}{poppink!12!white}
\colorlet{popgoldL}{popgold!14!white}

\pagecolor{popcream}

% ============================================================
% Typographie
% ============================================================
\defaultfontfeatures{Ligatures=TeX}
\setmainfont{PlusJakartaSans-Medium.ttf}[Path=../../../fonts/,BoldFont=PlusJakartaSans-Bold.ttf]
% Maths en sans-serif (Fira Math) avec chiffres et lettres droites de Plus
% Jakarta, pour que les formules se fondent dans le texte.
\usepackage[math-style=ISO,bold-style=ISO]{unicode-math}
\setmathfont{FiraMath-Regular.otf}[Path=../../../fonts/]
\setmathfont{PlusJakartaSans-Medium.ttf}[Path=../../../fonts/,range={up/num,up/{Latin,latin},\mathup}]
\usepackage{icomma} % virgule décimale sans espace en mode math
% Caractères mathématiques Unicode absents des polices de texte (Plus Jakarta,
% Archivo Black) : rendus par leur équivalent mathématique, en texte comme en
% formule — sinon ils s'affichent en carré vide (ex. « Arithmétique dans ℤ »).
\usepackage{newunicodechar}
\newunicodechar{ℝ}{\ensuremath{\mathbb{R}}}
\newunicodechar{ℤ}{\ensuremath{\mathbb{Z}}}
\newunicodechar{ℂ}{\ensuremath{\mathbb{C}}}
\newunicodechar{ℕ}{\ensuremath{\mathbb{N}}}
\newunicodechar{ℚ}{\ensuremath{\mathbb{Q}}}
\newunicodechar{𝔻}{\ensuremath{\mathbb{D}}}
\newunicodechar{α}{\ensuremath{\alpha}}
\newunicodechar{β}{\ensuremath{\beta}}
\newunicodechar{γ}{\ensuremath{\gamma}}
\newunicodechar{δ}{\ensuremath{\delta}}
\newunicodechar{ε}{\ensuremath{\varepsilon}}
\newunicodechar{θ}{\ensuremath{\theta}}
\newunicodechar{λ}{\ensuremath{\lambda}}
\newunicodechar{μ}{\ensuremath{\mu}}
\newunicodechar{σ}{\ensuremath{\sigma}}
\newunicodechar{τ}{\ensuremath{\tau}}
\newunicodechar{φ}{\ensuremath{\varphi}}
\newunicodechar{ψ}{\ensuremath{\psi}}
\newunicodechar{ω}{\ensuremath{\omega}}
\newunicodechar{Σ}{\ensuremath{\Sigma}}
% majuscules produites par \MakeUppercase dans les titres (α-aminés -> Α)
\newunicodechar{Α}{\ensuremath{\alpha}}
\newunicodechar{Β}{\ensuremath{\beta}}
\newunicodechar{Γ}{\ensuremath{\Gamma}}
\newunicodechar{Δ}{\ensuremath{\Delta}}
\newunicodechar{Ε}{\ensuremath{\varepsilon}}
\newunicodechar{Θ}{\ensuremath{\Theta}}
\newunicodechar{Λ}{\ensuremath{\Lambda}}
\newunicodechar{Μ}{\ensuremath{\mu}}
\newunicodechar{Τ}{\ensuremath{\tau}}
\newunicodechar{Φ}{\ensuremath{\Phi}}
\newunicodechar{Ψ}{\ensuremath{\Psi}}
\newunicodechar{Ω}{\ensuremath{\Omega}}
\newunicodechar{↦}{\ensuremath{\mapsto}}
\newunicodechar{∈}{\ensuremath{\in}}
\newunicodechar{∉}{\ensuremath{\notin}}
\newunicodechar{∧}{\ensuremath{\wedge}}
\newunicodechar{⇔}{\ensuremath{\Leftrightarrow}}
\newunicodechar{⟺}{\ensuremath{\Longleftrightarrow}}
\newunicodechar{⇒}{\ensuremath{\Rightarrow}}
\newunicodechar{⟹}{\ensuremath{\Longrightarrow}}
\newunicodechar{∘}{\ensuremath{\circ}}
\newunicodechar{∪}{\ensuremath{\cup}}
\newunicodechar{∩}{\ensuremath{\cap}}
\newunicodechar{≡}{\ensuremath{\equiv}}
\newunicodechar{⊂}{\ensuremath{\subset}}
\newunicodechar{⊥}{\ensuremath{\perp}}
\newunicodechar{∃}{\ensuremath{\exists}}
\newunicodechar{∀}{\ensuremath{\forall}}
\newunicodechar{∛}{\ensuremath{\sqrt[3]{\,}}}
\newunicodechar{∜}{\ensuremath{\sqrt[4]{\,}}}
\newunicodechar{⃗}{\ensuremath{{}^{\to}}}
\newunicodechar{ⁿ}{\ifmmode^{n}\else\textsuperscript{\ensuremath{n}}\fi}
\newunicodechar{ˣ}{\ifmmode^{x}\else\textsuperscript{\ensuremath{x}}\fi}
\newunicodechar{ᵇ}{\ifmmode^{b}\else\textsuperscript{\ensuremath{b}}\fi}
\newunicodechar{ʳ}{\ifmmode^{r}\else\textsuperscript{\ensuremath{r}}\fi}
\newunicodechar{ᵘ}{\ifmmode^{u}\else\textsuperscript{\ensuremath{u}}\fi}
\newunicodechar{ʸ}{\ifmmode^{y}\else\textsuperscript{\ensuremath{y}}\fi}
\newunicodechar{ᵃ}{\ifmmode^{a}\else\textsuperscript{\ensuremath{a}}\fi}
\newunicodechar{ᵉ}{\ifmmode^{e}\else\textsuperscript{\ensuremath{e}}\fi}
\newunicodechar{ᵏ}{\ifmmode^{k}\else\textsuperscript{\ensuremath{k}}\fi}
\newunicodechar{ᵖ}{\ifmmode^{p}\else\textsuperscript{\ensuremath{p}}\fi}
\newunicodechar{ᴺ}{\ifmmode^{N}\else\textsuperscript{\ensuremath{N}}\fi}
\newunicodechar{ᶜ}{\ifmmode^{c}\else\textsuperscript{\ensuremath{c}}\fi}
\newunicodechar{ʰ}{\ifmmode^{h}\else\textsuperscript{\ensuremath{h}}\fi}
\newunicodechar{ᵠ}{\ifmmode^{q}\else\textsuperscript{\ensuremath{q}}\fi}
\newunicodechar{ₙ}{\ifmmode_{n}\else\textsubscript{\ensuremath{n}}\fi}
\newunicodechar{ₐ}{\ifmmode_{a}\else\textsubscript{\ensuremath{a}}\fi}
\newunicodechar{ᵢ}{\ifmmode_{i}\else\textsubscript{\ensuremath{i}}\fi}
\newunicodechar{ᵣ}{\ifmmode_{r}\else\textsubscript{\ensuremath{r}}\fi}
\newunicodechar{ₖ}{\ifmmode_{k}\else\textsubscript{\ensuremath{k}}\fi}
\newunicodechar{ᵦ}{\ifmmode_{\beta}\else\textsubscript{\ensuremath{\beta}}\fi}
\newunicodechar{ₓ}{\ifmmode_{x}\else\textsubscript{\ensuremath{x}}\fi}
\newfontfamily\popdisplay{ArchivoBlack-Regular.ttf}[Path=../../../fonts/]
\newfontfamily\poplogo{SpaceGrotesk-Bold.ttf}[Path=../../../fonts/]
\newfontfamily\popsemi{PlusJakartaSans-SemiBold.ttf}[Path=../../../fonts/]
\newfontfamily\popxbold{PlusJakartaSans-ExtraBold.ttf}[Path=../../../fonts/]
\newfontfamily\popxboldls{PlusJakartaSans-ExtraBold.ttf}[Path=../../../fonts/,LetterSpace=3.0]

\setlist[enumerate]{leftmargin=1.7em,itemsep=5pt,topsep=4pt}
\setlist[itemize]{leftmargin=1.7em,itemsep=4pt,topsep=4pt,label={\color{poporange}\textbullet}}
\setlength{\parindent}{0pt}
\setlength{\parskip}{5pt}
\renewcommand{\arraystretch}{1.25}

% pgfplots : style Pop par défaut
\pgfplotsset{
  every axis/.append style={axis line style={popink,line width=1pt},tick style={popink},
    grid style={popline,line width=0.5pt},label style={font=\small},tick label style={font=\footnotesize}},
  popcurve/.style={poporange,line width=1.8pt,no markers,smooth},
}
% Même style disponible dans un \draw TikZ simple (hors pgfplots)
\tikzset{popcurve/.style={poporange,line width=1.8pt}}
\tikzset{popgrid/.style={popline,very thin}}
\colorlet{popcurve}{poporange} % le nom du style est parfois utilisé comme couleur

% ============================================================
% Pill / squiggle
% ============================================================
\newcommand{\poppill}[3]{%
  \tikz[baseline=-0.55ex]\node[draw=popink,line width=1.2pt,rounded corners=2.6pt,%
    fill=#2,inner xsep=6.5pt,inner ysep=3pt]{%
    \popxboldls\fontsize{7}{7}\selectfont\color{#3}\MakeUppercase{#1}};%
}
\newcommand{\popsquiggle}[1]{%
  \tikz[baseline=-0.4ex]{
    \draw[#1,line width=2.6pt,line cap=round]
      plot [smooth] coordinates {(0,0.09) (0.34,0.01) (0.68,0.21) (1.02,0.01) (1.36,0.10)};
  }%
}
% Mot-clé surligné (marqueur orange sous le texte)
\newcommand{\cle}[1]{{\popxbold\color{popdark}#1}}

% ============================================================
% En-tête / pied
% ============================================================
\pagestyle{fancy}
\fancyhf{}
\renewcommand{\headrulewidth}{0pt}
\renewcommand{\footrulewidth}{0pt}
\lhead{\raisebox{-0.15ex}{\poplogo\fontsize{13}{13}\selectfont\color{popink}OnBuch\color{poporange}+}}
\rhead{\poppill{\DOCMATIERE\ \textbullet\ \DOCNIVEAU\ \textbullet\ Leçon \DOCLECON}{white}{popink}}
\lfoot{\popxbold\fontsize{7.6}{7.6}\selectfont\color{popmuted}\MakeUppercase{Cours signé OnBuch+}\ \textbullet\ {\popsemi\DOCTITRE}}
\rfoot{\tikz[baseline=-0.6ex]\node[rounded corners=8.5pt,fill=popink,minimum height=17pt,minimum width=17pt,inner xsep=5pt,inner sep=0]{\popxbold\fontsize{8}{8}\selectfont\color{popcream}\thepage\,/\,\pageref*{LastPage}};}

% ============================================================
% Titres
% ============================================================
\newcommand{\chaptitre}[2]{%
  \begin{tcolorbox}[enhanced,colback=poporange,colframe=popink,boxrule=2.6pt,arc=11pt,
      drop shadow=popink,left=15pt,right=15pt,top=12pt,bottom=11pt,before skip=3pt,after skip=18pt]
    \poppill{#2}{popink}{popcream}\par\vspace{9pt}
    {\raggedright\hyphenpenalty=10000\exhyphenpenalty=10000\popdisplay\fontsize{22}{24}\selectfont\color{popcream}\MakeUppercase{#1}\par}
  \end{tcolorbox}
}
% Section numérotée : pastille orange avec le numéro + titre Archivo
\newcounter{popsec}
\newcommand{\coursec}[1]{%
  \stepcounter{popsec}\setcounter{popsub}{0}%
  \par\vspace{16pt}\noindent
  \begin{minipage}{\linewidth}
  \tikz[baseline=-0.7ex]\node[draw=popink,line width=1.6pt,fill=poporange,rounded corners=4pt,minimum size=22pt,inner sep=0]{\popdisplay\fontsize{12}{12}\selectfont\color{popcream}\thepopsec};%
  \hspace{8pt}{\popdisplay\fontsize{15}{17}\selectfont\color{popink}\MakeUppercase{#1}}\par
  \vspace{4pt}\noindent{\color{poporange}\rule{36pt}{3.6pt}}
  \end{minipage}\par\nopagebreak\vspace{8pt}
}
\newcounter{popsub}
\newcommand{\courssub}[1]{%
  \stepcounter{popsub}%
  \par\vspace{9pt}\noindent{\popxbold\fontsize{11.5}{13}\selectfont\color{popink}{\color{poporange}\thepopsec.\thepopsub}\hspace{6pt}#1}\par\nopagebreak\vspace{4pt}
}

% ============================================================
% Boîtes pédagogiques — même recette : fond blanc, bordure épaisse colorée,
% ombre dure encre, badge plein. Titre optionnel [#1].
% ============================================================
\tcbset{popbase/.style={breakable,enhanced,colback=white,boxrule=2.3pt,arc=9pt,
  shadow={3pt}{-3pt}{0pt}{popink},left=14pt,right=14pt,top=11pt,bottom=11pt,before skip=11pt,after skip=15pt,
  fontupper=\popsemi\fontsize{10.6}{14.5}\selectfont\color{popink},
  fontlower=\popsemi\fontsize{10.6}{14.5}\selectfont\color{popink}}}
% Coche dessinée (le glyphe ✓ n'existe pas dans les polices du document)
\newcommand{\popcheck}[1]{\tikz[baseline=-0.5ex]\draw[#1,line width=1.6pt,line cap=round,line join=round](-0.13,0.0)--(-0.04,-0.09)--(0.13,0.1);}
\providecommand{\coche}{\popcheck{popgreen}}
\providecommand{\resultat}[1]{\par\smallskip\noindent\fcolorbox{popgreen}{popgreenL}{\parbox{\dimexpr\linewidth-2\fboxsep-2\fboxrule\relax}{\textbf{Résultat.} #1}}\par\smallskip}
\newcommand{\popboxhead}[3]{\poppill{#1}{#2}{white}\ifx\relax#3\relax\else\hspace{6pt}{\popxbold\fontsize{10.6}{12}\selectfont\color{popink}#3}\fi\par\vspace{7pt}}

\newtcolorbox{aretenir}[1][]{popbase,colframe=popgreen,before upper={\popboxhead{À retenir}{popgreen}{#1}}}
\newtcolorbox{exemplebox}[1][]{popbase,colframe=poporange,before upper={\popboxhead{Exemple}{poporange}{#1}}}
\newtcolorbox{definition}[1][]{popbase,colframe=popblue,before upper={\popboxhead{Définition}{popblue}{#1}}}
\newtcolorbox{propriete}[1][]{popbase,colframe=poppurple,before upper={\popboxhead{Propriété}{poppurple}{#1}}}
\newtcolorbox{methode}[1][]{popbase,colframe=popgold,before upper={\popboxhead{Méthode}{popgold}{#1}}}
\newtcolorbox{attention}[1][]{popbase,colframe=poppink,before upper={\popboxhead{Attention}{poppink}{#1}}}
\newtcolorbox{experience}[1][]{popbase,colframe=popblue,colback=popblueL,before upper={\popboxhead{Expérience}{popblue}{#1}}}
% « Le savais-tu ? » : carte violette pleine (contexte, culture, Cameroun)
\newtcolorbox{savaistu}[1][]{popbase,colback=poppurple,colframe=popink,
  fontupper=\popsemi\fontsize{10.4}{14.2}\selectfont\color{white},
  before upper={\poppill{Le savais-tu\,?}{white}{poppurple}\ifx\relax#1\relax\else\hspace{6pt}{\popxbold\color{white}#1}\fi\par\vspace{7pt}}}
% Exercice résolu : énoncé en haut, \tcblower puis la solution sur fond crème
\newcounter{popexo}\newcounter{popexr}
\newtcolorbox{exoresolu}[1][]{popbase,colframe=poporange,colbacklower=popcream,
  segmentation style={popink,line width=1.2pt,dashed},
  before upper={\stepcounter{popexr}\popboxhead{Exercice résolu \thepopexr}{poporange}{#1}},
  before lower={\poppill{Solution}{popink}{popcream}\par\vspace{6pt}}}
% Exercice (non résolu en place) — la correction est dans la partie Corrigés
\newtcolorbox{exercice}[1][]{popbase,colframe=popink,
  before upper={\stepcounter{popexo}\popboxhead{Exercice \thepopexo}{popink}{#1}}}
% Corrigé
\newtcolorbox{corrige}[1][]{popbase,colframe=popgreen,colback=popgreenL,
  before upper={\popboxhead{Corrigé}{popgreen}{#1}}}
% Objectifs / prérequis / bilan
\newtcolorbox{objectifs}{popbase,colframe=popink,colback=poporangeL,
  before upper={\popboxhead{Objectifs de la leçon}{popink}{}}}
\newtcolorbox{prerequis}{popbase,colframe=popink,colback=popblueL,
  before upper={\popboxhead{Prérequis}{popblue}{}}}
\newtcolorbox{bilan}{popbase,colframe=popink,colback=white,boxrule=2.8pt,
  before upper={\popboxhead{Fiche bilan}{poporange}{L'essentiel à connaître pour l'examen}}}
% Figure encadrée avec légende
\newtcolorbox{popfigure}[1][]{enhanced,breakable=false,colback=white,colframe=popline,boxrule=1.4pt,arc=8pt,
  left=10pt,right=10pt,top=10pt,bottom=8pt,before skip=10pt,after skip=12pt,halign=center,
  lower separated=false,halign lower=center,
  fontlower=\popsemi\fontsize{9}{11.5}\selectfont\color{popmuted},
  #1}
\newcommand{\legende}[1]{\par\vspace{6pt}{\popsemi\fontsize{9}{11.5}\selectfont\color{popmuted}#1}}

% ============================================================
% Signature OnBuch+
% ============================================================
\newcommand{\poplogoinline}[1]{{\poplogo\fontsize{#1}{#1}\selectfont\color{popink}OnBuch\color{poporange}+}}
\newcommand{\signatureonbuch}{%
  \begin{tcolorbox}[enhanced,colback=popink,colframe=popink,boxrule=2.2pt,arc=10pt,
      drop shadow=poporange,left=14pt,right=14pt,top=10pt,bottom=10pt,before skip=0pt,after skip=0pt]
    \tikz[baseline=-0.7ex]\node[circle,fill=popgreen,draw=popcream,line width=1.3pt,minimum size=22pt,inner sep=0]{\popcheck{white}};%
    \hspace{9pt}\begin{minipage}[c]{0.62\linewidth}
      {\popxbold\fontsize{12}{13}\selectfont\color{popcream}Signé\ \textbullet\ {\poplogo OnBuch\color{poporange}+}}\par\vspace{2pt}
      {\raggedright\popsemi\fontsize{8}{10.5}\selectfont\color{popline}\MakeUppercase{Équipe pédagogique OnBuch+}\\\MakeUppercase{Conforme au programme officiel MINESEC}\par}
    \end{minipage}\hfill
    {\poplogo\fontsize{22}{22}\selectfont\color{popcream}OnBuch\color{poporange}+}
  \end{tcolorbox}%
}

% ============================================================
% Page de garde
% #1 titre · #2 matière · #3 niveau · #4 module · #5 n° leçon · #6 durée ·
% #7 description / objectifs (texte ou itemize)
% ============================================================
\newcommand{\pagegarde}[7]{%
  \thispagestyle{empty}
  \newgeometry{a4paper,margin=1.7cm,top=1.6cm,bottom=1.4cm}%
  \begin{tikzpicture}[remember picture,overlay]
    \fill[poporange!18] ([xshift=-3.2cm,yshift=-3.6cm]current page.north east) circle (4.6cm);
    \fill[poppurple!14] ([xshift=2.4cm,yshift=7.6cm]current page.south west) circle (3.8cm);
  \end{tikzpicture}%
  {\poplogo\fontsize{30}{30}\selectfont\color{popink}OnBuch\color{poporange}+}\hfill
  \raisebox{0.6ex}{\poppill{Cours signé \textbullet\ Premium}{popink}{popcream}}\par
  \vspace{1pt}\popsquiggle{poppurple}\par
  \vspace{8pt}
  \begin{tcolorbox}[enhanced,colback=poporange,colframe=popink,boxrule=2.8pt,arc=13pt,
      drop shadow=popink,left=18pt,right=18pt,top=12pt,bottom=13pt,after skip=10pt]
    \poppill{#2 \textbullet\ #3}{popink}{popcream}\hspace{5pt}\poppill{#4}{white}{popink}\par\vspace{8pt}
    {\popdisplay\fontsize{14}{14}\selectfont\color{popink}LEÇON #5}\par\vspace{4pt}
    {\raggedright\hyphenpenalty=10000\exhyphenpenalty=10000\popdisplay\fontsize{25}{27}\selectfont\color{popcream}\MakeUppercase{#1}\par}
    \vspace{8pt}
    {\popxbold\fontsize{9.5}{9.5}\selectfont\color{popink}\MakeUppercase{Durée conseillée : #6}}
  \end{tcolorbox}
  \begin{tcolorbox}[enhanced,colback=white,colframe=popink,boxrule=2.2pt,arc=10pt,
      drop shadow=popink,left=16pt,right=16pt,top=10pt,bottom=10pt,after skip=8pt]
    \poppill{Dans cette leçon}{poppurple}{white}\par\vspace{5pt}
    \popsemi\fontsize{9.3}{12.6}\selectfont\color{popink}#7
  \end{tcolorbox}
  \noindent\begin{minipage}[t]{0.487\linewidth}
    \begin{tcolorbox}[enhanced,colback=popgreen,colframe=popink,boxrule=2.2pt,arc=10pt,
        drop shadow=popink,left=11pt,right=11pt,top=10pt,bottom=10pt,height=3.1cm,valign=center]
      \tcbox[colback=white,colframe=popink,boxrule=1.3pt,arc=5pt,boxsep=0pt,nobeforeafter,
        left=3pt,right=3pt,top=3pt,bottom=3pt,valign=center]{\qrcode[height=1.9cm]{https://play.google.com/store/apps/details?id=com.luvvix.onbuch&hl=fr}}%
      \hspace{8pt}\begin{minipage}[c]{0.5\linewidth}
        {\popdisplay\fontsize{11}{12.5}\selectfont\color{white}\MakeUppercase{Télécharge OnBuch}}\par\vspace{4pt}
        {\raggedright\popsemi\fontsize{8.4}{11}\selectfont\color{white}Gratuit \textbullet\ Google Play\\Tuteur IA, annales, concours, résultats\par}
      \end{minipage}
    \end{tcolorbox}
  \end{minipage}\hfill
  \begin{minipage}[t]{0.487\linewidth}
    \begin{tcolorbox}[enhanced,colback=poppurple,colframe=popink,boxrule=2.2pt,arc=10pt,
        drop shadow=popink,left=11pt,right=11pt,top=10pt,bottom=10pt,height=3.1cm,valign=center]
      \tikz[baseline=-0.7ex]\node[circle,fill=white,draw=popink,line width=1.3pt,minimum size=1.6cm,inner sep=0]{\popdisplay\fontsize{24}{24}\selectfont\color{poppurple}+};%
      \hspace{8pt}\begin{minipage}[c]{0.6\linewidth}
        {\popdisplay\fontsize{11}{12.5}\selectfont\color{white}\MakeUppercase{Passe à OnBuch+}}\par\vspace{4pt}
        {\raggedright\popsemi\fontsize{8.4}{11}\selectfont\color{white}Tous les cours signés, Tuteur IA illimité, fiches PDF — onglet OnBuch+ de l'app\par}
      \end{minipage}
    \end{tcolorbox}
  \end{minipage}
  \vspace{8pt}
  \popcontenu
  \vfill
  \signatureonbuch
  \restoregeometry
  \newpage
}

% Rangée « Ce cours contient » (page de garde)
\newcommand{\poptile}[3]{% #1 couleur #2 gros texte #3 légende
  \begin{tcolorbox}[enhanced,colback=#1,colframe=popink,boxrule=2pt,arc=9pt,shadow={2.5pt}{-2.5pt}{0pt}{popink},
      left=6pt,right=6pt,top=9pt,bottom=9pt,halign=center,valign=center,height=1.75cm,width=\linewidth,nobeforeafter]
    {\popdisplay\fontsize{12.5}{13}\selectfont\color{white}\MakeUppercase{#2}}\par\vspace{3pt}
    {\centering\popsemi\fontsize{7.8}{9.5}\selectfont\color{white}#3\par}
  \end{tcolorbox}}
\newcommand{\popcontenu}{%
  \par\noindent{\popxboldls\fontsize{7.5}{7.5}\selectfont\color{popmuted}CE COURS SIGNÉ CONTIENT}\par\vspace{7pt}
  \noindent\begin{minipage}[t]{0.235\linewidth}\poptile{popblue}{Cours}{complet, illustré, pas à pas}\end{minipage}\hfill
  \begin{minipage}[t]{0.235\linewidth}\poptile{poporange}{Méthodes}{et exercices résolus}\end{minipage}\hfill
  \begin{minipage}[t]{0.235\linewidth}\poptile{poppink}{Exercices}{type examen + corrigés}\end{minipage}\hfill
  \begin{minipage}[t]{0.235\linewidth}\poptile{popgreen}{Fiche bilan}{l'essentiel à retenir}\end{minipage}\par}

% Sommaire
\newtcolorbox{sommaire}{enhanced,colback=white,colframe=popink,boxrule=2.2pt,arc=10pt,
  drop shadow=popink,left=18pt,right=18pt,top=13pt,bottom=13pt,before skip=6pt,after skip=20pt,
  before upper={\poppill{Sommaire}{poppurple}{white}\par\vspace{9pt}\popsemi\fontsize{10.6}{14.5}\selectfont\color{popink}}}

% Signature de fin de cours — poussée tout en bas de la dernière page
\newcommand{\findecours}{%
  \par\vfill\nopagebreak
  \begin{center}\popsquiggle{poporange}\end{center}\vspace{4pt}
  \signatureonbuch
}
\AtBeginDocument{\def\llbracket{\mathopen{[\mkern-3.5mu[}}\def\rrbracket{\mathclose{]\mkern-3.5mu]}}} % intervalles d'entiers (glyphe ⟦ absent de la police)
% Glyphes absents de Fira Math : redéfinis à partir de glyphes disponibles
\AtBeginDocument{%
  \def\setminus{\mathbin{\backslash}}\let\smallsetminus\setminus
  \def\vdots{\mathord{\vbox{\baselineskip3.2pt\lineskiplimit0pt\kern4pt\hbox{.}\hbox{.}\hbox{.}}}}%
  \def\ddots{\mathinner{\mkern1mu\raise6.5pt\hbox{.}\mkern2mu\raise3.6pt\hbox{.}\mkern2mu\raise0.7pt\hbox{.}\mkern1mu}}%
  \def\triangle{\mathord{\scalebox{0.9}{$\symup{\Delta}$}}}%
  \def\nabla{\mathord{\raisebox{\depth}{\scalebox{1}[-1]{$\symup{\Delta}$}}}}%
  \def\checkmark{\popcheck{popdark}}%
  \def\star{\mathbin{\ast}}\def\bigstar{\mathord{\ast}}%
  \def\diamond{\mathbin{\scalebox{0.6}{$\Diamond$}}}%
  \def\bigcup{\mathop{\mathchoice{\raisebox{-0.3em}{\scalebox{1.7}{$\cup$}}}{\scalebox{1.25}{$\cup$}}{\cup}{\cup}}\displaylimits}%
  \def\bigcap{\mathop{\mathchoice{\raisebox{-0.3em}{\scalebox{1.7}{$\cap$}}}{\scalebox{1.25}{$\cap$}}{\cap}{\cap}}\displaylimits}%
  \def\lesseqqgtr{\mathrel{\lessgtr}}\def\bigoplus{\mathop{\oplus}\displaylimits}%
}
\providecommand{\ssi}{si et seulement si} % abréviation fréquente
\AtBeginDocument{\providecommand{\wideparen}{\overparen}} % arc de cercle

%% FIGFIX-BEGIN v5 — correctifs globaux des figures (docs/onbuch-generation/tools/figfix)
\usepackage{adjustbox}\usepackage{etoolbox}\tcbuselibrary{raster}
% toute figure TikZ/pgfplots plus large que la ligne est réduite à la largeur disponible
\BeforeBeginEnvironment{tikzpicture}{\begin{adjustbox}{max width=\linewidth}}
\AfterEndEnvironment{tikzpicture}{\end{adjustbox}}
% légendes pgfplots lisibles par-dessus les courbes
\pgfplotsset{every axis/.append style={legend style={fill=white,fill opacity=0.92,text opacity=1,draw=black!15,inner sep=3pt}}}
% étiquettes d'arêtes (syntaxe edge["..."]) avec fond blanc
\tikzset{every edge quotes/.append style={fill=white,inner sep=1.2pt}}
%% FIGFIX-END
\begin{document}

\pagegarde{Qu'est-ce que la philosophie?}{Philosophie}{1re Tech}{Philosophie – classe de Première technique}{N15}{5 séances (fiche de progression harmonisée)}{Cette leçon introductive définit la philosophie, retrace son origine et précise son objet, sa méthode et son importance. Elle éclaire enfin le rapport entre philosophie et science, à travers des situations concrètes de la vie camerounaise.
\vspace{4pt}
\begin{multicols}{2}\small
\begin{itemize}
\item À la fin de cette leçon, je sais définir la philosophie à partir de son étymologie et des définitions des grands auteurs
\item À la fin de cette leçon, je sais expliquer l'origine de la philosophie (étonnement, doute, loisir) et identifier son objet
\item À la fin de cette leçon, je sais décrire la méthode philosophique (questionnement, analyse, argumentation, synthèse)
\item À la fin de cette leçon, je sais distinguer la démarche philosophique de la démarche scientifique
\item À la fin de cette leçon, je sais justifier l'importance de la philosophie dans la vie personnelle et sociale
\item À la fin de cette leçon, je sais appliquer la démarche philosophique à une situation concrète de la vie courante camerounaise
\end{itemize}\end{multicols}}

\begin{sommaire}
\begin{multicols}{2}
\begin{enumerate}
\item Définitions de la philosophie
\item Origine et objet de la philosophie
\item La méthode de la philosophie
\item L'importance de la philosophie
\item Philosophie et science
\item Synthèse et application à la vie courante
\item Activité d'intégration
\item Méthodes et exercices résolus
\item Exercices
\item Corrigés des exercices
\item Fiche bilan
\end{enumerate}
\end{multicols}
\end{sommaire}

\begin{objectifs}
\begin{itemize}
\item À la fin de cette leçon, je sais définir la philosophie à partir de son étymologie et des définitions des grands auteurs
\item À la fin de cette leçon, je sais expliquer l'origine de la philosophie (étonnement, doute, loisir) et identifier son objet
\item À la fin de cette leçon, je sais décrire la méthode philosophique (questionnement, analyse, argumentation, synthèse)
\item À la fin de cette leçon, je sais distinguer la démarche philosophique de la démarche scientifique
\item À la fin de cette leçon, je sais justifier l'importance de la philosophie dans la vie personnelle et sociale
\item À la fin de cette leçon, je sais appliquer la démarche philosophique à une situation concrète de la vie courante camerounaise
\item À la fin de cette leçon, je sais rédiger et justifier une réponse argumentée en suivant les étapes de la méthode philosophique
\end{itemize}
\end{objectifs}

\begin{prerequis}
\begin{itemize}
\item Notions de citoyenneté et de vie en société vues au collège (EMC)
\item Capacité à lire et commenter un texte court (français, seconde)
\item Notion de science et de démarche expérimentale (SVT/physique de seconde)
\item Vocabulaire de base : opinion, vérité, argument, preuve
\end{itemize}
\end{prerequis}

\newpage

\coursec{Définitions de la philosophie}

\courssub{Situation d'accroche et problématique}

Lors d'une assemblée de quartier à Yaoundé, un jeune élève de Première assiste à une scène qui le laisse pensif. Deux voisins se disputent la limite de leurs terrains. Le chef de quartier tranche en invoquant la \emph{coutume} : « Chez nous, la borne, c'est le vieux manguier planté par les ancêtres. » L'ingénieur des travaux publics, lui, propose un \emph{bornage scientifique} : « Sortons le GPS et le plan cadastral, la mesure ne ment pas. » Alors un vieux notable se lève et demande : « Mais au fond, qu'est-ce qui est \emph{juste} ? » Silence. Chacun, pourtant, a raisonné, questionné, justifié son point de vue. Sans le savoir, tous les trois font déjà de la \cle{philosophie} : ils cherchent à fonder rationnellement ce qui doit être cru et ce qui doit être fait.

Cette scène de la vie quotidienne camerounaise soulève les questions qui guident toute cette leçon : qu'est-ce que la philosophie exactement ? D'où vient ce mot ? Que cherche celui qui philosophie ? En quoi la philosophie se distingue-t-elle de la sagesse du notable, de la coutume du chef et de la science de l'ingénieur ? Pour répondre, nous partirons du mot lui-même (son étymologie grecque), puis nous examinerons les grandes définitions proposées par les philosophes, avant de dégager une définition synthétique.

\courssub{L'étymologie du mot « philosophie »}

\textbf{L'intuition d'abord.} Pour comprendre une chose, il est souvent utile de commencer par son nom. En français comme en grec, les mots composés gardent la trace de leur histoire. « Philosophie » vient du grec ancien et se décompose en deux racines :

\begin{itemize}
\item \emph{philein} (verbe) : \cle{aimer}, désirer, rechercher avec affection ;
\item \emph{sophia} (nom) : la \cle{sagesse}, le savoir, la science.
\end{itemize}

Littéralement, la philosophie est donc « l'\cle{amour de la sagesse} ». Celui qui philosophie (\emph{philosophos}) n'est pas celui qui \emph{possède} la sagesse, mais celui qui l'\emph{aime} et la \emph{recherche}. Cette nuance est capitale : le mot désigne un mouvement, un désir, une quête — non un acquis.

\begin{definition}[Définition étymologique de la philosophie]
Du grec \emph{philein} (aimer) et \emph{sophia} (sagesse, savoir), la philosophie signifie littéralement « \cle{amour de la sagesse} ». Le philosophe est celui qui, ne possédant pas la sagesse, la désire, la recherche et s'en rapproche par l'exercice de la raison.
\end{definition}

\begin{popfigure}
\begin{center}
\begin{tikzpicture}[
every node/.style={font=\small},
racine/.style={rectangle, rounded corners=3pt, draw=popblue, fill=popblueL, thick, align=center, inner sep=6pt},
sens/.style={rectangle, rounded corners=3pt, draw=poporange, fill=poporangeL, thick, align=center, inner sep=5pt},
mot/.style={rectangle, rounded corners=4pt, draw=popgreen, fill=popgreenL, very thick, align=center, inner sep=8pt, font=\bfseries},
fleche/.style={-{Stealth[length=3mm]}, thick, popdark}
]
\node[mot, align=center] (philo) at (0,0){\large PHILO-SOPHIE\\[2pt] \normalsize « amour de la sagesse »};
\node[racine, align=center] (philein) at (-4,2.2){\emph{philein}\\ (verbe grec)};
\node[racine, align=center] (sophia) at (4,2.2){\emph{sophia}\\ (nom grec)};
\node[sens, align=center] (aimer) at (-4,4.4){aimer, désirer,\\ rechercher};
\node[sens, align=center] (sagesse) at (4,4.4){sagesse,\\ savoir, science};
\draw[fleche] (philein) -- (philo);
\draw[fleche] (sophia) -- (philo);
\draw[fleche] (aimer) -- (philein);
\draw[fleche] (sagesse) -- (sophia);
\node[align=center, font=\itshape\small, popmuted] at (0,-1.4) {Le philosophe n'est pas le sage : il est celui qui cherche la sagesse.};
\end{tikzpicture}
\end{center}
\legende{Arbre étymologique du mot « philosophie » : les deux racines grecques \emph{philein} et \emph{sophia} se combinent pour désigner l'amour de la sagesse.}
\end{popfigure}

\textbf{Le mythe de Pythagore : naissance du mot.} D'où vient cette étrange idée de se dire « amoureux de la sagesse » plutôt que « sage » ? La tradition attribue l'invention du mot \emph{philosophos} à \textbf{Pythagore} (VI\textsuperscript{e} siècle av. J.-C.). Selon le récit rapporté par les Anciens (notamment Diogène Laërce et Cicéron), on demanda à Pythagore qui il était. Il répondit par une comparaison : la vie ressemble à une grande fête où l'on vient pour trois raisons. Certains viennent pour \emph{vendre} et s'enrichir (ce sont les ambitieux) ; d'autres viennent pour \emph{combattre} et gagner la gloire (ce sont les coureurs) ; d'autres enfin viennent simplement pour \emph{regarder} et comprendre le spectacle : ce sont les philosophes. Pythagore ajoutait que seuls les dieux possèdent la sagesse parfaite ; l'homme, lui, ne peut que l'aimer et la poursuivre. Ainsi le philosophe n'est \emph{ni un dieu} (qui possède la sagesse) \emph{ni un savant} (qui accumule des connaissances particulières) : il est un « \cle{amoureux de la sagesse} », un chercheur humble et infatigable.

\begin{attention}[Erreur fréquente : confondre philosophie et sagesse]
Le piège classique de la dissertation consiste à définir la philosophie comme « la sagesse » ou « la science de tout ». C'est faux. Le préfixe \emph{philo-} indique un \emph{désir}, non une possession : le philosophe \textbf{cherche} la sagesse, il ne la \textbf{possède} pas. Dire « le philosophe est un sage » est une erreur ; il faut dire : « le philosophe est celui qui tend vers la sagesse sans jamais prétendre l'avoir atteinte ». Cette humilité distingue le philosophe du prétendu savant (le \emph{sophiste}, que Socrate combattait).
\end{attention}

\begin{savaistu}[La première apparition du mot]
Le mot « philosophie » apparaît pour la première fois dans les écrits de l'historien grec \textbf{Hérodote}, au V\textsuperscript{e} siècle av. J.-C., sous la forme du verbe \emph{philosophein} (« philosopher »), qu'il emploie à propos du législateur Solon voyageant « par amour du savoir ». C'est ensuite Héraclide du Pont, disciple de Platon, qui fixe le récit de Pythagore inventeur du mot \emph{philosophos}. Le mot voyage ensuite dans toutes les langues : \emph{philosophia} en latin, \emph{philosophy} en anglais, « philosophie » en français.
\end{savaistu}

\courssub{Les grandes définitions de la philosophie chez les philosophes}

L'étymologie donne une première lumière, mais elle ne suffit pas : que cherche exactement le philosophe, et comment ? Les grands auteurs ont chacun proposé une définition qui éclaire un aspect de la philosophie.

\textbf{Platon : la philosophie comme recherche de la vérité et science des idées.} Pour \textbf{Platon} (427--348 av. J.-C.), disciple de Socrate, le monde que nous percevons par les sens n'est qu'une apparence changeante, comparable aux ombres projetées sur la paroi de la caverne (allégorie de la caverne, \emph{République}, VII). La philosophie est l'effort par lequel l'âme se détourne des apparences sensibles pour s'élever vers les \emph{Idées} (ou Formes) — les réalités intelligibles, éternelles et vraies (le Juste en soi, le Beau en soi, le Vrai en soi). La philosophie est ainsi \cle{recherche de la vérité} et \cle{science des idées} : elle vise un savoir certain et universel, au-delà des opinions (\emph{doxa}). Dans notre assemblée de quartier, le notable qui demande « qu'est-ce qui est juste ? » fait précisément un geste platonicien : il refuse de s'en tenir aux apparences (le manguier, le papier cadastral) et interroge la justice elle-même.

\textbf{Aristote : la philosophie comme science des premières causes et des premiers principes.} Pour \textbf{Aristote} (384--322 av. J.-C.), disciple de Platon, toute recherche commence par l'\emph{étonnement} (\emph{thaumazein}) : « C'est par l'étonnement que les hommes, aujourd'hui comme au commencement, ont commencé de philosopher » (\emph{Métaphysique}, A, 2). Mais la philosophie ne s'arrête pas à l'étonnement : elle remonte des effets vers les causes, et des causes proches vers les \cle{premières causes} et les \cle{premiers principes} de toutes choses. Là où le technicien sait \emph{comment} une machine fonctionne, le philosophe demande \emph{pourquoi} il y a quelque chose plutôt que rien, ce qu'est la cause en général, ce qui fonde toute explication. La philosophie est ainsi la science la plus haute, celle qu'Aristote appelle « philosophie première » ou théologie (au sens de science de l'être en tant qu'être).

\textbf{Kant : la philosophie en quatre questions.} À l'époque moderne, \textbf{Emmanuel Kant} (1724--1804) résume le champ entier de la philosophie en quatre questions fondamentales (\emph{Logique}, 1800) :
\begin{enumerate}
\item \emph{Que puis-je savoir ?} — question de la connaissance (théorique : métaphysique, épistémologie) ;
\item \emph{Que dois-je faire ?} — question de l'action (pratique : morale, éthique) ;
\item \emph{Que m'est-il permis d'espérer ?} — question de l'espérance (religion, postulates de la raison pratique) ;
\item \emph{Qu'est-ce que l'homme ?} — question qui réunit les trois premières (anthropologie philosophique).
\end{enumerate}
Cette définition montre que la philosophie n'est pas un savoir parmi d'autres : elle est l'interrogation totale de l'homme sur lui-même, sur ce qu'il peut connaître, ce qu'il doit faire et ce qu'il peut espérer. Toutes les questions de notre assemblée de quartier — que savons-nous du terrain ? que devons-nous décider ? qu'est-ce qui est juste pour l'homme ? — rentrent dans ce cadre kantien.

\textbf{Alain : « un effort pour penser juste ».} Le philosophe français \textbf{Alain} (Émile Chartier, 1868--1951) donne de la philosophie une définition volontairement modeste et pratique : « La philosophie n'est pas un savoir, c'est un \cle{effort pour penser juste}. » (\emph{Propos sur l'éducation}). Autrement dit, la philosophie n'est pas un stock de vérités à apprendre par cœur ; elle est une \emph{activité}, un exercice permanent de la pensée qui se corrige elle-même, qui examine les opinions, qui refuse les préjugés et qui exige des preuves. Philosopher, c'est moins « savoir » que « bien penser ».

\begin{exemplebox}[Le sage traditionnel camerounais et le philosophe grec]
Dans les villages du Cameroun, le \emph{notable}, le \emph{chef traditionnel} ou le \emph{griot} des Grassfields est réputé « sage » : il connaît la coutume, les généalogies, les proverbes, et sa parole tranche les conflits. Mais son savoir est \emph{reçu} des ancêtres et transmis sans être discuté : on ne demande pas au proverbe de prouver sa vérité. Le philosophe grec, lui, partage avec le sage l'amour de la vérité et le souci du juste, mais il s'en distingue par la \emph{méthode} : il questionne tout, exige des raisons, accepte le débat contradictoire et remet en cause même la tradition. Le notable dit : « Nos pères ont dit que... » ; le philosophe demande : « Est-ce vrai, et pourquoi ? ». La philosophie hérite ainsi de la sagesse traditionnelle son objet (le bien vivre, la justice), mais elle lui ajoute l'exigence rationnelle et critique.
\end{exemplebox}

\begin{center}
\begin{tabularx}{\linewidth}{|l|X|X|}\hline
\rowcolor{popblueL} \textbf{Auteur} & \textbf{Définition} & \textbf{Idée directrice} \\ \hline
Pythagore & Le philosophe est un « amoureux de la sagesse » & Humilité : la sagesse se désire, elle ne se possède pas \\ \hline
Platon & Recherche de la vérité, science des Idées & S'élever des apparences sensibles vers l'intelligible \\ \hline
Aristote & Science des premières causes et des premiers principes & Remonter de l'étonnement vers les causes ultimes \\ \hline
Kant & Réponse à quatre questions : savoir, agir, espérer, l'homme & La philosophie comme interrogation totale sur l'homme \\ \hline
Alain & « Pas un savoir, un effort pour penser juste » & La philosophie comme exercice critique de la pensée \\ \hline
\end{tabularx}
\end{center}

\courssub{Définition synthétique}

En croisant l'étymologie et ces grandes définitions, on peut dégager les traits communs de la philosophie. Platon et Aristote insistent sur la \emph{rationalité} (chercher le vrai par des raisons) ; Kant montre l'\emph{universalité} des questions (tout l'humain : connaître, agir, espérer) ; Alain souligne le caractère \emph{critique} de l'effort (penser juste, contre les préjugés) ; Pythagore enfin rappelle que cette recherche est infinie.

\begin{definition}[Définition synthétique de la philosophie]
La philosophie est une \cle{réflexion rationnelle}, \cle{critique} et \cle{universelle} sur l'\cle{homme}, le \cle{monde} et les \cle{valeurs}. \emph{Rationnelle} : elle procède par concepts, arguments et preuves, non par autorité ou révélation. \emph{Critique} : elle examine et discute toutes les opinions, y compris les siennes. \emph{Universelle} : elle ne s'occupe pas d'un objet particulier (comme la physique ou l'histoire), mais interroge le tout de l'expérience humaine : le vrai, le bien, le beau, le juste.
\end{definition}

\begin{aretenir}[L'essentiel de la section]
\begin{itemize}
\item Étymologie : \emph{philein} (aimer) + \emph{sophia} (sagesse) = amour de la sagesse.
\item Pythagore : le philosophe n'est ni dieu ni savant, mais un chercheur de sagesse.
\item Platon : recherche de la vérité, science des Idées ; Aristote : science des premières causes ; Kant : quatre questions sur l'homme ; Alain : effort pour penser juste.
\item Définition synthétique : réflexion rationnelle, critique et universelle sur l'homme, le monde et les valeurs.
\item Piège à éviter : ne jamais confondre philosophie et sagesse possédée.
\end{itemize}
\end{aretenir}

\begin{exoresolu}[Analyser une définition]
\textbf{Énoncé.} Un élève écrit dans sa copie : « La philosophie est la sagesse que possèdent les vieillards et les savants. » Relevez deux erreurs contenues dans cette définition et proposez une formulation correcte, en vous appuyant sur l'étymologie et sur Alain.
\tcblower
\textbf{Solution détaillée.} Première erreur : la philosophie n'est pas la sagesse \emph{possédée}. L'étymologie (\emph{philein} + \emph{sophia}) montre qu'elle est l'\emph{amour} de la sagesse, c'est-à-dire une recherche, un désir, un mouvement — non un acquis. Le mythe de Pythagore confirme que seuls les dieux possèdent la sagesse parfaite. Deuxième erreur : la philosophie n'est pas un \emph{savoir} réservé à une catégorie de personnes (vieillards, savants). Comme le dit Alain, « la philosophie n'est pas un savoir, c'est un effort pour penser juste » : elle est une activité ouverte à tout homme qui raisonne et s'interroge. Formulation correcte : « La philosophie est l'amour de la sagesse : une réflexion rationnelle et critique, accessible à tout homme, qui recherche la vérité sans prétendre la posséder. »
\end{exoresolu}

\coursec{Origine et objet de la philosophie}

\courssub{L'origine : l'étonnement et la naissance de la raison}

Si le mot est grec, l'acte de philosopher est né dans les cités ioniennes (Milet, Éphèse) au VI\textsuperscript{e} siècle av. J.-C. avec les \emph{présocratiques} (Thalès, Anaximandre, Héraclite...). Aristote identifie le moteur psychologique : l'\cle{étonnement} (\emph{thaumazein}). Face au spectacle du monde (le ciel étoilé, le changement des saisons, la diversité des êtres), l'homme ne se contente pas de subir ou d'utiliser ; il s'interroge : « Pourquoi y a-t-il quelque chose plutôt que rien ? » (Leibniz reprendra cette formule).

\begin{propriete}[Les trois conditions de naissance de la philosophie (Cornford / Vernant)]
\begin{enumerate}
\item \textbf{La cité (polis) :} un espace public où la parole est libre, où l'on débat, où la loi prime sur la force.
\item \textbf{L'écriture alphabétique :} elle permet de fixer la pensée, de la relire, de la critiquer, de la transmettre sans déformation.
\item \textbf{La sécularisation de la pensée :} on explique le monde non plus par le mythe (\emph{mythos}, récit sacré, arbitraire) mais par le \emph{logos} (discours rationnel, démonstration, causes naturelles).
\end{enumerate}
\end{propriete}

\begin{exemplebox}[Du mythe au logos : l'exemple de l'éclipse]
Dans la tradition mythique, l'éclipse solaire est l'œuvre d'un dragon qui dévore l'astre ; il faut faire du bruit pour le chasser. Chez les Milesiens, Thalès prédit l'éclipse de -585 en calculant les cycles astronomiques : l'éclipse devient un phénomène naturel, prévisible, explicable par des causes géométriques. C'est le passage du \emph{mythos} (récit consolateur) au \emph{logos} (explication rationnelle).
\end{exemplebox}

\courssub{L'objet : l'universel, le total, le fondamental}

Contrairement aux sciences particulières (biologie, physique, économie) qui ont un objet \emph{délimité} (la cellule, le mouvement, le marché), la philosophie n'a pas d'objet \emph{propre} et \emph{limité}. Son objet est \textbf{universel} et \textbf{transversal} :

\begin{itemize}
\item \textbf{L'être en tant qu'être} (métaphysique / ontologie) : qu'est-ce qui existe ? Qu'est-ce que exister ?
\item \textbf{La connaissance} (épistémologie / gnoséologie) : que puis-je savoir ? Comment ? Avec quelle certitude ?
\item \textbf{L'action et les valeurs} (éthique, esthétique, politique) : que dois-je faire ? Qu'est-ce que le bien, le beau, le juste ?
\item \textbf{L'homme} (anthropologie philosophique) : qu'est-ce que l'homme ? Quelle est sa place dans le cosmos ?
\end{itemize}

\begin{definition}[Objet de la philosophie]
L'objet de la philosophie est l'ensemble des \cle{premiers principes} et des \cle{questions ultimes} qui fondent l'expérience humaine : le Vrai (connaissance), le Bien (morale), le Beau (esthétique), l'Être (métaphysique). Elle interroge le \emph{sens} et non seulement le \emph{fonctionnement}.
\end{definition}

\coursec{La méthode de la philosophie}

\courssub{Caractéristiques de la démarche philosophique}

La philosophie n'a pas \emph{une} méthode unique (comme la méthode expérimentale en physique), mais elle se reconnaît à des \cle{exigences méthodologiques} communes :

\begin{propriete}[Les quatre piliers de la méthode philosophique]
\begin{enumerate}
\item \textbf{La conceptualisation :} créer, définir, ordonner des \cle{concepts} (notions générales, abstraites : « justice », « cause », « liberté », « technique ») pour penser au-delà du cas particulier.
\item \textbf{La problématisation :} ne pas prendre la question pour acquise ; la transformer en \cle{problème} (tension entre deux thèses opposées, ex. : « La technique libère-t-elle ou asservit-elle l'homme ? »).
\item \textbf{L'argumentation rationnelle :} avancer des \cle{thèses} soutenues par des \cle{arguments} (logiques, factuels, analogiques), anticiper les \cle{objections} et y répondre.
\item \textbf{L'exercice critique (doute méthodique) :} soumettre toute affirmation (même la sienne, même la tradition, même l'évidence) à l'épreuve de la raison. Refuser l'argument d'autorité (\emph{ipse dixit}).
\end{enumerate}
\end{propriete}

\begin{methode}[Démarche type pour une dissertation philosophique (Bac/Probatoire)]
\begin{enumerate}
\item \textbf{Analyser le sujet :} définir les termes, repérer la tension (problème), formuler la problématique.
\item \textbf{Construire le plan :} thèse / antithèse / synthèse (dialectique) OU analyse des concepts / enjeux / réponses possibles (analytique).
\item \textbf{Rédiger l'introduction :} accroche, définitions, problématique, annonce du plan.
\item \textbf{Développer :} un argument par paragraphe (idée directrice + explication + exemple/référence + transition).
\item \textbf{Conclure :} réponse synthétique à la problématique, ouverture.
\end{enumerate}
\end{methode}

\begin{attention}[Pièges méthodologiques au Bac]
\begin{itemize}
\item \textbf{Le catalogue d'auteurs :} aligner « Platon dit..., Kant dit... » sans articuler leurs thèses à \emph{votre} problème.
\item \textbf{L'opinion non argumentée :} « Je pense que... », « Tout le monde sait que... » sans preuve.
\item \textbf{Le hors-sujet :} traiter un thème voisin au lieu du sujet exact.
\item \textbf{L'absence de problématique :} traiter le sujet comme une question de cours (récitation) au lieu d'une question ouverte.
\end{itemize}
\end{attention}

\courssub{La philosophie, science ou sagesse ? Précision sur la méthode}

\begin{center}
\begin{tabularx}{\linewidth}{|l|X|X|X|}\hline
\rowcolor{popblueL} \textbf{Critère} & \textbf{Science (particulière)} & \textbf{Sagesse / Tradition} & \textbf{Philosophie} \\ \hline
\textbf{Objet} & Particulier, délimité, factuel & Existence, conduite de vie, sens & Universel, principes, sens \\ \hline
\textbf{Méthode} & Expérimentale, hypothético-déductive & Transmission, autorité, intuition & Conceptuelle, critique, dialectique \\ \hline
\textbf{Résultat} & Vérifiable, cumulatif, objectif & Adhésion, conviction, paix intérieure & Argumenté, discutable, rationnel \\ \hline
\textbf{Statut du vrai} & Provisoire, falsifiable (Popper) & Absolu, révélé, intangible & Régulateur, horizon de la raison \\ \hline
\end{tabularx}
\end{center}

\coursec{L'importance de la philosophie}

\courssub{Pour l'individu : autonomie et lucidité}

\begin{itemize}
\item \textbf{Former l'esprit critique :} apprendre à distinguer l'argument valable du sophisme, le fait de l'opinion, la preuve du préjugé. « Sapere aude ! » (Ose savoir / Aie le courage de te servir de ton propre entendement) — devise des Lumières (Kant).
\item \textbf{Conquérir l'autonomie intellectuelle et morale :} ne plus subir passivement les modes, les idéologies, les manipulations médiatiques ou publicitaires. Devenir auteur de ses propres jugements.
\item \textbf{Donner du sens à l'existence :} clarifier ses valeurs, choisir sa vie en connaissance de cause (Socrate : « Une vie sans examen ne vaut pas d'être vécue »).
\item \textbf{Maîtriser l'angoisse :} nommer ses peurs (mort, absurdité, liberté) pour les penser et les dépasser (Épicure, Stoïciens, Existentialistes).
\end{itemize}

\courssub{Pour le technicien et le citoyen : éthique et responsabilité}

En Première technique, l'élève se prépare à un métier (génie civil, électricité, informatique, comptabilité, secrétariat...). La philosophie n'est pas un « luxe » littéraire : elle est \cle{nécessaire} à l'exercice responsable de la technique.

\begin{exemplebox}[Cas concrets au Cameroun]
\begin{itemize}
\item \textbf{Génie civil / Topographie :} L'ingénieur calcule la résistance d'un pont (science). Le philosophe (et le technicien éthique) demande : « Ce pont, pour qui ? Déplace-t-il des populations ? Respecte-t-il l'environnement ? L'argent public est-il bien utilisé ? » (\emph{Éthique de la responsabilité}, Jonas).
\item \textbf{Informatique / IA :} L'algorithme optimise un tri (science/technique). La philosophie interroge : « Biais discriminants ? Vie privée ? Remplacement de l'humain ? Responsabilité en cas d'erreur ? » (\emph{Éthique du numérique}).
\item \textbf{Comptabilité / Gestion :} La norme comptable dit comment enregistrer (technique). La philosophie demande : « La légalité suffit-elle ? Quid de l'évasion fiscale légale mais immorale ? Justice sociale ? » (\emph{Éthique des affaires, RSE}).
\end{itemize}
\end{exemplebox}

\begin{aretenir}[Importance de la philosophie pour le technicien]
\begin{itemize}
\item Éviter le \emph{technicisme} (la technique comme fin en soi, sans regard sur l'humain).
\item Développer une \cle{réflexivité} sur sa pratique professionnelle (déontologie).
\item Savoir argumenter, rédiger des rapports clairs, négocier, décider en situation complexe.
\item Contribuer au débat public camerounais (décentralisation, environnement, éducation, santé) en citoyen éclairé.
\end{itemize}
\end{aretenir}

\coursec{Philosophie et science}

\courssub{Distinctions fondamentales}

La confusion historique est ancienne : en grec, \emph{epistémè} et \emph{sophia} se chevauchent. La séparation moderne (XVII\textsuperscript{e}-XVIII\textsuperscript{e} s.) s'opère avec Galilée, Descartes, Newton : les sciences deviennent autonomes (physique, biologie...). La philosophie devient « philosophie des sciences » (épistémologie) et garde les questions que les sciences ne peuvent pas résoudre par leur méthode propre.

\begin{propriete}[Différences clés (programme MINESEC)]
\begin{enumerate}
\item \textbf{L'objet :} La science étudie le \emph{comment} (mécanismes, lois, faits particuliers, mesurables). La philosophie étudie le \emph{pourquoi} (causes premières, sens, valeurs, conditions de possibilité de la science elle-même).
\item \textbf{La méthode :} Science = observation, hypothèse, expérience, modélisation mathématique, reproductibilité. Philosophie = analyse conceptuelle, déduction logique, dialectique, pensée expérimentale (expériences de pensée), herméneutique.
\item \textbf{La vérité :} Science = vérité \emph{provisoire}, falsifiable, intersubjective (Popper : « Toute science est conjecturale »). Philosophie = vérité \emph{rationnelle}, universelle en intention, mais toujours discutée (pas d'expérience décisive unique).
\item \textbf{Le progrès :} Science = cumulatif, linéaire (on sait plus qu'hier). Philosophie = \emph{recommencement} permanent (on reprend les mêmes questions ; les réponses s'affinent mais ne s'empilent pas comme des briques).
\end{enumerate}
\end{propriete}

\courssub{Rapports et complémentarité}

\begin{itemize}
\item \textbf{La philosophie éclaire la science :} \emph{Épistémologie} (qu'est-ce qu'une loi scientifique ? une théorie ? une preuve ?), \emph{Éthique des sciences} (limites de la manipulation du vivant, IA, énergie nucléaire), \emph{Histoire des sciences} (comment les paradigmes changent, Kuhn).
\item \textbf{La science nourrit la philosophie :} Les découvertes (relativité, quantique, neurosciences, évolution, big data) obligent la philosophie à réviser ses concepts (temps, causalité, liberté, nature humaine, vérité).
\item \textbf{Ni réductionnisme, ni obscurantisme :} La philosophie ne \emph{remplace} pas la science (elle ne fait pas de physique à la place des physiciens). La science ne \emph{supprime} pas la philosophie (elle ne peut pas dire si l'usage de la bombe atomique est juste, ni ce qu'est la justice).
\end{itemize}

\begin{exemplebox}[Le technicien face au rapport science/philosophie]
Un technicien en laboratoire analyse un échantillon d'eau (science : protocole, mesures, normes OMS). Il constate une pollution. La science lui dit \emph{quoi} faire pour dépolluer (techniques). La philosophie (éthique, droit de l'environnement) l'aide à décider \emph{s'il doit} alerter sa hiérarchie malgré la pression de son employeur, \emph{qui} est responsable, \emph{au nom de quoi} (valeur santé, droit à l'eau, principe de précaution). La compétence technique est nécessaire mais \emph{insuffisante} pour l'action responsable.
\end{exemplebox}

\coursec{Synthèse et application à la vie courante}

\courssub{Carte mentale récapitulative}

\begin{popfigure}
\begin{center}
\begin{center}\tcbox[colback=popblue,colframe=popblue,coltext=white,arc=9pt,boxsep=4pt,left=6pt,right=6pt]{\bfseries\small PHILOSOPHIE (Amour de la sagesse)}\end{center}
\vspace{-2pt}
\begin{tcbraster}[raster columns=2,raster equal height=rows,raster column skip=6pt,raster row skip=6pt,enhanced,blankest]
\begin{tcolorbox}[enhanced,colback=white,colframe=popgreen,boxrule=0.9pt,arc=3pt,title={DÉFINITIONS},fonttitle=\bfseries\scriptsize,coltitle=white,colbacktitle=popgreen,left=4pt,right=4pt,top=2pt,bottom=2pt,boxsep=1pt,toptitle=1.5pt,bottomtitle=1.5pt,halign=left]
\begin{itemize}[nosep,leftmargin=*,itemsep=1pt,font=\scriptsize]
\item Étymologique : Philein + Sophia
\item Platon : Vérité / Idées
\item Aristote : Causes / Principes
\item Kant : 4 Questions
\item Alain : Effort pour penser juste
\end{itemize}
\end{tcolorbox}
\begin{tcolorbox}[enhanced,colback=white,colframe=poporange,boxrule=0.9pt,arc=3pt,title={ORIGINE \& OBJET},fonttitle=\bfseries\scriptsize,coltitle=white,colbacktitle=poporange,left=4pt,right=4pt,top=2pt,bottom=2pt,boxsep=1pt,toptitle=1.5pt,bottomtitle=1.5pt,halign=left]
\begin{itemize}[nosep,leftmargin=*,itemsep=1pt,font=\scriptsize]
\item Origine : Étonnement, Polis, Logos
\item Objet : Universel (Être, Savoir, Agir, Homme)
\end{itemize}
\end{tcolorbox}
\begin{tcolorbox}[enhanced,colback=white,colframe=poppurple,boxrule=0.9pt,arc=3pt,title={MÉTHODE},fonttitle=\bfseries\scriptsize,coltitle=white,colbacktitle=poppurple,left=4pt,right=4pt,top=2pt,bottom=2pt,boxsep=1pt,toptitle=1.5pt,bottomtitle=1.5pt,halign=left]
\begin{itemize}[nosep,leftmargin=*,itemsep=1pt,font=\scriptsize]
\item Conceptualisation
\item Problématisation
\item Argumentation
\item Critique / Doute
\end{itemize}
\end{tcolorbox}
\begin{tcolorbox}[enhanced,colback=white,colframe=poppink,boxrule=0.9pt,arc=3pt,title={IMPORTANCE \& SCIENCE},fonttitle=\bfseries\scriptsize,coltitle=white,colbacktitle=poppink,left=4pt,right=4pt,top=2pt,bottom=2pt,boxsep=1pt,toptitle=1.5pt,bottomtitle=1.5pt,halign=left]
\begin{itemize}[nosep,leftmargin=*,itemsep=1pt,font=\scriptsize]
\item Autonomie, Citoyenneté
\item Éthique technique / Déontologie
\item Science : Comment / Provisoire
\item Philos : Pourquoi / Sens / Valeurs
\end{itemize}
\end{tcolorbox}
\end{tcbraster}

\end{center}
\legende{Synthèse visuelle de la leçon 15 : les quatre piliers de l'introduction à la philosophie en Première technique.}
\end{popfigure}

\courssub{Application à des situations concrètes (Entraînement Probatoire)}

\begin{exoresolu}[Situation type : Dissertation courte / Question de cours]
\textbf{Énoncé.} « La technique a-t-elle besoin de la philosophie ? » Vous répondrez en une page structurée (introduction, développement en deux ou trois parties, conclusion), en mobilisant les distinctions science/philosophie et l'importance de la philosophie pour le technicien.
\tcblower
\textbf{Corrigé guidé.}
\textbf{Introduction.} Accroche : Le technicien manie des outils, des normes, des calculs (savoir-faire). Mais tout acte technique s'inscrit dans un monde humain (valeurs, fins, responsabilités). Définitions : Technique = ensemble de procédés rationnels pour transformer la matière selon une fin. Philosophie = réflexion rationnelle, critique, universelle sur l'homme, le monde, les valeurs. Problématique : La technique, par sa méthode propre (efficacité, moyens), suffit-elle à déterminer ses propres fins et ses limites éthiques, ou a-t-elle besoin de la philosophie pour s'orienter ?
\textbf{Développement.}
\textbf{I. L'autonomie de la technique (non, elle a sa rationalité propre).} La technique obéit à des lois objectives (physique, maths). Le pont tient ou ne tient pas. La philosophie ne construit pas le pont. La science/technique a sa vérité propre, vérifiable, cumulative. Le technicien n'a pas besoin de Kant pour calculer une poutre.
\textbf{II. L'aveuglement du seul savoir-faire (oui, pour les fins et les valeurs).} La technique est efficacité (moyens). Elle ne dit pas \emph{quoi} faire ni \emph{pourquoi} (fins). Ex. : L'IA peut optimiser une sélection de CV (efficacité) mais reproduire des biais racistes/sexistes (injustice). La philosophie (éthique, épistémologie) interroge les présupposés, les biais, les impacts sociaux. Principe de responsabilité (Jonas) : la puissance technique exige une éthique de la prévision.
\textbf{III. La philosophie comme garde-fou et boussole (complémentarité).} Épistémologie : clarifier le statut des modèles (incertitude, validation). Déontologie : codes éthiques professionnels (ex. ingénieurs, informaticiens). Citoyenneté : participation au débat public sur les choix technologiques (OGM, 5G, ville intelligente à Douala/Yaoundé).
\textbf{Conclusion.} La technique n'a pas besoin de la philosophie pour \emph{fonctionner}, mais elle en a besoin pour \emph{s'orienter} et \emph{se justifier} humainement. Sans philosophie, la technique risque le technicisme (moyens devenus fins). Le technicien philosophe est un technicien responsable.
\end{exoresolu}

\begin{exercice}[Entraînement personnel — À rendre pour la séance prochaine]
\begin{enumerate}
\item \textbf{QCM / Vrai-Faux (Justifiez brièvement).}
\begin{enumerate}
\item « Le philosophe est un homme qui sait tout. » (Faux : étymologie, Pythagore, Alain).
\item « La philosophie progresse comme la physique : on sait plus aujourd'hui qu'hier. » (Faux : nature du progrès philosophique, recommencement).
\item « La science répond aux questions de fait ; la philosophie répond aux questions de valeur. » (Vrai : distinction fait/valeur, Hume, Weber).
\item « Un ingénieur n'a pas besoin de philosophie pour exercer son métier. » (Faux : éthique, responsabilité, sens de l'action technique).
\end{enumerate}
\item \textbf{Analyse de citation (5 lignes max).} « La science nous donne le pouvoir sur les choses ; la philosophie nous donne le pouvoir sur nous-mêmes. » (Alain, \emph{Propos sur les pouvoirs}). Expliquez la distinction opérée et donnez un exemple technique camerounais illustrant ce pouvoir sur soi.
\item \textbf{Mini-dissertation (1 page, 30 min chrono).} Sujet : « L'étonnement est-il le début de la philosophie ? » (Mobilisez Aristote, la naissance du logos, et un exemple personnel ou professionnel).
\end{enumerate}
\end{exercice}

\begin{corrige}[Exercice 1 - QCM]
\begin{enumerate}
\item \textbf{Faux.} Le philosophe (\emph{philosophos}) est l'\emph{amant} de la sagesse, non le possesseur (Pythagore). Alain précise : « pas un savoir, un effort ».
\item \textbf{Faux.} La science est cumulative (nouveaux faits, théories plus larges). La philosophie revisite les mêmes questions fondamentales (mort, liberté, vérité) ; les réponses s'affinent mais ne s'empilent pas.
\item \textbf{Vrai.} Distinction classique (Hume : « on ne peut déduire un \emph{il faut} d'un \emph{il est} »). Science = domaine du fait, de l'efficacité, du vrai vérifiable. Philosophie = domaine de la valeur, du sens, du juste, du fondement.
\item \textbf{Faux.} L'ingénieur a besoin de compétences techniques (savoir-faire) MAIS aussi de discernement éthique (choix des matériaux, impact environnemental, sécurité publique, propriété intellectuelle, management d'équipe). La philosophie fournit ce \emph{savoir-être} et ce \emph{savoir-juger}.
\end{enumerate}
\end{corrige}

\begin{corrige}[Exercice 2 - Citation Alain]
\textbf{Analyse.} La science (et la technique qui en découle) augmente notre \emph{puissance} matérielle : transformer la matière, guérir, communiquer, détruire. C'est le « pouvoir sur les choses ». Mais cette puissance est aveugle : elle ne dit pas \emph{comment} l'utiliser à bon escient. La philosophie, par l'exercice critique (Alain : « effort pour penser juste »), permet de maîtriser ses passions, ses préjugés, ses intérêts immédiats pour agir selon la raison et la justice. C'est le « pouvoir sur nous-mêmes » (autonomie).
\textbf{Exemple camerounais.} Un ingénieur en génie civil découvre un vice de construction sur un immeuble en chantier à Douala (pouvoir sur la chose : il sait le corriger). La pression du promoteur est forte pour cacher le défaut (pulsion, intérêt). La philosophie (éthique professionnelle, devoir de sécurité, responsabilité vis-à-vis des futurs habitants) lui donne la force intérieure de refuser la corruption et de signaler le danger (pouvoir sur soi).
\end{corrige}

\coursec{Qu'est-ce que la philosophie ?}

\courssub{Définitions de la philosophie : rappel et mise en perspective}

\begin{aretenir}[Les quatre visages de la définition philosophique]
Avant d'interroger l'origine et l'objet, rappelons les quatre acceptions majeures qui structurent l'identité de la discipline :
\begin{itemize}
    \item \textbf{Étymologique :} \cle{Philo-sophia} (amour de la sagesse) -- \emph{Pythagore}, \emph{Socrate}. La philosophie comme \emph{quête} inachevée, non possession.
    \item \textbf{Systématique :} \cle{Science universelle} -- \emph{Aristote} (\emph{Métaphysique}), \emph{Descartes} (\emph{Principes}, préface). La philosophie comme savoir premier, fondement des sciences particulières (\cle{premières causes}, \cle{premiers principes}).
    \item \textbf{Critique :} \cle{Critique des préjugés} et de l'illusion -- \emph{Kant} (\emph{Critique de la raison pure}), \emph{Nietzsche} (\emph{Généalogie de la morale}). La philosophie comme \cle{éclaircissement} des limites de la raison.
    \item \textbf{Analytique :} \cle{Clarification des concepts} -- \emph{Wittgenstein} (\emph{Tractatus}, \emph{Investigations}), école analytique. La philosophie comme activité de \emph{désembrouillage} logique du langage.
\end{itemize}
Ces définitions ne s'excluent pas : elles dessinent le \textbf{champ de tension} où s'exerce la philosophie -- entre \emph{quête de sens} (sagesse), \emph{exigence de rigueur} (science), \emph{rupture avec l'évidence} (critique) et \emph{précision du dire} (analyse).
\end{aretenir}

\courssub{Origine et objet de la philosophie}

\paragraph{La naissance de la philosophie en Grèce : Milet et Thalès}

\paragraph{Le contexte historique : le miracle grec.}
La philosophie naît officiellement à \cle{Milet}, cité ionienne d'Asie Mineure (actuelle Turquie), vers le \textsc{vi}\up{e} siècle av. J.-C. Le premier philosophe de l'histoire, \cle{Thalès de Milet} (v. 625--547 av. J.-C.), marque la rupture. Mais cette naissance n'est pas un accident : elle repose sur trois conditions matérielles et politiques réunies pour la première fois dans la \cle{cité grecque} (\emph{polis}).

\begin{itemize}
    \item \textbf{Le loisir (\cle{skholè}) :} Contrairement aux sociétés où l'homme est entièrement absorbé par la survie, la cité grecque, reposant sur l'esclavage et une agriculture relativement prospère, libère du temps pour une classe de citoyens. Ce \emph{loisir} n'est pas oisiveté : c'est la disponibilité d'esprit pour contempler le monde sans utilité immédiate. \cle{Aristote} le formulera plus tard : « C'est lorsqu'on a pourvu aux besoins de la vie qu'on commence à philosopher » (\emph{Métaphysique}, A, 2, 982b).
    \item \textbf{L'écriture alphabétique :} L'adoption de l'alphabet phénicien (voyelles comprises) permet de fixer la pensée, de la relire, de la critiquer, de la transmettre sans déformation. Le \cle{logos} écrit devient un objet d'analyse autonome, indépendant de celui qui le prononce.
    \item \textbf{La cité démocratique et le débat public (\cle{agora}) :} À Milet comme à Athènes plus tard, les décisions politiques se prennent par la discussion, la persuasion (\emph{peithô}) et le vote, non par l'ordre du roi. Le citoyen doit argumenter, définir la justice, le bien commun. La philosophie naît comme l'exercice théorique de cette pratique citoyenne : chercher la vérité par la raison partagée, non par l'autorité.
\end{itemize}

\begin{savaistu}[Thalès et l'éclipse de 585 av. J.-C.]
Selon Hérodote (\emph{Histoires}, I, 74), Thalès aurait prédit l'éclipse solaire du 28 mai 585 av. J.-C. qui interrompit une bataille entre Mèdes et Lydiens. Qu'il ait utilisé des cycles babyloniens (Saros) ou une intuition géométrique, l'essentiel est symbolique : \textbf{le ciel n'est plus le domaine exclusif des dieux capricieux, il devient lisible par la raison humaine}. C'est l'acte de naissance de la \cle{rationalité scientifique} et philosophique : le monde est un \cle{cosmos} (ordre), non un \cle{chaos} mythique.
\end{savaistu}

\begin{popfigure}
\begin{tikzpicture}[scale=0.9, every node/.style={font=\small}]
% Axe principal
\draw[popdark, thick, ->] (0,0) -- (15,0) node[right] {Temps};
% Graduations
\foreach \x/\label in {0/VIII\textsuperscript{e} av. J.-C., 3/VII\textsuperscript{e}, 6/VI\textsuperscript{e} (Milet), 9/V\textsuperscript{e}, 12/IV\textsuperscript{e} (Athènes)} {
    \draw (\x, -0.1) -- (\x, 0.1) node[below=2pt, align=center] {\label};
}
% Zones
\fill[poporange!15] (0,0.2) rectangle (5.5,1.8);
\node[poporange, align=center, font=\bfseries] at (2.7,1.2) {Ère \cle{Mythique}\\(Homère, Hésiode)\\\textit{Explication par les dieux,}\\\textit{arbitraire, récit}};
\fill[popblue!15] (5.5,0.2) rectangle (15,1.8);
\node[popblue, align=center, font=\bfseries] at (10.2,1.2) {Ère \cle{Rationnelle} (\cle{Logos})\\\textit{Explication par causes naturelles,}\\\textit{lois universelles, démonstration}};
% Marqueur Thalès
\draw[popgreen, very thick] (6,0) -- (6,2.2) node[above, popgreen, font=\bfseries] {Thalès de Milet (v. 585)};
\node[popgreen, align=center] at (6,2.5) {Premier \cle{physiologue}\\\cle{Archè} = Eau};
% Marqueur Socrate/Platon/Aristote
\draw[poppurple, very thick, dashed] (12,0) -- (12,2.2) node[above, poppurple, font=\bfseries] {Socrate / Platon / Aristote};
\node[poppurple, align=center] at (12,2.5) {Philosophie\\maturity};
% Flèche transition
\draw[popdark, decorate, decoration={snake, amplitude=1pt, segment length=4pt}, thick] (5.5, -0.8) -- (5.5, -1.5) node[below, align=center, font=\itshape] {Passage\\Mythe $\to$ Logos};
\end{tikzpicture}
\legende{Frise chronologique : du mythe au logos. La philosophie naît à Milet (VI\textsuperscript{e} s.) avec Thalès, marquant le passage de l'explication mythique (arbitraire divin) à l'explication rationnelle (causes naturelles, recherche de l'archè).}
\end{popfigure}

\paragraph{Les racines subjectives : l'étonnement et le doute}

Si les conditions sociales expliquent la \emph{possibilité} de la philosophie, qu'en est-il de sa \emph{nécessité} intérieure ?

\paragraph{L'étonnement (\cle{thaumazein}) : le point de départ.}
Pour \cle{Platon} (\emph{Théétète}, 155d), « c'est par l'étonnement que les hommes commencent à philosopher, maintenant comme jadis ». \cle{Aristote} confirme (\emph{Métaphysique}, A, 2) : « C'est par l'étonnement que les hommes ont commencé à philosopher, d'abord devant les difficultés évidentes, puis peu à peu ils ont soulevé des problèmes plus grands. »

\begin{definition}[L'étonnement philosophique (\emph{thaumazein})]
Ce n'est pas la simple surprise face à l'inattendu (comme un tour de magie), ni la peur superstitieuse. C'est un \textbf{étonnement intellectuel et métaphysique} : la prise de conscience brutale que \textbf{le monde n'est pas transparent}, que l'évidence du quotidien cache une opacité fondamentale. L'étonnement philosophique surgit quand \emph{ce qui va de soi ne va plus de soi}. Il transforme l'habitude en question.
\end{definition}

\begin{exemplebox}[Situations d'étonnement dans la vie camerounaise]
\begin{itemize}
    \item \textbf{La mort d'un proche :} Le corps est là, identique, mais \emph{la vie} a disparu. Qu'est-ce qui anime la matière ? Où va la conscience ? L'étonnement face à l'énigme de la mort fonde la métaphysique de l'âme.
    \item \textbf{Les inégalités sociales :} Pourquoi certains naissent dans le luxe à Bastos ou Akwa, d'autres dans la précarité à Nkolbisson ou au village ? L'évidence de l'injustice suscite l'étonnement éthique : \emph{Qu'est-ce que la justice ?} La société est-elle naturelle ou conventionnelle ?
    \item \textbf{Les phénomènes naturels :} Le tonnerre qui gronde au-dessus du Mont Cameroun, la pluie qui tombe ou ne tombe pas en saison sèche. Avant la météorologie, l'étonnement face à la régularité ou la violence du ciel pousse à chercher une \cle{cause} (archè) : est-ce le dieu du tonnerre (mythe) ou un mouvement d'air chaud/froid (logos) ?
\end{itemize}
\end{exemplebox}

\paragraph{Le doute et l'expérience de l'incertitude.}
L'étonnement ouvre la brèche ; le \cle{doute} y entre. Se rendre compte que nos opinions (\emph{doxa}) sont fragiles, que les sens trompent (un bâton courbé dans l'eau), que la tradition se contredit, c'est faire l'expérience de l'\cle{incertitude}. Le doute n'est pas nihilisme : c'est le \textbf{moteur de la recherche}. Comme le dira plus tard \cle{Descartes}, douter, c'est déjà penser. La philosophie naît du refus de l'évidence non examinée.

\begin{attention}[Piège fréquent : l'universalité de la naissance]
\emph{Erreur :} « La philosophie est née partout en même temps (Chine, Inde, Grèce). »
\emph{Correction :} Si la \textbf{sagesse} (réflexion sur la conduite de vie) existe en Chine (Confucius, Lao-Tseu) et en Inde (Upanishads, Bouddha) au même moment (\textsc{vi}\up{e}--\textsc{v}\up{e} s., « âge axial » selon Jaspers), la \textbf{philosophie} au sens strict — \emph{recherche théorique de la vérité par le logos autonome, démonstration rationnelle, concept — a un berceau unique : la Grèce ionienne}. Confucius dit « J'ai transmis, je n'ai pas créé » ; Thalès dit « L'eau est le principe de toutes choses » et tente de le \emph{prouver}. La différence est épistémologique : \textbf{argumentation rationnelle vs tradition aphoristique}.
\end{attention}

\paragraph{L'objet de la philosophie : la totalité du réel}

Si la méthode est le \cle{logos} (raison critique, conceptuelle), quel est le \cle{champ} d'investigation ? Contrairement aux sciences particulières (physique, biologie, sociologie, droit) qui découpent le réel en \cle{objets partiels} (la matière vivante, les structures sociales, les normes juridiques), la philosophie vise la \cle{totalité}.

\begin{propriete}[L'objet philosophique : la totalité et les questions ultimes]
La philosophie a pour objet l'ensemble des \textbf{questions fondamentales} que l'homme se pose sur sa condition et sur l'univers. On les regroupe classiquement en grands domaines :
\begin{center}
\begin{tabularx}{\linewidth}{|l|X|}\hline
\rowcolor{poppurpleL} \textbf{Domaine} & \textbf{Questions types (Objets partiels de la totalité)} \\ \hline
\cle{L'Homme (Anthropologie philosophique)} & Qu'est-ce que l'homme ? Suis-je un corps, une âme, une conscience, un cerveau ? Suis-je libre ou déterminé ? \\ \hline
\cle{Le Monde (Métaphysique / Ontologie / Cosmologie)} & Pourquoi y a-t-il quelque chose plutôt que rien ? Le monde a-t-il un sens, une finalité ? Qu'est-ce que l'être ? \\ \hline
\cle{Dieu (Théodicée / Philosophie de la religion)} & Dieu existe-t-il ? Quel rapport entre foi et raison ? Le mal est-il compatible avec un Dieu bon et tout-puissant ? \\ \hline
\cle{La Vérité (Épistémologie / Logique)} & Qu'est-ce que savoir ? La vérité est-elle absolue ou relative ? Comment distinguer science et opinion ? \\ \hline
\cle{La Justice / Le Droit (Philosophie politique)} & Qu'est-ce qu'une société juste ? L'État a-t-il tous les droits ? Loi naturelle vs loi positive ? \\ \hline
\cle{La Liberté (Éthique / Métaphysique)} & Suis-je auteur de mes actes ? Responsabilité vs déterminisme (social, biologique, divin). \\ \hline
\cle{Le Bonheur (Éthique / Morale)} & Quel est le bien suprême ? Plaisir, vertu, devoir, accomplissement ? Le bonheur est-il le but de la vie ? \\ \hline
\end{tabularx}
\end{center}
\end{propriete}

\begin{aretenir}[La spécificité de l'objet philosophique : la Totalité]
\begin{itemize}
    \item \textbf{Non-spécialisation :} La philosophie ne se contente pas d'étudier \emph{un} aspect du réel (le vivant, le social, le langage), elle interroge \emph{les liens} entre tous les aspects et le \emph{sens de l'ensemble}.
    \item \textbf{Questions « ultimes » :} Elle pose les questions que les autres sciences \emph{présupposent} sans les démontrer (ex. : « La nature a-t-elle des lois ? » pour la physique ; « L'homme est-il libre ? » pour le droit ; « La vie a-t-elle un sens ? » pour la biologie).
    \item \textbf{Secondarité réflexive :} La philosophie pense \emph{sur} les autres savoirs (philosophie des sciences, de l'histoire, de l'art, du droit, de la technique). Elle est une \cle{méta-connaissance}.
\end{itemize}
\end{aretenir}

\courssub{La méthode de la philosophie : le logos critique}

Le passage du mythe au logos est le \cle{geste fondateur} de la philosophie. Il ne s'agit pas de nier le mythe, mais de changer de \cle{régime d'explication}.

\begin{definition}[Le \cle{Logos} (parole raisonnée, discours rationnel)]
Étymologiquement : \emph{parole, discours, raison, ratio}. En philosophie, le \cle{logos} désigne :
\begin{enumerate}
    \item Le \textbf{discours structuré} : énoncé articulé, syntaxiquement correct, susceptible de vrai ou de faux.
    \item La \textbf{raison} : la faculté de lier les concepts, d'inférer, de démontrer.
    \item Le \textbf{principe d'intelligibilité du monde} : le cosmos est \emph{logique}, il a une structure rationnelle accessible à l'esprit humain (Héraclite : « Tout se passe selon le Logos »).
\end{enumerate}
Le logos s'oppose au \cle{mythos} (récit traditionnel, autorité du poète/prêtre, explication par des volontés divines anthropomorphes, arbitraires).
\end{definition}

\begin{popfigure}
\begin{tikzpicture}[scale=0.95, every node/.style={font=\small}]
% Deux colonnes
% Colonne Gauche : MYTHE
\begin{scope}[xshift=0cm]
\node[poporange, font=\bfseries\large, align=center] at (3.5, 6.5) {Explication \cle{Mythique} (\emph{Mythos})};
\draw[poporange, thick] (0.5,6.2) -- (6.5,6.2);
% Boîte phénomène
\node[draw=poporange, thick, rounded corners, fill=poporange!10, align=center, minimum width=5cm, minimum height=1cm] (phenom) at (3.5, 5.3) {Phénomène : \textbf{La Foudre / Le Tonnerre}};
% Flèche
\draw[poporange, ->, thick] (phenom.south) -- ++(0,-0.5) node[below=0.2cm, align=center, font=\itshape] (q1) {Pourquoi ?};
% Réponse Mythique
\node[draw=poporange, thick, rounded corners, fill=poporange!5, align=left, minimum width=5.5cm, minimum height=3.5cm] (myth) at (3.5, 3.5) {
\textbf{Récit :} Zeus (ou Heuso, dieu du tonnerre Bamiléké) est en colère.\\[4pt]
\textbf{Cause :} Volonté divine, arbitraire, personnelle.\\[4pt]
\textbf{Mode :} Anthropomorphisme (dieux = hommes puissants).\\[4pt]
\textbf{Validation :} Autorité de la tradition, du prêtre, du devin.\\[4pt]
\textbf{Finalité :} Apaiser le dieu (sacrifice, rituel).
};
\node[poporange, font=\bfseries] at (3.5, 1.5) {Savoir : \cle{Opinion} (\emph{Doxa}) / Croyance};
\end{scope}
% Colonne Droite : LOGOS
\begin{scope}[xshift=8.5cm]
\node[popblue, font=\bfseries\large, align=center] at (3.5, 6.5) {Explication \cle{Rationnelle} (\emph{Logos})};
\draw[popblue, thick] (0.5,6.2) -- (6.5,6.2);
% Boîte phénomène
\node[draw=popblue, thick, rounded corners, fill=popblue!10, align=center, minimum width=5cm, minimum height=1cm] (phenom2) at (3.5, 5.3) {Phénomène : \textbf{La Foudre / Le Tonnerre}};
% Flèche
\draw[popblue, ->, thick] (phenom2.south) -- ++(0,-0.5) node[below=0.2cm, align=center, font=\itshape] (q2) {Comment ? / Quelle cause ?};
% Réponse Rationnelle (Thalès / Anaximandre / Science moderne)
\node[draw=popblue, thick, rounded corners, fill=popblue!5, align=left, minimum width=5.5cm, minimum height=3.5cm] (logos) at (3.5, 3.5) {
\textbf{Hypothèse :} Séparation de charges électriques dans les nuages (Thalès : électricité statique / ambre).\\[4pt]
\textbf{Cause :} Mécanisme naturel, impersonnel, universel (lois physique).\\[4pt]
\textbf{Mode :} Concepts abstraits (charge, potentiel, champ), mathématisation.\\[4pt]
\textbf{Validation :} Observation, expérience, démonstration, prévision.\\[4pt]
\textbf{Finalité :} Comprendre, prévoir, se protéger (paratonnerre).
};
\node[popblue, font=\bfseries] at (3.5, 1.5) {Savoir : \cle{Connaissance} (\emph{Epistémè}) / Science};
\end{scope}
% Flèche de transition centrale
\draw[popdark, <->, very thick, decorate, decoration={snake, amplitude=2pt, segment length=6pt}] (7.0, 4.0) -- (8.0, 4.0) node[midway, above, font=\bfseries\small] {Rupture épistémologique};
\end{tikzpicture}
\legende{Schéma comparatif : même phénomène (la foudre), deux régimes d'explication. Le mythe raconte une histoire à visée pratique (apaiser) ; le logos propose un modèle causal à visée théorique (comprendre).}
\end{popfigure}

\begin{methode}[La démarche philosophique (Logos) en 4 temps]
La méthode philosophique ne se réduit pas à « donner son avis ». Elle suit un parcours rigoureux :
\begin{enumerate}
    \item \textbf{Problématiser :} Transformer un thème en \cle{problème} (tension entre thèses opposées). Ex. : « La technique libère-t-elle ou aliène-t-elle l'homme ? »
    \item \textbf{Conceptualiser :} Définir rigoureusement les termes (technique, liberté, aliénation) pour éviter les équivoques.
    \item \textbf{Argumenter :} Construire un \cle{raisonnement} (déductif, inductif, analogique, dialectique) en articulant \emph{thèses}, \emph{arguments}, \emph{exemples} et \emph{contre-arguments}.
    \item \textbf{Conclure provisoirement :} Apporter une réponse nuancée, consciente de ses limites, qui rouvre le questionnement.
\end{enumerate}
\end{methode}

\courssub{L'importance de la philosophie}

\begin{aretenir}[Pourquoi philosopher ? Trois ordres de nécessité]
\begin{itemize}
    \item \textbf{Nécessité existentielle (Sagesse) :} « Soigner l'âme » (Socrate, Épicure, Stoïciens). Donner du sens à la vie, à la souffrance, à la mort. Apprendre à \emph{bien vivre} et \emph{bien mourir}.
    \item \textbf{Nécessité citoyenne (Démocratie) :} Former l'esprit critique, résister à la manipulation (propagande, fake news, populisme), fonder le \emph{vivre-ensemble} sur la discussion rationnelle (\emph{Habermas}) et non sur la force. \emph{Ex. camerounais :} Débat sur le code électoral, la décentralisation, le bilinguisme.
    \item \textbf{Nécessité épistémique (Fondement des sciences) :} Clarifier les concepts (espace, temps, cause, vie, information), épurer les présupposés, arbitrer les conflits inter-disciplinaires (ex. : neurosciences vs phénoménologie sur la conscience). La philosophie est la \cle{conscience réflexive} des sciences.
\end{itemize}
\end{aretenir}

\courssub{Philosophie et science : distinction et articulation}

\begin{exemplebox}[Récits traditionnels camerounais vs Explication rationnelle : L'origine du monde]
\begin{center}
\begin{tabularx}{\linewidth}{|l|X|X|}\hline
\rowcolor{popblueL} \textbf{Aspect} & \textbf{Récit Traditionnel (Ex. Bassa, Bamoun, Bamiléké)} & \textbf{Explication Rationnelle (Cosmologie / Biologie)} \\ \hline
\textbf{Origine} & Un Être Suprême (Njambi, Si, Ngwè) façonne le monde, souvent à partir d'un chaos primordial (œuf, eau, vide). & \cle{Big Bang} (expansion d'une singularité il y a 13,8 Ga) ; formation solaire par accrétion (4,5 Ga). \\ \hline
\textbf{Apparition de l'homme} & Modelage dans l'argile / glaise, souffle de vie, sortie d'un trou / grotte / arbre sacré. & \cle{Évolution} : sélection naturelle, ancêtres communs avec les grands singes (\emph{Homo sapiens} $\approx$ 300 000 ans). \\ \hline
\textbf{Statut du récit} & \cle{Vérité sacrée}, fondatrice de l'identité, transmissible par initiation/oral. & \cle{Hypothèse scientifique} falsifiable, révisable, validée par preuves (fossiles, ADN, rayonnement fossile). \\ \hline
\textbf{Question philosophique} & \emph{Quel est le sens de cette création ? Quelle est ma place dans l'ordre voulu par l'Ancêtre ?} & \emph{Qu'est-ce que la matière ? Le hasard a-t-il un sens ? Qu'est-ce que la conscience émergeant du cerveau ?} \\ \hline
\end{tabularx}
\end{center}
\emph{Note :} La philosophie ne tranche pas « pour » ou « contre » le mythe ; elle \textbf{interroge le statut de vérité} de l'un et de l'autre, et le rapport entre \cle{mythe fondateur} (sens social) et \cle{logos explicatif} (validité universelle).
\end{exemplebox}

\begin{propriete}[Rapports Philosophie / Sciences : ni confusion, ni séparation]
\begin{center}
\begin{tabularx}{\linewidth}{|l|X|X|}\hline
\rowcolor{popgreenL} \textbf{Critère} & \textbf{Philosophie} & \textbf{Sciences (particulières)} \\ \hline
\textbf{Objet} & \cle{Totalité} / Questions ultimes (Sens, Valeur, Être) & \cle{Partiel} / Domaine délimité (Vivant, Matière, Social) \\ \hline
\textbf{Méthode} & \cle{Logos} : Conceptualisation, argumentation, critique a priori & \cle{Méthode expérimentale} : Observation, hypothèse, expérience, modélisation \\ \hline
\textbf{Vérité} & \cle{Universalisée} par la raison (valide pour tout sujet rationnel) & \cle{Objectivée} par le fait (valide par reproductibilité intersubjective) \\ \hline
\textbf{Statut} & \cle{Seconde} (réfléchit sur les présupposés des sciences) & \cle{Première} (produit des connaissances positives sur le réel) \\ \hline
\textbf{Exemple} & « Qu'est-ce que le temps ? » (Bergson, Heidegger) & « Comment mesurer le temps ? » (Horlogerie atomique, Relativité) \\ \hline
\end{tabularx}
\end{center}
\emph{Articulation :} La philosophie \textbf{conditionne} les sciences (clarification des concepts, éthique de la recherche) et les sciences \textbf{alimentent} la philosophie (nouvelles données pour repenser l'homme, la nature, la technique).
\end{propriete}

\courssub{Synthèse et application à la vie courante}

\begin{exoresolu}[Analyse philosophique d'une situation concrète : « La mort soudaine du chef du village »]
\textbf{Énoncé :} Dans un village de l'Ouest-Cameroun, le chef supérieur décède brutalement d'une crise cardiaque. La foule murmure : « C'est la sorcellerie du clan rival ! » Le médecin légiste parle « d'infarctus du myocarde ». Le notable traditionnel ordonne des rites d'apaisement des ancêtres. Le fils aîné, étudiant en philosophie à l'université, reste silencieux.

\textbf{Question :} Identifiez dans cette situation les trois attitudes : \emph{mythique}, \emph{scientifique}, \emph{philosophique}. Justifiez.

\tcblower
\textbf{Solution détaillée :}

\begin{enumerate}
    \item \textbf{Attitude mythique (La foule / Notable) :} « C'est la sorcellerie / Les ancêtres sont fâchés. »
        \begin{itemize}
            \item \emph{Type d'explication :} Cause intentionnelle, invisible, personnelle (volonté d'un agent surnaturel).
            \item \emph{Validation :} Autorité de la tradition, consensus social, efficacité rituelle (apaiser le trouble social).
            \item \emph{Fonction :} Donner un \emph{sens moral et social} à l'absurde de la mort ; maintenir la cohésion du groupe par la désignation d'un coupable (bouc émissaire).
        \end{itemize}
    \item \textbf{Attitude scientifique (Le médecin) :} « Infarctus du myocarde, athérosclérose, hypertension non traitée. »
        \begin{itemize}
            \item \emph{Type d'explication :} Cause naturelle, impersonnelle, universelle (lois de la physiologie, facteurs de risque).
            \item \emph{Validation :} Preuves objectives (autopsie, ECG, biologie), reproductibilité, prévisibilité statistique.
            \item \emph{Fonction :} Expliquer le \emph{mécanisme} (comment), permettre la prévention (hygiène de vie), soigner les vivants.
        \end{itemize}
    \item \textbf{Attitude philosophique (L'étudiant silencieux / La philosophie) :} Elle ne choisit pas « pour » ou « contre ». Elle \emph{problématise} :
        \begin{itemize}
            \item \emph{Sur le savoir :} Pourquoi deux explications incompatibles coexistent-elles ? (Épistémologie : pluralisme des régimes de vérité).
            \item \emph{Sur l'homme :} La mort réduit-elle l'homme à un corps biologique ? Qu'en est-il de sa \emph{dignité}, de son \emph{nom}, de sa \emph{fonction symbolique} ? (Anthropologie philosophique).
            \item \emph{Sur la justice :} Accuser le clan rival sans preuve est-il juste ? La « vérité » judiciaire (procès de sorcellerie) vaut-elle la vérité scientifique ? (Philosophie du droit / Éthique).
            \item \emph{Sur le sens :} La mort du chef a-t-elle un \emph{sens} pour la communauté au-delà de sa cause biologique ? (Métaphysique / Philosophie de la culture).
        \end{itemize}
    \item \textbf{Synthèse :} La philosophie naît ici de l'\cle{étonnement} face au choc des interprétations. Elle ne nie ni la douleur culturelle (mythe), ni la réalité biologique (science), mais elle \textbf{interroge les conditions de possibilité, de validité et de coexistence de ces vérités}. C'est l'exercice du \cle{logos} sur la \cle{totalité} de la situation.
\end{enumerate}
\end{exoresolu}

\begin{exercice}[Application personnelle]
Choisissez un événement marquant de votre vie ou de l'actualité camerounaise récente (ex. : éboulement à Dschang, réussite/échec au Probatoire, conflit foncier, innovation technique locale).
\begin{enumerate}
    \item Décrivez l'explication \emph{spontanée / mythique / senso-comune} qu'on en donne souvent.
    \item Décrivez l'explication \emph{scientifique / technique / juridique} qu'en donnent les experts.
    \item Formulez \textbf{trois questions philosophiques} distinctes (sur l'homme, la vérité, la justice, la liberté, le bonheur, le sens) que cet événement vous inspire.
\end{enumerate}
\end{exercice}

\begin{corrige}[Exercice n°1 -- Pistes de correction]
\emph{Exemple : Échec au Probatoire.}
\begin{enumerate}
    \item \textbf{Mythique / Spontanée :} « C'est la malchance / le mauvais sort / les profs ne m'aiment pas / le village m'a jeté un sort. » (Cause intentionnelle, externe, non contrôlable).
    \item \textbf{Scientifique / Technique :} « Notes insuffisantes dans les matières à coefficient élevé / mauvaise gestion du temps / lacunes méthodologiques / stress non géré. » (Causes factuelles, internes/externes, mesurables, corrigeables).
    \item \textbf{Questions philosophiques :}
        \begin{enumerate}
            \item (Liberté/Responsabilité) \emph{Suis-je pleinement responsable de mon échec, ou suis-je déterminé par mon milieu social, mon QI, mes professeurs ?}
            \item (Vérité/Justice) \emph{L'examen est-il un critère juste de la valeur d'un homme ? La note mesure-t-elle l'intelligence ou la conformité ?}
            \item (Bonheur/Sens) \emph{Cet échec détruit-il ma vie ou peut-il être une épreuve fondatrice (résilience) ? Le bonheur dépend-il de la réussite sociale ?}
        \end{enumerate}
\end{enumerate}
\emph{L'essentiel n'est pas la réponse, mais la mise en œuvre du \cle{logos} (questionnement rationnel) sur la \cle{totalité} de l'expérience vécue.}
\end{corrige}

\coursec{La méthode de la philosophie}

Nous avons vu, dans la section précédente, que la philosophie est née de l'\cle{étonnement} et qu'elle a pour objet la recherche de la vérité sur l'homme, le monde et les valeurs. Mais une question demeure : comment le philosophe procède-t-il concrètement pour mener cette recherche ? On ne philosophe pas n'importe comment, en disant tout ce qui nous passe par la tête. La philosophie possède une \cle{méthode}, c'est-à-dire une manière ordonnée et rigoureuse de conduire sa pensée. C'est cette méthode que nous allons maintenant étudier pas à pas, car elle constitue l'outil principal de l'élève de Première technique pour réussir la dissertation au Probatoire et au Baccalauréat technique.

\courssub{Notion de méthode}

Commençons par une intuition simple. Un mécanicien qui répare un moteur ne démonte pas les pièces au hasard : il suit un ordre précis, diagnostique la panne, teste les pièces une à une, puis remonte. De même, un élève en atelier de menuiserie ne coupe pas le bois avant d'avoir mesuré et tracé. Dans tous les métiers techniques, il existe une \emph{procédure} à respecter pour obtenir un bon résultat. La pensée philosophique obéit à la même exigence.

\begin{definition}[La méthode philosophique]
La \cle{méthode} (du grec \emph{methodos}, « chemin, poursuite ») est un \textbf{ensemble ordonné de démarches et de règles} permettant de conduire correctement sa pensée vers la vérité. La méthode philosophique est la démarche rationnelle par laquelle le philosophe pose un problème, examine les opinions, définit les notions, argumente et conclut de manière justifiée.
\end{definition}

La méthode philosophique comporte classiquement \textbf{cinq étapes} que nous allons détailler : le questionnement, le doute, l'analyse, l'argumentation et la synthèse.

\begin{popfigure}
\begin{center}
\begin{tikzpicture}[
  node distance=0.55cm,
  etape/.style={rectangle, rounded corners=4pt, draw=popblue, fill=popblueL, thick, text width=10.5cm, align=center, font=\small, minimum height=0.9cm},
  fleche/.style={-{Stealth[length=3mm]}, thick, popdark}
]
\node[etape] (e1) {\textbf{Étape 1 — L'étonnement et le questionnement} : poser un problème};
\node[etape, below=of e1] (e2) {\textbf{Étape 2 — Le doute méthodique} : mettre en question les opinions reçues};
\node[etape, below=of e2] (e3) {\textbf{Étape 3 — L'analyse des notions} : définir clairement les termes};
\node[etape, below=of e3] (e4) {\textbf{Étape 4 — L'argumentation} : avancer des raisons, examiner les objections};
\node[etape, below=of e4] (e5) {\textbf{Étape 5 — La synthèse} : conclure de manière nuancée et justifiée};
\draw[fleche] (e1) -- (e2);
\draw[fleche] (e2) -- (e3);
\draw[fleche] (e3) -- (e4);
\draw[fleche] (e4) -- (e5);
\end{tikzpicture}
\end{center}
\legende{Organigramme des cinq étapes de la démarche philosophique. Chaque étape prépare la suivante : on ne peut pas argumenter (étape 4) sans avoir défini les termes (étape 3), ni douter utilement (étape 2) sans avoir d'abord posé un problème (étape 1).}
\end{popfigure}

\courssub{Étape 1 : l'étonnement et le questionnement}

Toute démarche philosophique commence par un \cle{étonnement} : quelque chose nous surprend, nous intrigue, cesse d'aller de soi. Un élève de Yaoundé constate que deux camarades punis pour la même faute n'ont pas reçu la même sanction : « Est-ce juste ? » se demande-t-il. Cet étonnement se transforme alors en \cle{question}, puis en \cle{problème}.

\begin{definition}[Le problème philosophique]
Un \cle{problème} est une question précise, formulée clairement, qui met en tension des réponses possibles mais contradictoires. Poser un problème, c'est passer de l'étonnement confus (« c'est bizarre ») à une question bien formulée (« faut-il toujours obéir à la loi pour être juste ? »).
\end{definition}

L'art de philosopher commence donc par l'art de \textbf{bien poser les questions}. Une bonne question philosophique est générale (elle ne concerne pas seulement un cas particulier), elle porte sur un concept (la justice, la liberté, la vérité) et elle appelle une réponse justifiée, pas une simple information.

\courssub{Étape 2 : le doute méthodique (Descartes)}

Une fois le problème posé, la tentation est grande de répondre immédiatement avec ce que l'on a toujours entendu dire : les opinions de la famille, du quartier, des réseaux sociaux, les proverbes, les préjugés. Or ces opinions reçues ne sont pas nécessairement vraies. Le philosophe français René \textbf{Descartes} (1596--1650) a donc proposé de commencer par \emph{douter}.

\begin{definition}[Le doute méthodique]
Le \cle{doute méthodique} est la démarche par laquelle on suspend provisoirement son adhésion à toutes les opinions reçues, afin d'examiner si elles reposent sur des raisons solides. Chez Descartes, ce doute est \textbf{provisoire} (on doute pour mieux fonder, non pour rester dans le doute) et \textbf{volontaire} (c'est un outil de travail, non un état d'esprit sceptique).
\end{definition}

Descartes compare ce doute au tri d'un panier de fruits : pour écarter les fruits pourris, on vide tout le panier, puis on ne remet que les fruits sains après examen. De même, le philosophe « vide » son esprit de ses opinions pour ne retenir que celles qui résistent à l'examen rationnel.

\begin{attention}[Ne pas confondre doute méthodique et scepticisme]
Le doute méthodique n'est pas le scepticisme. Le \emph{sceptique} doute définitivement et affirme qu'on ne peut rien connaître ; celui qui pratique le doute méthodique doute \emph{temporairement}, comme on nettoie un outil avant de s'en servir, précisément \textbf{pour} atteindre une connaissance certaine. Au Bac, écrire « Descartes doute de tout donc il ne sait rien » est une erreur d'interprétation.
\end{attention}

\courssub{Étape 3 : l'analyse des notions}

Beaucoup de disputes sont de simples malentendus : on se dispute parce qu'on n'entend pas la même chose par le même mot. Si deux élèves débattent de la « liberté » alors que l'un pense à la liberté de mouvement et l'autre à la liberté d'expression, ils ne pourront jamais se comprendre. D'où la troisième étape : \textbf{définir clairement les termes du problème}.

Analyser une notion, c'est :
\begin{enumerate}
\item dégager son \emph{sens ordinaire} (ce que le mot veut dire dans la langue courante) ;
\item distinguer ses \emph{sens différents} et lever les ambiguïtés ;
\item proposer une \emph{définition précise} qui servira de référence pour toute la réflexion.
\end{enumerate}

\begin{exemplebox}[Analyser la notion de « vérité »]
Sens ordinaire : dire ce qui est. Sens différents : la vérité comme exactitude d'un fait (« il est 14 h »), la vérité comme sincérité (« dire ce que l'on pense »), la vérité scientifique (vérifiable par l'expérience). Définition de travail retenue : la vérité est l'accord entre ce que l'on dit et ce qui est réellement. Cette définition permet ensuite d'argumenter sans confusion.
\end{exemplebox}

\courssub{Étape 4 : l'argumentation}

C'est le cœur de la méthode. \cle{Argumenter}, c'est avancer des \textbf{raisons} pour soutenir une thèse, et non simplement affirmer ce que l'on pense. Une argumentation philosophique comporte :
\begin{itemize}
\item une \textbf{thèse} : la position que l'on défend (« il faut toujours dire la vérité ») ;
\item des \textbf{arguments} : les raisons qui la justifient (« sans vérité, la confiance est impossible ») ;
\item des \textbf{exemples} : des cas concrets qui illustrent les arguments ;
\item l'examen des \textbf{objections} : les arguments de la thèse adverse, auxquels on répond.
\end{itemize}

L'examen des objections est essentiel : philosopher, ce n'est pas écraser l'adversaire, c'est prendre au sérieux les raisons de celui qui pense autrement, pour tester la solidité de sa propre position. C'est ce mouvement thèse -- objection -- réponse qui fait progresser la pensée.

\courssub{Étape 5 : la synthèse et la conclusion nuancée}

La démarche se termine par une \cle{synthèse} : on rassemble les résultats obtenus et on formule une \textbf{conclusion nuancée}. Nuancée ne veut pas dire vague ou hésitante : cela signifie que la conclusion tient compte de la complexité du problème, reconnaît la part de vérité des positions adverses et précise les conditions dans lesquelles la thèse retenue est valable. Par exemple : « Dire la vérité est un devoir, mais la manière et le moment de la dire doivent tenir compte d'autrui. »

\courssub{Les exigences de la méthode philosophique}

Quatre exigences traversent toutes ces étapes :

\begin{aretenir}[Les quatre exigences de la méthode philosophique]
\begin{itemize}
\item \textbf{La rationalité} : ne s'appuyer que sur des raisons accessibles à tout être pensant, et non sur l'émotion, l'autorité ou la tradition seule.
\item \textbf{L'universalité} : viser des conclusions valables pour tous, partout, et non seulement pour soi ou pour son groupe.
\item \textbf{L'esprit critique} : ne rien accepter sans examen, y compris ses propres opinions.
\item \textbf{La rigueur du langage} : employer des termes précis, définis, sans ambiguïté, car une pensée claire exige un langage clair.
\end{itemize}
\end{aretenir}

\begin{attention}[Erreur fréquente : l'opinion sans argument]
Donner son opinion personnelle sans la justifier n'est \textbf{pas} philosopher. Écrire « moi je pense que mentir c'est mal parce que je n'aime pas ça » relève du sentiment, pas de la philosophie. La philosophie exige de transformer l'opinion (\emph{doxa}) en connaissance justifiée : chaque affirmation doit être accompagnée d'une raison que n'importe qui peut examiner et discuter. À l'examen, une copie qui ne fait qu'exprimer des opinions personnelles, même sincères, ne peut pas obtenir une bonne note.
\end{attention}

\courssub{La maïeutique de Socrate : l'art d'accoucher les esprits}

La méthode philosophique trouve son premier modèle historique chez \textbf{Socrate} (Athènes, v. 470--399 av. J.-C.). Socrate n'écrivait rien : il arpentait les rues d'Athènes et interrogeait ses concitoyens. Sa mère était sage-femme : elle accouchait les corps. Socrate disait pratiquer, lui, la \cle{maïeutique} (du grec \emph{maieutikè}, « art d'accoucher ») : l'art d'\textbf{accoucher les esprits}, c'est-à-dire d'aider chacun, par le dialogue, à mettre au monde la vérité qu'il porte en lui sans le savoir.

La maïeutique repose sur le \textbf{dialogue question--réponse} : Socrate pose une question, l'interlocuteur répond avec son opinion ; Socrate montre alors, par une objection, que cette réponse mène à une contradiction ; l'interlocuteur affine sa réponse ; et ainsi de suite, jusqu'à une définition plus solide.

\begin{popfigure}
\begin{center}
\begin{tikzpicture}[
  node distance=0.5cm,
  q/.style={rectangle, rounded corners=3pt, draw=popblue, fill=popblueL, thick, text width=4.6cm, align=center, font=\small},
  r/.style={rectangle, rounded corners=3pt, draw=poporange, fill=poporangeL, thick, text width=4.6cm, align=center, font=\small},
  fleche/.style={-{Stealth[length=3mm]}, thick, popdark}
]
\node[q] (q1) {\textbf{Question du maître} : « Qu'est-ce que la justice ? »};
\node[r, right=1.2cm of q1] (r1) {\textbf{Réponse de l'élève} : opinion première};
\node[q, below=1.1cm of r1] (q2) {\textbf{Objection} : cas où cette réponse échoue};
\node[r, left=1.2cm of q2] (r2) {\textbf{Réponse affinée} : définition corrigée};
\draw[fleche] (q1) -- (r1);
\draw[fleche] (r1) -- (q2);
\draw[fleche] (q2) -- (r2);
\draw[fleche, dashed] (r2.south) to[out=-90,in=-90] node[below, font=\small\itshape]{nouvelle question, jusqu'à une définition solide} (q1.south);
\end{tikzpicture}
\end{center}
\legende{Le dialogue socratique : par une suite de questions, de réponses et d'objections, l'interlocuteur abandonne ses opinions confuses et parvient lui-même à une définition rigoureuse. La vérité n'est pas imposée de l'extérieur : elle est « accouchée » par l'élève.}
\end{popfigure}

\begin{exemplebox}[Dialogue socratique reconstitué : maître et élève sur la justice]
\textbf{Maître} : Dis-moi, qu'est-ce qu'être juste ?\par
\textbf{Élève} : Être juste, c'est rendre à chacun ce qu'on lui doit.\par
\textbf{Maître} : Bien. Si un ami t'a prêté sa machette et qu'entre-temps il est devenu furieux et dangereux, la justice exige-t-elle de la lui rendre ?\par
\textbf{Élève} : Non... ce serait lui nuire, et nuire aux autres.\par
\textbf{Maître} : Ta première définition est donc insuffisante. Que faut-il ajouter ?\par
\textbf{Élève} : Être juste, c'est rendre à chacun ce qu'on lui doit \emph{en veillant à ne nuire à personne}.\par
\textbf{Maître} : Voilà une définition meilleure. Examinons-la encore...
\end{exemplebox}

\courssub{Application : rédiger une réponse philosophique en cinq étapes}

\begin{methode}[Rédiger une réponse philosophique en cinq étapes]
\begin{enumerate}
\item \textbf{Poser le problème} : transformer le sujet en une question claire et montrer pourquoi elle est difficile (tension entre deux réponses possibles).
\item \textbf{Douter des opinions reçues} : écarter provisoirement les réponses toutes faites (préjugés, proverbes, « on dit que... »).
\item \textbf{Définir les notions} : analyser et définir précisément les termes essentiels du sujet.
\item \textbf{Argumenter} : défendre une thèse par des raisons et des exemples, puis examiner au moins une objection et y répondre.
\item \textbf{Conclure par une synthèse nuancée} : formuler une réponse justifiée qui tient compte de la complexité du problème.
\end{enumerate}
\end{methode}

\begin{experience}[TP guidé : « Faut-il toujours dire la vérité ? »]
\textbf{Situation de départ.} À Douala, une élève de Première, Amina, découvre que sa meilleure amie a triché à une évaluation. Le professeur interroge Amina : sait-elle qui a triché ? Dire la vérité, c'est dénoncer son amie ; se taire, c'est mentir par silence. Faut-il toujours dire la vérité ?\par
\textbf{Étape 1 — Questionnement.} Problème : la vérité est-elle un devoir absolu, ou y a-t-il des situations où il est permis, voire juste, de la taire ?\par
\textbf{Étape 2 — Doute.} Opinions reçues à suspendre : « il faut toujours dire la vérité » (principe reçu) et « il ne faut jamais dénoncer un ami » (loyauté du groupe). Aucune n'est encore justifiée.\par
\textbf{Étape 3 — Analyse.} Définir : la \emph{vérité} (accord du discours avec le réel), le \emph{devoir} (obligation morale), \emph{toujours} (sans aucune exception).\par
\textbf{Étape 4 — Argumentation.} Thèse 1 : il faut dire la vérité, car sans elle ni confiance ni justice ne sont possibles (la tricherie lèse les élèves honnêtes). Objection : la loyauté envers l'amitié est aussi une valeur, et dénoncer peut détruire une personne. Réponse : on peut dire la vérité en protégeant autrui (encourager l'amie à avouer elle-même).\par
\textbf{Étape 5 — Synthèse.} Conclusion nuancée : la vérité est un devoir, mais elle n'exclut ni la délicatesse ni le souci de la justice envers tous ; le silence complice n'est pas une loyauté authentique.
\end{experience}

\begin{exoresolu}[Exercice d'application de la méthode]
\textbf{Énoncé.} Un élève affirme : « La méthode philosophique ne sert à rien : chacun a son opinion et toutes les opinions se valent. » En mobilisant les étapes et les exigences de la méthode philosophique, montrez en quoi cette affirmation est elle-même contradictoire.
\tcblower
\textbf{Solution détaillée.} \emph{Étape 1 (problème)} : toutes les opinions se valent-elles ? \emph{Étape 2 (doute)} : l'élève présente sa propre opinion comme vraie pour tous, donc comme supérieure aux opinions contraires. \emph{Étape 3 (analyse)} : distinguer l'\emph{opinion} (croyance non justifiée) de la \emph{connaissance} (thèse justifiée par des raisons). \emph{Étape 4 (argumentation)} : si toutes les opinions se valaient, alors l'opinion « toutes les opinions ne se valent pas » vaudrait autant que la sienne : sa thèse se détruit elle-même. De plus, affirmer que la méthode ne sert à rien est déjà un jugement qui prétend être rationnellement fondé : il utilise donc la méthode qu'il nie. \emph{Étape 5 (synthèse)} : les opinions ne se valent pas toutes : seules celles qui résistent à l'examen rationnel (doute, définition, argumentation) méritent d'être retenues. L'élève confond le droit d'avoir une opinion avec la valeur de vérité de cette opinion.
\end{exoresolu}

\begin{exercice}[Pour s'entraîner]
Appliquez les cinq étapes de la méthode philosophique à la question suivante, inspirée de la vie scolaire au Cameroun : « La réussite scolaire dépend-elle uniquement du travail de l'élève ? » Rédigez un paragraphe pour chaque étape, en veillant à définir les notions de \emph{réussite} et de \emph{travail}, et à examiner au moins une objection (rôle du milieu familial, des conditions matérielles, du professeur).
\end{exercice}

\begin{corrige}[Exercice]
\emph{Étape 1} : problème — le mérite personnel suffit-il à expliquer la réussite, ou celle-ci dépend-elle aussi de conditions extérieures ? \emph{Étape 2} : opinions reçues à suspendre — « qui veut peut » et « tout dépend de la chance ». \emph{Étape 3} : définitions — la \emph{réussite} (atteinte d'un objectif scolaire mesurable), le \emph{travail} (effort régulier et méthodique). \emph{Étape 4} : thèse — le travail est la condition principale (argument : sans effort, aucun avantage extérieur ne suffit ; exemple : un élève boursier qui ne travaille pas échoue). Objection — les conditions (électricité pour réviser le soir, soutien familial, qualité de l'encadrement) ne sont pas égales pour tous. Réponse — ces conditions influencent la pente, mais ne suppriment pas le rôle décisif de l'effort. \emph{Étape 5} : synthèse nuancée — la réussite résulte de la rencontre entre le travail personnel, qui reste déterminant, et des conditions extérieures qui peuvent le faciliter ou le freiner ; reconnaître ces conditions ne déresponsabilise pas l'élève, mais invite aussi la société à rendre les chances plus égales.
\end{corrige}

Ainsi, la méthode philosophique transforme l'étonnement spontané en réflexion ordonnée : questionner, douter, définir, argumenter, conclure. Mais à quoi sert concrètement cette exigence dans la vie d'un élève, d'un technicien, d'un citoyen camerounais ? C'est la question de l'\emph{importance} de la philosophie, que nous aborderons dans la section suivante.

\coursec{L'importance de la philosophie}

Nous venons de voir que la philosophie possède une méthode propre : le doute méthodique, l'analyse des concepts, l'argumentation rigoureuse et le dialogue. Mais à quoi sert-elle concrètement ? Beaucoup d'élèves, au premier abord, se demandent pourquoi on leur impose une matière qui ne fabrique ni machine, ni médicament, ni application mobile. Cette question, il faut le remarquer, est déjà une question philosophique : elle interroge la \emph{valeur} d'un savoir. Toute cette section va y répondre en montrant que la philosophie est utile à cinq niveaux : individuel, moral, civique, intellectuel et professionnel. Loin d'être un luxe réservé aux intellectuels, elle est une école de la liberté, indispensable à quiconque veut penser par lui-même et vivre en société.

\courssub{Sur le plan individuel : former l'esprit critique et l'autonomie de jugement}

Commençons par une situation quotidienne. Vous recevez sur WhatsApp un message affirmant qu'un concours de la fonction publique est reporté, ou qu'un remède traditionnel guérit telle maladie. Deux attitudes sont possibles : croire et transférer immédiatement, ou bien s'arrêter, examiner, vérifier. Cette seconde attitude, c'est exactement ce que la philosophie entraîne.

\begin{definition}[L'esprit critique]
L'\cle{esprit critique} est l'aptitude à examiner les opinions, les informations et les raisonnements avant de les accepter : on demande des preuves, on vérifie les sources, on repère les contradictions, et l'on refuse de croire sur la seule autorité du nombre ou du plus fort. Il ne consiste pas à tout rejeter, mais à n'accepter que ce qui résiste à l'examen rationnel.
\end{definition}

L'esprit critique conduit à l'\cle{autonomie de jugement}, c'est-à-dire la capacité de penser par soi-même. Le philosophe Emmanuel Kant résumait cette exigence par la formule : \emph{« Aie le courage de te servir de ton propre entendement »} (\emph{Qu'est-ce que les Lumières ?}). Celui qui juge par lui-même ne répète pas servilement l'opinion du groupe, du chef ou de l'influenceur ; il se fait sa propre conviction, fondée sur des raisons.

\begin{exemplebox}[Vérifier une information virale sur WhatsApp]
Un message circule dans un groupe de quartier à Douala : « Boire de l'eau salée chaque matin protège contre toutes les infections. Partagez pour sauver des vies ! » L'esprit critique applique la méthode philosophique en quatre questions :
\begin{enumerate}
\item \textbf{Quelle est la source ?} Le message ne cite ni médecin, ni ministère de la Santé, ni étude.
\item \textbf{Y a-t-il des preuves ?} Aucun chiffre, aucune expérience, aucun raisonnement n'est donné.
\item \textbf{Le raisonnement est-il cohérent ?} Affirmer qu'un seul produit protège de « toutes » les infections est une généralisation abusive.
\item \textbf{Quel est le mobile ?} La consigne « partagez » cherche l'émotion et la peur, non la vérité.
\end{enumerate}
Conclusion : on ne croit pas, on ne transfère pas, et l'on vérifie auprès d'une source fiable (MINSANTE, OMS). Voilà la philosophie à l'œuvre dans la vie quotidienne.
\end{exemplebox}

Face aux rumeurs, aux réseaux sociaux et aux manipulations (désinformation, fausses images, discours démagogiques), la philosophie est donc un véritable bouclier : elle apprend à distinguer le vrai du faux, le fait de l'opinion, l'argument de la simple affirmation.

\courssub{Sur le plan moral : réfléchir aux valeurs qui guident nos actions}

Chaque jour, nous accomplissons des actes que nous jugeons bons ou mauvais : rendre un portefeuille trouvé, tricher à un examen, respecter ou humilier un camarade. Mais pourquoi le bien est-il bien, et le mal est-il mal ? La philosophie nous oblige à ne pas nous contenter d'habitudes ou de coutumes reçues sans examen : elle nous invite à \emph{réfléchir aux valeurs} — le bien, le mal, la justice, le devoir — qui doivent guider nos actions.

Socrate affirmait : \emph{« Une vie sans examen ne vaut pas la peine d'être vécue »} (\emph{Apologie de Socrate}). Examiner sa vie, c'est se demander : mes actes sont-ils justes ? Suis-je honnête par conviction ou seulement par peur d'être puni ? Grâce à cette réflexion, la morale cesse d'être une obéissance aveugle pour devenir un choix libre et responsable. La philosophie forme ainsi la \cle{conscience morale} : elle ne donne pas des ordres, elle apprend à se donner à soi-même des règles réfléchies.

\courssub{Sur le plan social et civique : préparer le citoyen au débat démocratique}

Vivre en société, c'est vivre avec des personnes qui ne pensent pas comme nous. Comment décider ensemble sans violence ? Par le \cle{débat} : échanger des arguments, écouter l'autre, accepter d'avoir tort. Or le débat démocratique exige précisément les vertus que la philosophie entraîne :

\begin{itemize}
\item \textbf{argumenter} au lieu d'insulter ou de menacer ;
\item \textbf{écouter} la position adverse avant de la réfuter ;
\item \textbf{distinguer} la personne de ses idées (on critique une thèse, on n'attaque pas celui qui la défend) ;
\item \textbf{pratiquer la tolérance}, c'est-à-dire respecter la liberté de pensée d'autrui, même quand on la désapprouve.
\end{itemize}

La philosophie prépare ainsi le citoyen : celui qui sait discuter peut voter en connaissance de cause, participer aux débats publics, résister aux discours de haine et aux manipulations politiques. Une démocratie sans citoyens capables de raisonner ensemble est une démocratie en danger.

\courssub{Sur le plan intellectuel : donner du sens aux savoirs scientifiques et techniques}

Vous êtes des élèves de l'enseignement technique : vous apprenez des savoir-faire précis (électricité, mécanique, comptabilité, informatique, construction...). Mais la science et la technique répondent à la question « comment ? », non à la question « pourquoi ? » ni « pour quoi faire ? ». Un technicien sait construire un pont ; mais faut-il construire ce pont-là, à cet endroit, pour quelle finalité humaine ? Un informaticien sait traiter des données ; mais jusqu'où peut aller la surveillance des personnes ?

La philosophie \emph{donne du sens} aux savoirs : elle interroge leurs fondements (qu'est-ce qu'une preuve ? qu'est-ce que la vérité ?) et leurs finalités (le progrès technique est-il toujours un progrès humain ?). Sans elle, le savant risque de devenir, selon l'expression célèbre, « un savant ignorant » : très compétent dans sa spécialité, mais incapable de réfléchir au sens et aux conséquences de ce qu'il fait. La philosophie réunit ce que les spécialités séparent : elle offre une vision d'ensemble.

\courssub{Sur le plan professionnel : rigueur, analyse et argumentation}

Contrairement à une idée reçue, la philosophie développe des compétences directement utiles à tout métier :

\begin{itemize}
\item la \textbf{rigueur} : définir ses termes, éviter les confusions, raisonner sans contradiction ;
\item l'\textbf{analyse} : décomposer un problème complexe en questions simples, repérer l'essentiel ;
\item l'\textbf{argumentation} : défendre un projet, rédiger un rapport, convaincre un client ou un jury avec des raisons ordonnées ;
\item la \textbf{prise de décision} : peser le pour et le contre avant d'agir, compétence précieuse pour tout responsable.
\end{itemize}

Un technicien qui sait analyser une panne méthodiquement, un comptable qui justifie ses choix, un chef de chantier qui tranche après examen des arguments : tous pratiquent, sans le savoir, les habitudes formées par la philosophie.

\courssub{La philosophie, école de la liberté : se libérer des préjugés}

\begin{definition}[Le préjugé]
Un \cle{préjugé} est un jugement adopté \emph{avant} tout examen (du latin \emph{praejudicium}) : une opinion toute faite, héritée du groupe, de la coutume ou de l'habitude, que l'on répète sans l'avoir vérifiée. Exemples : « les gens de telle région sont tous pareils », « les filles ne sont pas faites pour les métiers techniques ».
\end{definition}

Le préjugé est une prison de l'esprit : il nous fait juger sans connaître, exclure sans comprendre. La philosophie est une \emph{école de la liberté} parce qu'elle nous libère de ces chaînes invisibles : en examinant nos opinions, nous découvrons lesquelles sont fondées et lesquelles ne sont que des préjugés. Être libre, ce n'est pas faire ce qu'on veut ; c'est d'abord \emph{penser par soi-même} au lieu de laisser penser les autres à notre place. Comme le montre Descartes avec le doute, il faut parfois tout remettre en question pour ne garder que ce qui est solide.

\begin{popfigure}
\begin{center}
\begin{tikzpicture}[scale=1]
\node[circle, draw=popdark, fill=popgoldL, text=popdark, font=\bfseries, align=center, minimum size=2.6cm] (C) at (0,0) {Utilité de la\\philosophie};
\node[draw=popblue, fill=popblueL, rounded corners, align=center, font=\small, text width=3.1cm] (N1) at (90:3.4) {\textbf{Individuelle}\\esprit critique,\\autonomie de jugement};
\node[draw=poppurple, fill=poppurpleL, rounded corners, align=center, font=\small, text width=3.1cm] (N2) at (162:3.4) {\textbf{Morale}\\réflexion sur le bien,\\le mal, la justice};
\node[draw=popgreen, fill=popgreenL, rounded corners, align=center, font=\small, text width=3.1cm] (N3) at (234:3.4) {\textbf{Civique}\\débat démocratique,\\tolérance};
\node[draw=poporange, fill=poporangeL, rounded corners, align=center, font=\small, text width=3.1cm] (N4) at (306:3.4) {\textbf{Professionnelle}\\rigueur, analyse,\\argumentation};
\node[draw=poppink, fill=poppinkL, rounded corners, align=center, font=\small, text width=3.1cm] (N5) at (18:3.4) {\textbf{Intellectuelle}\\sens des savoirs\\scientifiques et techniques};
\draw[-{Stealth}, thick, popdark] (C) -- (N1);
\draw[-{Stealth}, thick, popdark] (C) -- (N2);
\draw[-{Stealth}, thick, popdark] (C) -- (N3);
\draw[-{Stealth}, thick, popdark] (C) -- (N4);
\draw[-{Stealth}, thick, popdark] (C) -- (N5);
\end{tikzpicture}
\end{center}
\legende{Diagramme en étoile des cinq dimensions de l'utilité de la philosophie. Chaque branche correspond à un plan d'importance étudié dans cette section ; toutes convergent vers un même centre : former un homme libre et responsable.}
\end{popfigure}

\courssub{Application : la philosophie et la résolution pacifique des conflits communautaires au Cameroun}

Le Cameroun compte plus de 250 groupes ethniques et deux grandes traditions linguistiques officielles. Cette diversité est une richesse, mais elle peut aussi nourrir des tensions : conflits entre agriculteurs et éleveurs dans les régions du Nord, de l'Adamaoua et du Nord-Ouest, rivalités foncières entre villages, stéréotypes ethniques amplifiés par les réseaux sociaux. Comment la philosophie aide-t-elle à résoudre ces conflits pacifiquement ?

\begin{enumerate}
\item \textbf{En déconstruisant les préjugés ethniques.} L'esprit critique montre qu'un jugement comme « tous les membres de telle communauté sont... » est une généralisation abusive : on ne peut juger des millions de personnes sur le comportement de quelques-uns.
\item \textbf{En substituant le dialogue à la violence.} La méthode philosophique apprend à exposer son grief par des arguments, à écouter la partie adverse, à chercher une solution juste plutôt qu'une victoire. C'est l'esprit du « palaver » traditionnel, que la philosophie éclaire et renforce.
\item \textbf{En fondant la justice.} Un conflit ne se règle durablement que si la solution est reconnue équitable par les deux parties ; or réfléchir à ce qui est juste est précisément le travail du philosophe.
\item \textbf{En formant des citoyens tolérants.} La tolérance n'est pas l'indifférence : c'est le respect raisonné de l'autre, fondé sur l'idée que la vérité se cherche ensemble.
\end{enumerate}

\begin{exoresolu}[Dissertation guidée : la philosophie est-elle utile ?]
\textbf{Sujet.} Certains affirment : « la philosophie ne sert à rien, elle ne fabrique rien ». Discutez cette opinion en rédigeant une introduction et un plan détaillé.
\tcblower
\textbf{Introduction (modèle).} La philosophie n'enseigne ni à construire une maison, ni à soigner une maladie, ni à réparer un moteur. Faut-il en conclure, avec l'opinion commune, qu'elle « ne sert à rien » ? Cette question suppose que seul est utile ce qui produit un objet matériel. Mais n'existe-t-il pas une utilité d'un autre ordre : celle qui forme l'esprit et rend l'homme libre ? Nous montrerons que si la philosophie est inutile au sens pratique et immédiat, elle est indispensable au sens intellectuel, moral et civique.

\textbf{Plan détaillé.}
\begin{itemize}
\item \textbf{I. Le reproche d'inutilité} : la philosophie ne produit aucun bien matériel ; ses questions semblent sans réponse (exemple : les débats abstraits). Thèse du sens commun.
\item \textbf{II. Mais la philosophie possède une utilité supérieure} : elle forme l'esprit critique face aux rumeurs et aux manipulations (exemple de l'information virale) ; elle éclaire nos valeurs morales ; elle prépare au débat démocratique ; elle donne du sens aux savoirs techniques ; elle développe rigueur et argumentation utiles à tout métier.
\item \textbf{III. La vraie utilité est la liberté} : se libérer des préjugés, penser par soi-même (Kant) ; application : résolution pacifique des conflits communautaires au Cameroun. L'inutile au sens marchand est le plus nécessaire à l'homme.
\end{itemize}
\textbf{Justification de la démarche.} Le plan suit la progression dialectique exigée par le verbe « discutez » : thèse adverse (I), réfutation et thèse personnelle (II), dépassement (III).
\end{exoresolu}

\begin{attention}[Erreur fréquente : « la philosophie ne sert à rien »]
Dans une dissertation, ne répétez jamais sans examen que « la philosophie ne sert à rien ». Cette affirmation confond deux sens du mot \emph{utile} : l'\textbf{utilité pratique} (produire un objet, gagner de l'argent immédiatement) et l'\textbf{utilité intellectuelle et morale} (former le jugement, rendre libre, donner du sens). La philosophie est inutile au premier sens, mais précisément par là, elle est libre et libératrice au second sens. Un bon devoir distingue toujours ces deux niveaux avant de conclure.
\end{attention}

\begin{savaistu}[Une Journée mondiale de la philosophie]
L'UNESCO a institué en 2002 une \textbf{Journée mondiale de la philosophie}, célébrée chaque année le troisième jeudi du mois de novembre. Son objectif est de promouvoir la réflexion critique et le dialogue rationnel dans le monde entier, en particulier auprès des jeunes. Au Cameroun, des lycées et des universités organisent ce jour-là des débats publics, des concours de dissertation et des cafés philo : preuve que la philosophie est reconnue, au niveau mondial, comme un instrument de paix et de formation du citoyen.
\end{savaistu}

\begin{aretenir}[L'essentiel de la section]
\begin{itemize}
\item Sur le plan \cle{individuel}, la philosophie forme l'esprit critique et l'autonomie de jugement, armes contre les rumeurs et les manipulations.
\item Sur le plan \cle{moral}, elle nous fait réfléchir aux valeurs (bien, mal, justice) au lieu d'obéir aveuglément.
\item Sur le plan \cle{civique}, elle prépare au débat démocratique et à la tolérance.
\item Sur le plan \cle{intellectuel}, elle donne du sens aux savoirs scientifiques et techniques.
\item Sur le plan \cle{professionnel}, elle développe rigueur, analyse et argumentation.
\item École de la liberté, elle nous délivre des préjugés ; au Cameroun, elle sert la résolution pacifique des conflits communautaires.
\end{itemize}
\end{aretenir}

Ainsi comprise, l'importance de la philosophie n'est plus un slogan : c'est une évidence démontrée. Mais une objection demeure : la science, elle aussi, forme l'esprit et combat l'erreur. Qu'apporte alors la philosophie de spécifique ? C'est toute la question de la section suivante : le rapport entre la philosophie et la science.

\coursec{Philosophie et science}

Après avoir examiné l'importance de la philosophie dans la formation de l'esprit critique et la construction d'une existence lucide, il est légitime de se demander quelle place elle occupe face aux sciences. Longtemps confondues, la philosophie et les sciences ont suivi des chemins distincts, mais elles n'ont jamais cessé de se croiser. Comprendre leurs rapports, c'est saisir comment l'esprit humain organise son savoir : d'un côté, la conquête de l'explicable par l'expérience ; de l'autre, l'interrogation sur le sens de l'explicable.

\courssub{Une origine commune : la philosophie naturelle}

Dans l'Antiquité grecque, le mot \emph{philosophia} (amour de la sagesse) désignait l'ensemble des connaissances rationnelles. Les premiers philosophes — Thalès, Héraclite, Parménide, Aristote — étaient à la fois métaphysiciens, physiciens, biologistes, astronomes et logiciens. Aristote, par exemple, a écrit sur la \emph{Métaphysique} (la philosophie première), mais aussi sur la \emph{Physique}, l'\emph{Éthique}, la \emph{Politique}, l'\emph{Histoire des animaux} et la \emph{Météorologie}. Il n'y avait alors aucune frontière étanche : la \cle{philosophie naturelle} (ou \emph{physiologie}) cherchait les principes premiers (\emph{archai}) de tout ce qui existe, du mouvement des astres à la croissance des plantes.

Cette unité du savoir a perduré jusqu'à la Renaissance. Galilée, Descartes, Leibniz, Newton se pensaient encore comme des « philosophes de la nature ». La rupture moderne s'opère lorsque des domaines particuliers acquièrent leur autonomie méthodologique et épistémologique, devenant des \cle{sciences} à part entière.

\begin{savaistu}[Newton et la « philosophie naturelle »]
Isaac Newton intitula son œuvre majeure, publiée en 1687, \emph{Philosophiæ Naturalis Principia Mathematica} (\emph{Principes mathématiques de la philosophie naturelle}). Pour lui, la mécanique céleste et terrestre relevait encore de la philosophie, mais d'une philosophie devenue \emph{mathématique} et \emph{expérimentale}. Le terme « scientist » (scientifique) n'apparaîtra qu'en 1834, sous la plume de William Whewell, pour désigner les cultivateurs des sciences spécialisées, distincts du « philosophe » spéculatif.
\end{savaistu}

\courssub{La séparation progressive : la naissance des sciences autonomes}

À partir du XVII\textsuperscript{e} siècle, la \cle{révolution scientifique} opère une dissociation progressive. La \emph{mathématique} (Galilée : « La nature est écrite en langage mathématique »), la \emph{physique} (lois du mouvement, gravitation), la \emph{chimie} (Lavoisier, fin XVIII\textsuperscript{e}), la \emph{biologie} (Darwin, XIX\textsuperscript{e}), puis la \emph{psychologie} (Wundt, fin XIX\textsuperscript{e}) et la \emph{sociologie} (Comte, Durkheim) se constituent en disciplines distinctes.

Chaque science se dote :
\begin{itemize}
    \item d'un \textbf{objet propre} délimité (l'atome, la cellule, la société, l'inconscient) ;
    \item d'une \textbf{méthode spécifique} (expérience contrôlée, modélisation mathématique, observation statistique) ;
    \item de \textbf{critères de vérité} intersubjectifs (répétabilité, falsifiabilité, consensus de la communauté scientifique).
\end{itemize}

La philosophie se retrouve alors recentrée sur les questions que les sciences ne peuvent pas (ou pas encore) résoudre : le sens de l'existence, la nature de la conscience, les fondements de la morale, la justice, la beauté, la vérité elle-même.

\begin{propriete}[La philosophie, « mère des sciences »]
Historiquement et logiquement, la philosophie est la « mère des sciences ». Elle a enfanté les sciences en posant les premiers problèmes, en forgeant les concepts de base (cause, substance, espace, temps, vie, matière) et en exigeant la rationalité. Une science naît souvent le jour où une question philosophique trouve une méthode de résolution propre (ex. : la psychologie se détache de la philosophie de l'âme quand Wundt crée le premier laboratoire de psychologie expérimentale en 1879). Cette propriété n'est pas un théorème à démontrer, mais un fait historique et épistémologique majeur.
\end{propriete}

\courssub{Points communs : la rationalité et la quête de vérité}

Malgré leur divergence, philosophie et science partagent un même idéal : la \cle{raison}. Elles refusent l'argument d'autorité, le dogme, le mythe ou l'opinion non fondée (\emph{doxa}).

\begin{popfigure}
\begin{tikzpicture}[scale=1.1, every node/.style={font=\small}]
% Couleurs
\definecolor{philoc}{RGB}{70,130,180} % popblue
\definecolor{scic}{RGB}{220,100,50}   % poporange
\definecolor{interc}{RGB}{100,180,100} % popgreen

% Cercles
\begin{scope}[blend group=soft light]
    \fill[philoc!60] (-1.5,0) circle (3cm);
    \fill[scic!60] (1.5,0) circle (3cm);
\end{scope}

% Contours
\draw[philoc, thick] (-1.5,0) circle (3cm);
\draw[scic, thick] (1.5,0) circle (3cm);

% Étiquettes principales
\node[philoc, font=\bfseries\large] at (-3.5, 2.8) {Philosophie};
\node[scic, font=\bfseries\large] at (3.5, 2.8) {Science};

% Zone Philosophie seule (gauche)
\node[align=center, text width=3.5cm] at (-3.5, 0.5) {
    \textbf{Spécifique à la philosophie}\\
    \textbullet\ Objet total / sens du réel\\
    \textbullet\ Valeurs, éthique, esthétique, politique\\
    \textbullet\ Méthode : analyse conceptuelle, argumentation, pensée critique\\
    \textbullet\ Résultats : thèses discutables, interprétations ouvertes\\
    \textbullet\ Vérité : cohérence, pertinence, universalité claimée mais non démontrable empiriquement
};

% Zone Science seule (droite)
\node[align=center, text width=3.5cm] at (3.5, 0.5) {
    \textbf{Spécifique à la science}\\
    \textbullet\ Objets partiels, observables, mesurables\\
    \textbullet\ Faits, lois, mécanismes causaux\\
    \textbullet\ Méthode : observation, hypothèse, expérience, modélisation\\
    \textbullet\ Résultats : théories validées, prédictions vérifiables\\
    \textbullet\ Vérité : correspondance aux faits, falsifiabilité, consensus provisoire
};

% Intersection (centre)
\node[align=center, text width=4cm, font=\bfseries, interc] at (0, -0.2) {
    \large \textbf{Points communs}\\
    \normalsize
    \textbullet\ Quête de la vérité objective\\
    \textbullet\ Rationalité, logique, cohérence\\
    \textbullet\ Exigence de justification / preuve\\
    \textbullet\ Esprit critique, remise en cause\\
    \textbullet\ Universalité visée du discours\\
    \textbullet\ Usage du langage conceptuel rigoureux
};
\end{tikzpicture}
\legende{Diagramme de Venn : relations entre philosophie et science. L'intersection (vert) montre l'idéal rationaliste partagé ; les zones distinctes (bleu/orange) marquent les différences d'objet, de méthode et de régime de vérité.}
\end{popfigure}

\courssub{Différences fondamentales : objet, méthode, résultats}

C'est sur le plan de la \textbf{spécificité} que se joue la distinction opératoire. Le tableau suivant, à compléter par l'élève lors du TP, synthétise les trois écarts majeurs.

\begin{popfigure}
\begin{center}
\begin{tabularx}{\linewidth}{|>{\bfseries}l|X|X|}
\hline
\rowcolor{popblueL}
\textbf{Critère de distinction} & \textbf{Philosophie} & \textbf{Science} \\
\hline
\textbf{Objet} &
Le \textbf{total} (le réel dans sa globalité), le \textbf{sens}, les \textbf{valeurs}, l'\textbf{absolu}, l'\textbf{inconditionné} (Dieu, Liberté, Immortalité, Bien, Beau, Juste). Questions « ultimes » : \emph{Pourquoi y a-t-il quelque chose plutôt que rien ? Qu'est-ce que le bien ?} &
Des \textbf{objets partiels}, \textbf{observables}, \textbf{mesurables}, \textbf{quantifiables}. Domaines délimités : la matière inerte (physique/chimie), le vivant (biologie), le social (sociologie/économie), le psychique (psychologie). Questions « factuelles » : \emph{Comment fonctionne la photosynthèse ? Quelle est la masse du boson de Higgs ?} \\
\hline
\textbf{Méthode} &
\textbf{Réflexion conceptuelle}, \textbf{analyse logique}, \textbf{argumentation rationnelle}, \textbf{pensée critique}, \textbf{herméneutique} (interprétation des textes/sens). Pas d'expérience cruciale pour trancher. Le « laboratoire » est le concept et le dialogue. &
\textbf{Méthode expérimentale} (Galilée) : Observation $\to$ Hypothèse $\to$ Expérience contrôlée / Modélisation mathématique $\to$ Validation / Falsification (Popper). Reproductibilité, quantification, outillage technique. Le laboratoire est physique. \\
\hline
\textbf{Résultat / Type de vérité} &
Des \textbf{thèses}, des \textbf{systèmes}, des \textbf{arguments}. Vérité de \textbf{cohérence} et de \textbf{pertinence} (universalité claimée mais \textbf{non démontrable} empiriquement). Résultats \textbf{problématiques, discutables, révisibles} indéfiniment. Pas de consensus définitif. &
Des \textbf{lois}, des \textbf{théories}, des \textbf{modèles}. Vérité de \textbf{correspondance} (adéquation au réel) et de \textbf{validité intersubjective}. Résultats \textbf{démontrables, vérifiables, falsifiables}. Consensus scientifique provisoire mais robuste (progrès cumulatif). \\
\hline
\end{tabularx}
\end{center}
\legende{Tableau comparatif philosophie / sciences. À compléter en TP : l'élève doit expliciter chaque case à partir de deux textes fournis (ex. : extrait du \emph{Discours de la méthode} de Descartes et extrait d'un manuel de physique sur la méthode expérimentale).}
\end{popfigure}

\begin{definition}[La science (connaissance scientifique)]
La science est un ensemble de connaissances \textbf{exactes}, \textbf{objectives}, \textbf{vérifiables} et \textbf{systématiques}, portant sur des objets déterminés, obtenues par une \textbf{méthode rigoureuse} (observation, hypothèse, expérimentation, modélisation mathématique) et soumises au critère de \textbf{falsifiabilité} (Karl Popper). Elle vise l'universalité et la nécessité (lois), et non l'opinion ou la croyance.
\end{definition}

\courssub{Complémentarité : le « comment » et le « pourquoi »}

L'erreur la plus fréquente est de croire que science et philosophie s'excluent. En réalité, elles sont \textbf{complémentaires} : la science décrit \emph{comment} les choses fonctionnent (mécanismes, causes efficientes, lois) ; la philosophie interroge \emph{pourquoi} elles existent ainsi et \emph{à quelles fins} nous devons les utiliser (finalités, valeurs, éthique).

\begin{exemplebox}[Le vaccin : efficacité scientifique et questionnement éthique]
\emph{Situation concrète au Cameroun : campagne de vaccination contre le Covid-19 ou la poliomyélite.}
\begin{itemize}
    \item \textbf{La science (biologie, immunologie, épidémiologie)} établit l'efficacité du vaccin : taux de séroconversion, diminution de la létalité, effet de groupe (immunité collective). Elle fournit des \textbf{données chiffrées} (ex. : « 95\% d'efficacité », « $R_0$ réduit sous 1 »). Elle dit \textbf{comment} le vaccin agit (ARNm, vecteur viral, protéine Spike).
    \item \textbf{La philosophie (éthique, philosophie du droit, politique)} interroge : \textbf{Faut-il rendre la vaccination obligatoire ?} Le passeport vaccinal est-il juste ? Comment arbitrer entre \cle{liberté individuelle} (disposer de son corps) et \cle{solidarité collective} (protéger les vulnérables) ? La science ne peut pas trancher ce dilemme \emph{éthique} : elle informe le choix, elle ne le dicte pas.
\end{itemize}
\textbf{Enseignement :} Sans science, la décision est aveugle (irrationnelle) ; sans philosophie, la décision est muette sur le sens et la justice (technocratique).
\end{exemplebox}

\begin{attention}[Erreur fréquente : opposer totalement philosophie et science]
\emph{Piège de rédaction :} Écrire « La science donne des certitudes, la philosophie ne donne que des opinions. » ou « La philosophie est inutile depuis que la science explique tout. »
\emph{Correction :}
\begin{enumerate}
    \item La science ne donne pas des « certitudes absolues » mais des \textbf{valides provisoires} (falsifiables). Newton a été « corrigé » par Einstein.
    \item La philosophie ne produit pas des « opinions » subjectives, mais des \textbf{arguments rationnels} sur des questions que la science \emph{ne peut pas} traiter (le sens, la valeur, le fondement).
    \item \textbf{Au Bac :} Dans une dissertation sur « Science et philosophie », il faut \textbf{articuler} : distinction (objets/méthodes différents) $\to$ non-opposition (même exigence de raison) $\to$ complémentarité (éclairer l'action humaine).
\end{enumerate}
\end{attention}

\courssub{TP guidé : construire un tableau comparatif et un diagramme de Venn}

\textbf{Consigne :} À partir des deux textes ci-après, compléter le tableau comparatif (voir figure ci-dessus) et construire un diagramme de Venn sur votre cahier.

\textbf{Texte 1 (Descartes, \emph{Discours de la méthode}, 1637) :} « Car il n'y a rien qui soit si éloigné de la philosophie que la vérité... Je me mis à chercher une méthode... La première [règle] était de ne recevoir jamais aucune chose pour vraie que je ne la connusse évidemment être telle... La seconde, de diviser chacune des difficultés que j'examinerais, en autant de parcelles qu'il se pourrait... La troisième, de conduire par ordre mes pensées... La quatrième, de faire partout des dénombrements si entiers et des revues si générales, que je fusse assuré de ne rien omettre. »

\textbf{Texte 2 (Claude Bernard, \emph{Introduction à l'étude de la médecine expérimentale}, 1865) :} « L'expérimentateur... ne raisonne pas sur ce qu'il voit, il raisonne sur ce qu'il ne voit pas... L'idée scientifique naît de l'observation... L'expérimentateur provoque les phénomènes... Il les fait varier... Il les soumet à des conditions connues... La science expérimentale est la connaissance des rapports des phénomènes entre eux et avec les conditions où ils se produisent. »

\begin{exoresolu}[Analyse des textes : philosophie (Descartes) vs science (Bernard)]
\textbf{Énoncé :} Identifier dans chaque texte la conception de la méthode et de la vérité.

\textbf{Solution détaillée :}
\begin{enumerate}
    \item \textbf{Descartes (Philosophie rationaliste) :} La méthode est \textbf{intellectuelle, a priori, déductive}. Règles : \textbf{évidence} (critère de vérité = clarté et distinction), \textbf{analyse} (division), \textbf{synthèse} (ordre), \textbf{énumération} (exhaustivité). Le sujet pensant (\emph{cogito}) est la source de la certitude. Pas d'expérience externe requise pour les vérités premières. \rightarrow \textit{Correspond à la colonne « Philosophie » du tableau : méthode conceptuelle, vérité de cohérence/évidence.}
    \item \textbf{Bernard (Science expérimentale) :} La méthode est \textbf{empirique, a posteriori, hypotético-déductive}. L'observation \textbf{provoquée} (expérience) est la source. On raisonne sur l'\textbf{invisible} (cause, mécanisme) à partir du \textbf{visible} (effet, phénomène). Variation des conditions (contrôle). Vérité = \textbf{correspondance aux faits}, \textbf{objectivité} (indépendance du sujet). \rightarrow \textit{Correspond à la colonne « Science » du tableau : méthode expérimentale, vérité de correspondance/falsifiabilité.}
    \item \textbf{Point commun (Intersection Venn) :} Les deux exigent \textbf{rigueur, ordre, rationalité, rejet du hasard et de l'autorité}. Descartes « divise », Bernard « varie les conditions » : même souci d'analyse systématique.
\end{enumerate}
\end{exoresolu}

\courssub{Synthèse et application à la vie courante}

\begin{aretenir}[Synthèse : Philosophie et Science]
\begin{itemize}
    \item \textbf{Origine commune :} Philosophie antique = science universelle (philosophie naturelle).
    \item \textbf{Séparation :} Sciences = autonomie par spécialisation de l'objet et de la méthode (XVII\textsuperscript{e}--XIX\textsuperscript{e} s.).
    \item \textbf{Commun :} Rationalité, quête de vérité, esprit critique, universalité du discours.
    \item \textbf{Différences :}
    \begin{description}
        \item[Objet :] Total/Sens (Philo) vs Partiel/Fait (Science).
        \item[Méthode :] Conceptuelle/Argumentative (Philo) vs Expérimentale/Mathématique (Science).
        \item[Vérité :] Cohérence/Interprétation ouverte (Philo) vs Correspondance/Vérification close (Science).
    \end{description}
    \item \textbf{Rapport :} \textbf{Complémentarité}. Science = \emph{pouvoir} (savoir-faire, technique) ; Philosophie = \emph{devoir} (savoir-vivre, éthique). L'intelligence artificielle, le génome, l'écologie, la bioéthique exigent les deux.
\end{itemize}
\end{aretenir}

\begin{methode}[Réussir une question sur « Philosophie et Science » (Bac/Probatoire)]
\begin{enumerate}
    \item \textbf{Définir} les deux termes (Science : connaissance exacte, vérifiable, expérimentale ; Philosophie : interrogation rationnelle sur le sens total).
    \item \textbf{Distinguer} sans opposer : utiliser le tableau (Objet / Méthode / Vérité).
    \item \textbf{Montrer l'origine commune} (philosophie naturelle) et la filiation (philosophie = mère des sciences).
    \item \textbf{Illustrer la complémentarité} par un exemple concret et actuel (OGM, IA, vaccin, environnement, fin de vie, justice algorithmique).
    \item \textbf{Conclure} sur la nécessité d'un dialogue : la science sans conscience n'est que ruine de l'âme (Rabelais) ; la philosophie sans science est vide de contenu réel.
\end{enumerate}
\end{methode}

\begin{exercice}[Application : Intelligence Artificielle et responsabilité]
Un hôpital de Yaoundé utilise un algorithme d'IA pour prioriser l'accès aux lits de réanimation. L'algorithme (science informatique, statistiques) maximise le « nombre d'années de vie sauvées ».
\begin{enumerate}
    \item Quel type de vérité l'algorithme produit-il ? (Scientifique / Philosophique ? Justifier).
    \item Le choix du critère « années de vie sauvées » est-il une décision purement scientifique ? Pourquoi ?
    \item Quelle question philosophique (éthique) se pose ici que la science ne peut trancher seule ?
    \item Proposez un argument philosophique pour ou contre l'usage de cet algorithme sans contrôle humain.
\end{enumerate}
\end{exercice}

\begin{corrige}[Exercice 1 : Intelligence Artificielle et responsabilité]
\begin{enumerate}
    \item \textbf{Vérité scientifique.} L'algorithme traite des données observables (âge, pathologies, taux de survie statistiques) selon un modèle mathématique validé. Son résultat est \textbf{vérifiable, falsifiable, reproductible} (mêmes données $\to$ même score). C'est une \emph{efficacité technique}, pas une \emph{valeur morale}.
    \item \textbf{Non.} Le choix du critère « années de vie sauvées » (utilitarisme) exclut d'autres critères possibles : « priorité aux soignants », « tirage au sort (équité) », « priorité aux plus jeunes (potentiel) », « priorité aux plus démunis (justice sociale) ». Ce choix est un \textbf{arbitrage axiologique} (choix de valeur), relevant de l'éthique philosophique et du débat démocratique, pas de la science qui ne dit \emph{comment} calculer, pas \emph{quoi} maximiser.
    \item \textbf{Question philosophique :} « La vie d'un homme de 80 ans vaut-elle moins que celle d'un enfant de 10 ans ? » / « L'efficacité collective justifie-t-elle l'injustice individuelle ? » / « Peut-on déléguer une décision de vie ou de mort à une machine non responsable ? » (Problème de la \cle{responsabilité}, de la \cle{dignité}, de la \cle{justice}).
    \item \textbf{Argument (ex. contre) :} La délégation à l'IA crée une \emph{boîte noire} irresponsable. Si l'algorithme a un biais (données d'entraînement biaisées sur les minorités), il reproduit l'injustice sous couvert d'objectivité scientifique. La \textbf{responsabilité morale} exige un \textbf{sujet humain} qui assume le choix tragique (Jonas, \emph{Principe responsabilité}). La science fournit l'outil ; la philosophie (et le droit) doit en encadrer l'usage.
\end{enumerate}
\end{corrige}

\coursec{Synthèse et application à la vie courante}

Dans la section précédente, nous avons montré que la philosophie et la science ne sont ni rivales ni identiques : la science décrit et explique les faits avec des méthodes vérifiables, tandis que la philosophie interroge le sens, les valeurs et les fondements. Mais à quoi sert tout cela dans la vie de tous les jours, à Ngaoundéré, à Douala ou dans un village du Nord-Ouest ? C'est la question de cette dernière section : après avoir dressé le bilan complet de la leçon, nous montrerons que \cle{philosopher est une activité quotidienne}, puis nous appliquerons la démarche philosophique à deux situations concrètes, avant de rédiger ensemble une conclusion argumentée, comme l'exige le savoir-faire du programme.

\courssub{Bilan de la leçon : les cinq savoirs réunis}

Reprenons méthodiquement ce que nous avons établi.

\begin{aretenir}[Bilan : les cinq réponses à la question « Qu'est-ce que la philosophie ? »]
\begin{enumerate}
\item \textbf{Définition.} La philosophie est l'amour de la sagesse (\emph{philo-sophia}) : une activité rationnelle de questionnement, de problématisation et de justification qui vise la vérité et la sagesse, c'est-à-dire une vie éclairée par la raison.
\item \textbf{Origine et objet.} Née en Grèce au VI\textsuperscript{e} siècle avant J.-C. (avec Thalès de Milet) du passage du mythe au \emph{logos}, elle a pour objet la totalité de l'expérience humaine : le monde, l'homme, la connaissance, l'action, les valeurs.
\item \textbf{Méthode.} Elle procède par le doute méthodique, la problématisation, l'analyse des concepts, l'argumentation, la discussion rationnelle et l'examen critique des préjugés.
\item \textbf{Importance.} Elle forme le jugement critique, rend libre par l'autonomie de la raison, apprend à vivre ensemble dans le respect du droit et prépare aux responsabilités de la cité.
\item \textbf{Rapport à la science.} Même exigence de rigueur et de rationalité, mais objets et méthodes différents : la science explique les faits par des lois causales vérifiables, la philosophie interroge le sens, les valeurs et les fondements. Elles se complètent : la technique donne des moyens, la philosophie interroge les fins.
\end{enumerate}
\end{aretenir}

La carte mentale ci-dessous résume visuellement l'ensemble de la leçon.

\begin{popfigure}
\begin{center}\tcbox[colback=popblue,colframe=popblue,coltext=white,arc=9pt,boxsep=4pt,left=6pt,right=6pt]{\bfseries\small Qu'est-ce que la philosophie ?}\end{center}
\vspace{-2pt}
\begin{tcbraster}[raster columns=3,raster equal height=rows,raster column skip=6pt,raster row skip=6pt,enhanced,blankest]
\begin{tcolorbox}[enhanced,colback=poporange,colframe=poporange,coltext=white,boxrule=0.9pt,arc=3pt,left=4pt,right=4pt,top=3pt,bottom=3pt,boxsep=1pt,halign=left]\bfseries\scriptsize Définitions : amour de la sagesse, questionnement rationnel\end{tcolorbox}
\begin{tcolorbox}[enhanced,colback=popgreen,colframe=popgreen,coltext=white,boxrule=0.9pt,arc=3pt,left=4pt,right=4pt,top=3pt,bottom=3pt,boxsep=1pt,halign=left]\bfseries\scriptsize Origine et objet : Grèce, mythe vers logos, totalité de l'expérience\end{tcolorbox}
\begin{tcolorbox}[enhanced,colback=poppurple,colframe=poppurple,coltext=white,boxrule=0.9pt,arc=3pt,left=4pt,right=4pt,top=3pt,bottom=3pt,boxsep=1pt,halign=left]\bfseries\scriptsize Méthode : doute, problématisation, analyse, argumentation\end{tcolorbox}
\begin{tcolorbox}[enhanced,colback=poppink,colframe=poppink,coltext=white,boxrule=0.9pt,arc=3pt,left=4pt,right=4pt,top=3pt,bottom=3pt,boxsep=1pt,halign=left]\bfseries\scriptsize Importance : liberté, jugement critique, vie commune\end{tcolorbox}
\begin{tcolorbox}[enhanced,colback=popgold,colframe=popgold,coltext=white,boxrule=0.9pt,arc=3pt,left=4pt,right=4pt,top=3pt,bottom=3pt,boxsep=1pt,halign=left]\bfseries\scriptsize Philosophie et science : rigueur commune, objets distincts\end{tcolorbox}
\end{tcbraster}

\legende{Carte mentale récapitulative : au centre, la question de la leçon ; autour, les cinq savoirs du programme, chacun associé à une couleur.}
\end{popfigure}

\courssub{La philosophie comme activité quotidienne}

Commençons par une intuition simple. Un élève de Première technique qui demande à son maître d'apprentissage : « Pourquoi faut-il respecter les normes de sécurité même quand le client ne regarde pas ? » ne demande pas seulement une information technique : il demande une \emph{justification}. Il philosophe déjà, sans le savoir. De même, le commerçant de Mokolo qui se demande s'il est juste d'augmenter ses prix pendant la période de ramadan, ou la mère de famille qui hésite entre deux écoles pour son enfant, pratiquent l'essentiel de la démarche philosophique : \cle{questionner}, \cle{problématiser} puis \cle{justifier}.

\begin{definition}[Philosopher au quotidien]
Philosopher, c'est, face à une situation concrète : (1) s'étonner et poser une question au lieu de subir l'évidence ; (2) analyser les mots et les enjeux (problématiser) ; (3) examiner plusieurs réponses possibles (thèses) ; (4) justifier son choix par des raisons que d'autres peuvent discuter (argumentation) ; (5) assumer la conclusion et ses conséquences.
\end{definition}

Autrement dit, la philosophie n'est pas réservée aux professeurs ni aux livres. Elle commence dès qu'un homme refuse de répéter « on a toujours fait ainsi » et demande « pourquoi ? ». La différence entre le philosophe de métier et l'homme de la rue n'est pas de nature, mais de degré : le premier rend cette démarche systématique, rigoureuse et explicite.

\begin{methode}[Traiter une situation concrète avec la démarche philosophique]
\begin{enumerate}
\item \textbf{Décrire la situation} avec précision : qui ? quoi ? où ? quel est le conflit ou le problème ?
\item \textbf{Dégager la question philosophique} (problématisation) : transformer le problème pratique en question générale (justice, liberté, devoir, vérité, bonheur).
\item \textbf{Identifier les thèses en présence} : quelles réponses défendent les différents acteurs ? Quels arguments et quels présupposés ?
\item \textbf{Examiner les arguments} (analyse critique) : sont-ils valides ? reposent-ils sur des préjugés, sur la seule autorité, sur la tradition, ou sur des raisons rationnelles ?
\item \textbf{Conclure de manière justifiée} : prendre position en donnant ses raisons, reconnaître les limites de sa réponse, et indiquer ce que la conclusion change dans la pratique.
\end{enumerate}
\end{methode}

\courssub{Étude de cas guidée n°1 : le débat de quartier sur la délimitation d'un terrain}

\textbf{La situation.} Dans un quartier de Bafoussam, deux familles se disputent une bande de terrain de trois mètres de large. La famille N. invoque la \cle{coutume} : « Nos grands-parents cultivaient jusqu'au manguier ; la frontière a toujours été là. » La famille M. brandit un \cle{titre foncier} obtenu à la conservation : « Le droit moderne dit que la limite est au ruisseau. » Le chef du quartier propose une réunion de conciliation. Les jeunes du quartier s'interrogent : qui a raison ?

\textbf{Étape 1 — Description.} Le conflit oppose deux sources de légitimité : la tradition orale (coutume) et le droit écrit de l'État (légalité formelle).

\textbf{Étape 2 — La question philosophique.} Derrière le problème pratique se cache une question générale : \emph{qu'est-ce qui rend une règle juste ?} La justice se réduit-elle à la loi écrite (légalisme), ou à la coutume (traditionalisme), ou à l'équité ?

\textbf{Étape 3 — Les thèses en présence.}
\begin{itemize}
\item \emph{Thèse traditionaliste} : la coutume est juste parce qu'elle a fait ses preuves, qu'elle unit la communauté et qu'elle incarne la sagesse des ancêtres.
\item \emph{Thèse légaliste} : seule la loi écrite, égale pour tous et vérifiable, garantit la justice contre l'arbitraire.
\item \emph{Thèse conciliatrice (équité)} : la justice exige parfois d'accorder le droit à l'équité, c'est-à-dire d'adapter la règle au cas particulier.
\end{itemize}

\textbf{Étape 4 — Examen critique.} L'argument « on a toujours fait ainsi » est un \emph{argument d'autorité} (ou de tradition) : il prouve l'ancienneté, non la justice (l'esclavage aussi était ancien et coutumier). Inversement, le titre foncier peut être juridiquement valide (\cle{légal}) mais moralement contestable s'il a été obtenu sans enquête de voisinage sérieuse ; il manque alors de \cle{légitimité}. On retrouve ici la distinction classique entre le \textbf{légal} (conforme à la loi positive) et le \textbf{légitime} (conforme à la justice idéale). Platon déjà, dans le \emph{Gorgias}, refusait d'identifier la justice à la loi du plus fort ou au simple usage établi.

\textbf{Étape 5 — Conclusion justifiée.} Une solution philosophiquement défendable consiste à exiger que la conciliation respecte deux critères rationnels : la \emph{vérification} (enquête auprès des anciens, bornes, documents cadastraux) et l'\emph{équité} (aucune des deux familles ne doit être ruinée par la décision). La coutume et le droit ne s'opposent pas absolument : la coutume est respectable lorsqu'elle est juste, et la loi n'est pleinement légitime que lorsqu'elle est appliquée avec équité.

\begin{attention}[Piège à éviter dans ce type de cas]
Ne confondez jamais \emph{décrire} une coutume et la \emph{justifier}. Dire « la coutume exige ceci » est un constat sociologique ou juridique ; dire « il est juste de faire ceci » demande des raisons éthiques. La philosophie commence précisément à l'endroit où la simple description ne suffit plus.
\end{attention}

\courssub{Étude de cas guidée n°2 : le choix d'une filière après la Première}

\textbf{La situation.} Amina, élève de Première technique à Garoua, réussit bien. Son père veut la voir choisir la filière conduisant au génie civil, « parce que c'est là qu'il y a du travail ». Sa grande sœur lui conseille la comptabilité, « parce qu'elle aime les chiffres ». Amina, elle, rêve de froid et climatisation : elle veut installer et réparer. Elle se sent partagée entre obéissance, bonheur personnel et utilité.

\textbf{Étape 1 — Description.} Trois pressions s'exercent : l'autorité parentale (devoir filial), l'avis familial (conseil), le désir personnel (vocation). L'enjeu est un choix d'orientation, donc un choix de vie engageant l'avenir professionnel.

\textbf{Étape 2 — La question philosophique.} \emph{Que doit guider un choix de vie : le devoir envers les parents, le bonheur personnel ou l'utilité sociale ?} Autrement dit : suis-je \cle{libre} de choisir, et qu'est-ce qu'un \emph{bon} choix (rationnel et assumé) ?

\textbf{Étape 3 — Les thèses en présence.}
\begin{itemize}
\item \emph{Thèse du devoir filial} : on doit respect et obéissance à ses parents, qui ont l'expérience du terrain et qui financent les études (argument de la gratitude, de l'expérience et de la solidarité familiale).
\item \emph{Thèse du bonheur/efficacité} : on ne réussit durablement que dans ce qu'on aime ; un métier subi produit des professionnels médiocres et malheureux (argument de l'épanouissement et de la performance).
\item \emph{Thèse de l'utilité sociale} : le pays a besoin de techniciens dans certains secteurs porteurs ; le choix individuel doit tenir compte des besoins de la communauté (argument du bien commun et de l'employabilité).
\end{itemize}

\textbf{Étape 4 — Examen critique.} L'argument parental est fort mais pas absolu : les parents connaissent le marché d'hier, pas toujours celui de demain, et ils ne vivront pas le métier à la place de leur enfant. L'argument du bonheur est réel mais fragile : un désir d'adolescent peut être une mode passagère ; il faut distinguer le goût véritable (éprouvé par l'expérience, informé par la réalité du métier) du caprice. L'argument d'utilité est légitime mais ne doit pas écraser la personne : une société n'est pas servie par des techniciens résignés. Kant rappelle (\emph{Fondements de la métaphysique des mœurs}) que la personne humaine doit toujours être traitée aussi comme une fin, jamais seulement comme un moyen : sacrifier totalement Amina au « besoin du pays » ou à l'ambition paternelle serait la traiter comme un instrument.

\textbf{Étape 5 — Conclusion justifiée.} Le bon choix est celui qu'Amina peut \emph{assumer avec des raisons} : vérifier son goût par des stages ou essais en atelier, s'informer sur les débouchés réels du froid et climatisation au Cameroun (secteur porteur : conservation alimentaire, transport frigorifique, santé, confort), dialoguer avec son père en présentant un projet argumenté (plan de formation, perspectives d'embauche, création d'entreprise) plutôt qu'un simple refus. La liberté n'est pas de faire ce qu'on veut sans raison : c'est de choisir en connaissance de cause. On retrouve ici l'importance de la philosophie : elle transforme un conflit d'autorités en débat de raisons.

\courssub{Savoir-faire : rédiger une conclusion argumentée en cinq étapes}

Le programme exige que l'élève sache « rédiger et justifier une démarche, une réponse ou une conclusion ». Voici la méthode, appliquée au cas d'Amina.

\begin{methode}[Rédiger une conclusion argumentée en cinq étapes]
\begin{enumerate}
\item \textbf{Rappeler le problème} en une phrase (« Le choix d'Amina oppose le devoir familial, la recherche du bonheur personnel et l'utilité sociale »).
\item \textbf{Résumer les thèses examinées} et dire pourquoi aucune ne suffit seule (limites de chacune).
\item \textbf{Annoncer clairement sa position} (« Nous soutenons que... »).
\item \textbf{Donner la ou les raisons décisives} qui fondent cette position (l'argument clé).
\item \textbf{Ouvrir sur la portée pratique} : ce que la conclusion change dans la conduite concrète, ou la question nouvelle qu'elle soulève.
\end{enumerate}
\end{methode}

\begin{exoresolu}[Rédaction d'une conclusion argumentée (cas d'Amina)]
\textbf{Énoncé.} En cinq à huit lignes, rédigez la conclusion de la réflexion sur le choix de filière d'Amina, en suivant les cinq étapes de la méthode.
\tcblower
\textbf{Solution rédigée.} Le choix d'Amina opposait trois exigences : l'obéissance aux parents, la recherche du bonheur et l'utilité pour la société. Aucune ne peut décider seule : l'autorité parentale peut se tromper d'époque, le désir peut être un caprice, et l'utilité ne doit pas écraser la personne. Nous soutenons donc qu'Amina doit choisir librement, mais d'une liberté éclairée : vérifier son goût par l'expérience, s'informer sur les débouchés réels, et convaincre son père par un projet argumenté. La raison décisive est que seul un choix assumé pour des raisons engage durablement celui qui le fait. Ainsi la philosophie ne choisit pas à la place d'Amina : elle lui apprend à bien choisir. Reste ouverte une question plus large : jusqu'où la famille peut-elle orienter sans décider à la place de l'individu ?
\end{exoresolu}

\begin{exemplebox}[Mini-dissertation rédigée : « Faut-il toujours obéir à la coutume ? »]
\textbf{Introduction.} Dans nos villages comme dans nos villes, la coutume règle les mariages, les successions, les conflits fonciers. Elle apparaît comme la sagesse héritée des ancêtres, garante de l'ordre social. Mais l'ancienneté d'une règle suffit-elle à la rendre juste ? Faut-il toujours obéir à la coutume ? Nous montrerons d'abord que la coutume mérite respect pour sa fonction sociale, ensuite qu'elle doit être examinée par la raison critique, enfin que l'obéissance véritable est une obéissance éclairée.

\textbf{I. La coutume mérite respect.} Elle a fait ses preuves historiques : elle maintient la paix sociale, transmet des valeurs cardinales (solidarité, respect des aînés, hospitalité) et donne à chacun une identité culturelle. Refuser en bloc la coutume, c'est couper l'homme de ses racines ; le conciliabule traditionnel, par exemple, règle souvent les conflits plus vite, moins cher et plus humainement que la procédure judiciaire étatique.

\textbf{II. Mais la coutume doit être examinée par la raison.} L'argument « nos pères faisaient ainsi » est un argument d'autorité (ou de tradition) : il prouve l'ancienneté, non la justice. Certaines pratiques coutumières (lévirat imposé, exclusion systématique des filles de l'héritage foncier, mariages forcés précoces) blessent la dignité humaine et les droits fondamentaux. La raison doit donc juger la tradition elle-même : une coutume n'est respectable que si elle est juste, c'est-à-dire si elle respecte la personne humaine.

\textbf{III. Obéir, oui, mais en connaissance de cause.} La vraie fidélité aux ancêtres n'est pas la répétition aveugle : c'est de poursuivre leur recherche de la « vie bonne » avec les moyens de notre temps et les acquis du droit. L'obéissance éclairée conserve ce qui est juste dans la coutume et réforme ce qui ne l'est plus.

\textbf{Conclusion.} Il ne faut donc ni obéir toujours à la coutume (traditionalisme clos), ni la rejeter toujours (modernisme destructeur) : il faut l'interroger. La philosophie n'est pas l'ennemie de la tradition ; elle en est le discernement critique et le ferment de progrès.
\end{exemplebox}

\begin{attention}[Erreurs fréquentes en dissertation (Probatoire / Bac Tech)]
Deux fautes reviennent sans cesse dans les copies. Premièrement, \textbf{l'annonce du plan est une étape obligatoire de l'introduction} au Cameroun. Elle doit être fluide et logique (ex. : « Nous analyserons d'abord... nous verrons ensuite... nous montrerons enfin... »), non une liste scolaire. Ne pas l'annoncer pénalise la note de structure. Deuxièmement, \textbf{la conclusion ne doit pas répéter l'introduction} : elle répond à la question posée, synthétise ce qui a été établi par la démonstration, et peut ouvrir sur une question nouvelle ou une portée pratique. Une conclusion qui recopie l'introduction montre que la réflexion n'a pas avancé.
\end{attention}

\courssub{Complément utile à l'examen : vers la philosophie africaine}

\begin{savaistu}[Ce que l'épreuve valorise]
À l'examen du Probatoire technique (comme au Baccalauréat), l'épreuve de philosophie ne récompense pas la récitation de citations : elle valorise la \cle{clarté de la démarche} (problématisation, articulation logique, argumentation) autant que les connaissances. Une copie qui définit rigoureusement les termes du sujet, construit un véritable problème, oppose des thèses avec des exemples précis (camerounais si possible) et conclut en justifiant sa position obtient une excellente note, même avec peu de références d'auteurs. Apprenez donc à raisonner, pas seulement à citer des noms.
\end{savaistu}

\textbf{Articulation avec la leçon suivante (signalé pour l'examen).} Nos deux études de cas ont déjà mobilisé la coutume, la parole des anciens, la sagesse communautaire, le palaver : autant de questions qui annoncent la leçon suivante, consacrée à la \cle{philosophie africaine}. Existe-t-il une philosophie proprement africaine ? La sagesse orale, les proverbes, les palabres sont-ils de la philosophie ? Vous disposez désormais des outils pour y répondre : une définition rigoureuse de la philosophie (questionnement rationnel, problématisé et justifié), une méthode (analyse conceptuelle, argumentation, critique des préjugés) et un critère (la justification par des raisons universellement discutables, non par la seule autorité). Si la philosophie est l'amour de la sagesse, alors toute tradition qui interroge le monde, l'homme et la société et qui justifie ses réponses par le dialogue rationnel participe de l'aventure philosophique — c'est précisément ce que nous examinerons.

\begin{exercice}[Pour s'entraîner]
\textbf{Exercice 1.} Dans votre quartier ou votre village, identifiez un débat récent (partage d'un héritage, choix d'un chef, sanction d'un manquement à la règle). Traitez-le avec la démarche en cinq étapes (description, problématisation, thèses, examen critique, conclusion), puis rédigez une conclusion argumentée de cinq lignes en appliquant la méthode.

\textbf{Exercice 2.} Sujet de dissertation : « La philosophie est-elle utile au technicien ? » Rédigez uniquement l'introduction (accroche, définitions des termes, problématique, \textbf{annonce du plan}) et la conclusion (réponse justifiée, synthèse, ouverture), sans faire le développement.
\end{exercice}

\begin{corrige}[Exercice n°2 — éléments de correction attendus]
\textbf{Introduction attendue.} \emph{Accroche :} Le technicien est formé pour agir efficacement sur la matière : réparer, construire, calculer, programmer. \emph{Définitions :} La philosophie comme activité rationnelle de questionnement sur le sens, les valeurs et les fondements ; l'utilité comme capacité à produire un effet recherché, à servir une fin. \emph{Problématique :} L'utilité se mesure-t-elle seulement à l'efficacité technique immédiate ? La philosophie, qui ne fabrique aucun objet matériel, est-elle pour autant inutile au technicien ? \emph{Annonce du plan :} Nous montrerons d'abord que la philosophie est utile au technicien pour fonder son éthique professionnelle, ensuite qu'elle éclaire sa responsabilité citoyenne, enfin qu'elle est indispensable pour donner du sens à son métier.

\textbf{Conclusion attendue.} \emph{Réponse :} La philosophie est utile au technicien à un titre supérieur : elle ne lui apprend pas \emph{comment} faire (la technique), mais \emph{pourquoi} et \emph{pour qui} faire (sécurité, honnêteté, responsabilité, respect de l'environnement, sens du service). \emph{Justification :} Un technicien sans jugement critique devient un exécutant dangereux ; la technique donne des moyens (puissance), la philosophie interroge les fins (sagesse). \emph{Ouverture possible :} Toute technique puissante n'appelle-t-elle pas une éthique proportionnée à sa puissance ? Le technicien de demain ne sera-t-il pas d'abord un « philosophe praticien » ?
\end{corrige}

\newpage

\coursec{Activité d'intégration}

\coursec{Leçon 15 : Qu'est-ce que la philosophie ?}

\courssub{Objectifs de l'activité}
\begin{itemize}
    \item Mobiliser l'ensemble des notions de la leçon (définition, origine, objet, méthode, importance, rapport philosophie/science) pour analyser une situation concrète camerounaise.
    \item Identifier un problème philosophique à partir d'un fait social brut : la justice foncière.
    \item Distinguer clairement trois types de rationalité : la réponse \cle{scientifique} (bornage, droit positif), la réponse \cle{coutumière} (tradition, autorité), la réponse \cle{philosophique} (recherche du juste, universalisation, argumentation rationnelle).
    \item Construire un diagramme de Venn comparant philosophie et science pour visualiser leurs zones de convergence et de divergence.
    \item Rédiger une conclusion argumentée structurée en cinq étapes (questionnement, doute, analyse, argumentation, synthèse), conforme aux exigences de la dissertation au Probatoire.
\end{itemize}

\courssub{Situation : Un conflit de terrain au quartier « Ngodi » (Douala) — Philosopher ou trancher ?}
\begin{center}
\begin{tabularx}{\linewidth}{|X|}\hline
\rowcolor{popblueL}
\textbf{Dossier de travail — Conflit foncier au quartier Ngodi (Douala IV)} \\ \hline
\textbf{Document 1 : Récit du conflit (extrait du registre du chef de quartier, mai 2024)} \\
\emph{« Les familles \textbf{Mballa} et \textbf{Essomba} se disputent une parcelle de 450 m² située à l'entrée du quartier, bordant la route principale. La famille Mballa y a construit un hangar de stockage (tôles, bois) il y a 12 ans, sans titre foncier, sur un terrain considéré comme « vacant » par le chef de quartier d'alors. La famille Essomba produit un \textbf{certificat d'occupation} (non titré) datant de 1998, signé par l'ancien chef de village, attestant que ce terrain leur a été alloué pour culture vivrière. Depuis le bitumage de la route (2022), la valeur marchande a explosé (estimation : 15 millions FCFA).} \\
\textbf{Deux procédures parallèles s'engagent :} \\
1. \textbf{Devant le Chef de quartier actuel (M. Nkodo)} : Audition des anciens, serment sur les ancêtres, consultation du « Ngondo » local. Décision : « La terre appartient aux premiers occupants coutumiers. La famille Essomba a le droit coutumier. La famille Mballa doit démolir le hangar sous 30 jours ou payer une indemnité de 3 millions FCFA à la famille Essomba pour occuper le sol. » \\
2. \textbf{Devant un Géomètre-expert assermenté (Mme Fouda)} : Bornage GPS (précision ± 2 cm), vérification au cadastre (plan cadastral section Ngodi, feuille 12). Constat : Le certificat d'occupation des Essomba ne correspond pas aux coordonnées GPS du terrain litigieux (décalage de 15 m vers le nord). Le hangar des Mballa empiète sur la voirie publique (emprise routière de 5 m). Rapport technique : « Aucune des deux parties n'a de titre foncier parfait. Le terrain litigieux relève du domaine privé de l'État (non immatriculé). L'empiètement sur la voirie est avéré. » \\ \hline
\end{tabularx}
\end{center}

\vspace{0.5cm}
\begin{center}
\begin{tabularx}{\linewidth}{|X|}\hline
\rowcolor{poporangeL}
\textbf{Document 2 : Extrait de Kant, \emph{Logique} (1800), Introduction — Les quatre questions de la philosophie} \\
\emph{« Le champ de la philosophie peut se ramener aux questions suivantes : 1. \textbf{Que puis-je savoir ?} 2. \textbf{Que dois-je faire ?} 3. \textbf{Que m'est-il permis d'espérer ?} 4. \textbf{Qu'est-ce que l'homme ?} La première question est répondue par la métaphysique, la seconde par la morale, la troisième par la religion, la quatrième par l'anthropologie. Mais on pourrait tout réunir sous l'anthropologie, car les trois premières questions se rapportent à la dernière. [...] La philosophie est la science des rapports de toute connaissance avec les fins nécessaires de la raison. »} \\
\textbf{Repères :} Kant distingue le \cle{savoir} (connaissance théorique, science), le \cle{devoir} (obligation morale, philosophie pratique), l'\cle{espérance} (foi) et la \cle{nature humaine} (anthropologie). La philosophie ne donne pas de recettes techniques, elle interroge la \cle{valeur} et la \cle{légitimité} de nos actions. \\ \hline
\end{tabularx}
\end{center}

\vspace{0.5cm}
\begin{center}
\begin{tabularx}{\linewidth}{|X|}\hline
\rowcolor{popgreenL}
\textbf{Document 3 : Modes de règlement des conflits fonciers au Cameroun — Données MINDAF / INS (Annuaire statistique 2023, synthèse)} \\
\begin{tabular}{|l|c|c|c|}\hline
\textbf{Mode de règlement} & \textbf{Nombre de dossiers traités (2023)} & \textbf{\% du total} & \textbf{Taux d'exécution des décisions} \\ \hline
Tribunaux de grande instance (Droit positif) & 1 842 & 28 \% & 65 \% \\ \hline
Chefferies traditionnelles / Chefs de quartier (Coutume) & 3 510 & 53 \% & 78 \% \\ \hline
Médiation / Conciliation (Tierce partie neutre) & 985 & 15 \% & 85 \% \\ \hline
Autres (Arbitrage, Notaires) & 263 & 4 \% & 90 \% \\ \hline
\textbf{Total} & \textbf{6 600} & \textbf{100 \%} & \textbf{—} \\ \hline
\end{tabular} \\
\textbf{Observation clé :} La coutume reste le mode dominant (53 \%), mais son taux d'exécution (78 \%) masque des tensions : non-reconnaissance par l'État (absence de titre), risque d'arbitraire. La science (bornage, cadastre) apporte la \cle{certitude technique} (où est la limite ?) mais pas la \cle{légitimité normative} (à qui revient le droit ?). La médiation tente la synthèse. \\ \hline
\end{tabularx}
\end{center}

\courssub{Consignes}
\begin{enumerate}
    \item \textbf{Analyse documentaire (Groupe de 4) :} À partir du Document 1, identifiez : (a) le \cle{fait brut} (conflit d'intérêts) ; (b) la \cle{réponse coutumière} (autorité, tradition, sanction symbolique) ; (c) la \cle{réponse scientifique} (mesure, coordonnées, droit positif) ; (d) le \cle{problème philosophique} latent (formulez-le sous forme de question : « Qu'est-ce qui fait qu'une décision est \emph{juste} ? »).
    \item \textbf{Exploitation du Document 2 (Kant) :} Rattachez chaque décision (Chef de quartier, Géomètre) à l'une des quatre questions kantiennes. Laquelle des quatre questions reste \emph{non traitée} par les deux procédures ? Justifiez.
    \item \textbf{Construction du diagramme de Venn (Groupe) :} Sur une affiche, tracez deux cercles sécants : « \cle{Science} (Géomètre / Droit positif) » et « \cle{Philosophie} (Réflexion sur le juste) ». Placez dans l'intersection : « Recherche de vérité », « Argumentation rationnelle », « Universalisme ». Dans la partie « Science seule » : « Mesure empirique », « Procédure codifiée », « Certitude technique ». Dans « Philosophie seule » : « Question du sens », « Critique des valeurs », « Fondement de la norme », « Universalisation morale (impératif catégorique) ».
    \item \textbf{Rédaction individuelle (Évaluation) :} Rédigez une conclusion argumentée en \textbf{cinq étapes numérotées} (maximum 300 mots) répondant à la question : \emph{« Au-delà du bornage et de la coutume, en quoi la démarche philosophique est-elle indispensable pour dire le droit dans ce conflit ? »}
    \begin{enumerate}
        \item \textbf{Questionnement} : Reformulez le problème philosophique.
        \item \textbf{Doute} : Montrez l'insuffisance des réponses scientifique et coutumière prises séparément.
        \item \textbf{Analyse} : Mobilisez les notions de la leçon (objet, méthode, importance) pour définir l'apport spécifique de la philosophie.
        \item \textbf{Argumentation} : Appuyez-vous sur Kant (Doc 2) et les données statistiques (Doc 3) pour défendre une thèse.
        \item \textbf{Synthèse} : Concluez sur le rôle irremplaçable de la philosophie dans la cité.
    \end{enumerate}
    \item \textbf{Restitution orale :} Un rapporteur par groupe présente le diagramme de Venn ; deux élèves lisent leur conclusion. Débat contradictoire dirigé par le professeur au tableau.
\end{enumerate}

\courssub{Ressources à mobiliser}
\begin{definition}[Notions clés de la leçon (Unité 1)]
\textbf{Définition :} La philosophie est l'\cle{amour de la sagesse} (Philon + Sophia) : recherche rationnelle, critique et systématique des principes premiers et du sens de l'existence.
\textbf{Origine :} \cle{Étonnement} (Aristote, \emph{Métaphysique} : « C'est par l'étonnement que les hommes commencent à philosopher ») ; passage du \cle{mythe} au \cle{logos} (explication par la raison, non par le récit sacré).
\textbf{Objet :} L'\cle{universel} (l'être en tant qu'être, le bien, le vrai, le beau, l'homme) — pas un objet particulier comme la science.
\textbf{Méthode :} \cle{Rationnelle} (concepts, arguments, logique), \cle{critique} (examen des préjugés, Kant : « Sapere aude »), \cle{systématique} (cohérence), \cle{problématisation} (transformation d'un fait en question).
\textbf{Importance :} Former l'\cle{esprit critique}, fonder l'\cle{éthique} et le \cle{droit}, donner du \cle{sens} à l'action technique et scientifique, éduquer à la \cle{citoyenneté}.
\textbf{Rapport Philosophie / Science :} La science étudie le \cle{fait} (comment ? lois, mesure, prévision) ; la philosophie interroge la \cle{valeur} et le \cle{sens} (pourquoi ? fondement, éthique, épistémologie). La science \cle{presuppose} la philosophie (validité de la raison, éthique de la recherche) ; la philosophie \cle{s'appuie} sur la science (données empiriques). Ni confusion, ni séparation.
\end{definition}

\courssub{Démarche et corrigé commenté}
\begin{exoresolu}[Corrigé type de l'activité d'intégration — Conflit Ngodi]
\textbf{Consignes 1 \& 2 : Analyse documentaire et rattachement kantien (Travail de groupe)}

\textbf{a) Identification des registres (Doc 1)}
\begin{itemize}
    \item \textbf{Fait brut :} Conflit d'appropriation d'un espace physique (450 m²) à forte plus-value (15 M FCFA) entre deux groupes sociaux (familles Mballa / Essomba).
    \item \textbf{Réponse coutumière (Chef Nkodo) :} Fondée sur l'\cle{autorité traditionnelle} et la \cle{première occupation} symbolique. Procédure : oralité, serment, ancêtres. Sanction : démolition ou indemnité transactionnelle (3 M FCFA). \emph{Registre :} Légitimité historique / communautaire / sacrée.
    \item \textbf{Réponse scientifique (Géomètre Fouda) :} Fondée sur la \cle{mesure objective} (GPS, cadastre) et le \cle{droit positif} (domaine privé de l'État, emprise routière). Résultat : « Nul n'a de titre parfait », « Empiètement voirie avéré ». \emph{Registre :} Vérité technique / légalité formelle / description factuelle.
    \item \textbf{Problème philosophique :} \cle{« Qu'est-ce qu'une décision juste dans ce conflit ? »} ou \cle{« Le droit coutumier ou le droit positif suffisent-ils à fonder la justice foncière ? »}. Ce n'est pas « où est la limite ? » (science) ni « qui a l'autorité ? » (coutume), mais « quel principe universel de justice s'applique ici ? ».
\end{itemize}

\textbf{b) Rattachement aux quatre questions de Kant (Doc 2)}
\begin{center}
\begin{tabularx}{\linewidth}{|l|X|X|}\hline
\rowcolor{poppurpleL}
\textbf{Question kantienne} & \textbf{Réponse du Chef de quartier (Coutume)} & \textbf{Réponse du Géomètre (Science)} \\ \hline
\textbf{1. Que puis-je savoir ?} (Métaphysique / Épistémologie) & « Je sais par la tradition orale que les Essomba sont les ayants droit. » (Savoir basé sur l'autorité/foi) & « Je sais les coordonnées exactes (X, Y), l'emprise routière, l'absence de titre. » (\cle{Savoir scientifique}, certitude technique) \\ \hline
\textbf{2. Que dois-je faire ?} (\cle{Morale} / Philosophie pratique) & « Tu dois payer 3 M FCFA ou démolir. » (Injonction normative basée sur la coutume) & « Le rapport constate l'illégalité. » (\textbf{Ne dit pas ce qu'il faut faire}, seulement ce qui \emph{est}). \textbf{Lacune majeure.} \\ \hline
\textbf{3. Que m'est-il permis d'espérer ?} (Religion / Espérance) & Espérance de paix sociale par le respect des ancêtres. & Aucune réponse (hors champ scientifique). \\ \hline
\textbf{4. Qu'est-ce que l'homme ?} (Anthropologie) & L'homme comme \cle{sujet de droit coutumier}, membre de lignage, tenu par les ancêtres. & L'homme comme \cle{usager de l'espace}, titulaire potentiel de droit positif, géomètre. \\ \hline
\end{tabularx}
\end{center}
\emph{Commentaire :} La science (Géomètre) répond parfaitement à « Que puis-je savoir ? » (le fait brut). La coutume (Chef) tente de répondre à « Que dois-je faire ? » mais par autorité, non par raison universelle. \textbf{La philosophie intervient précisément sur la question 2 (Morale) : fonder le "doit" par la raison, non par la force ou la technique.}

\textbf{Consigne 3 : Diagramme de Venn — Philosophie vs Science (Production collective attendue)}

\begin{center}
\begin{tikzpicture}[scale=0.9, every node/.style={font=\small}]
% Cercles
\begin{scope}[blend group=soft light]
    \fill[popblue!60] (-2.0,0.0) circle (3.0cm);
    \fill[poporange!60] (2.0,0.0) circle (3.0cm);
\end{scope}
% Légendes cercles
\node[popblue, font=\bfseries] at (-3.5, 2.5) {SCIENCE};
\node[poporange, font=\bfseries] at (3.5, 2.5) {PHILOSOPHIE};
% Intersection
\node[align=center, text width=4.0cm] at (0.0,0.0) {
    \textbf{INTERSECTION}\\
    \textit{Recherche de vérité}\\
    \textit{Argumentation rationnelle}\\
    \textit{Exigence d'universalité}\\
    \textit{Usage du concept / logos}
};
% Science seule
\node[align=center, text width=3.5cm, popblue] at (-3.5, -1.0) {
    \textbf{SCIENCE SEULE}\\
    \textit{Observation / Mesure empirique}\\
    \textit{Lois causales / Prévision}\\
    \textit{Procédure codifiée (bornage)}\\
    \textit{Certitude technique (GPS ± 2 cm)}\\
    \textit{Description du fait (empiètement)}
};
% Philosophie seule
\node[align=center, text width=3.5cm, poporange] at (3.5, -1.0) {
    \textbf{PHILOSOPHIE SEULE}\\
    \textit{Question du sens / de la valeur}\\
    \textit{Critique des normes (pourquoi cette loi ?)}\\
    \textit{Fondement de la justice (le Juste)}\\
    \textit{Universalisation morale (Kant)}\\
    \textit{Problématisation (fait $\to$ problème)}
};
% Flèches explicatives
\draw[->, popmuted] (-3.5, -2.5) -- (-3.5, -3.0) node[below, align=center] {\textit{Le "Comment ?"}\\(Fait, Technique)};
\draw[->, popmuted] (3.5, -2.5) -- (3.5, -3.0) node[below, align=center] {\textit{Le "Pourquoi ?"}\\(Valeur, Sens, Devoir)};
\end{tikzpicture}
\legende{Diagramme de Venn : Distinction et articulation Science / Philosophie dans le conflit foncier.}
\end{center}

\textbf{Consigne 4 : Conclusion argumentée type (Rédaction individuelle — 5 étapes)}

\begin{enumerate}
    \item \textbf{Questionnement} : Le conflit Ngodi ne pose pas seulement un problème technique (bornage) ni seulement un problème d'autorité (coutume), mais un \textbf{problème philosophique} : \emph{« Quel principe universel de justice permet de trancher entre un droit coutumier ancestral sans titre et une occupation effective sans base légale, dans un contexte de valorisation foncière spéculative ? »}
    \item \textbf{Doute} : La réponse \textbf{scientifique} (Géomètre) apporte la \cle{vérité des faits} (coordonnées, emprise routière, absence de titre) mais reste \emph{muette sur la norme} : elle dit « ce qui est », non « ce qui doit être » (Kant, Q2). Elle aboutit à une impasse juridique (« nul n'a de titre parfait »). La réponse \textbf{coutumière} (Chef) tranche le « doit faire » mais par \cle{l'autorité} et la \cle{tradition}, sans \cle{justification rationnelle universalisable} : pourquoi 3 M FCFA ? Pourquoi l'ancienneté prime-t-elle sur l'investissement (hangar) ? Elle risque l'arbitraire (Doc 3 : 53 \% des dossiers, mais taux d'exécution 78 \% seulement).
    \item \textbf{Analyse} : C'est ici que la \textbf{démarche philosophique} (Leçon 15) devient \cle{indispensable}. Son \textbf{objet} n'est pas le terrain (fait), mais la \cle{Justice} (valeur universelle). Sa \textbf{méthode} (rationnelle, critique, systématique) transforme le fait brut en \cle{problème} : elle soumet la coutume et la loi positive à l'épreuve de la \cle{raison} (universalisation, non-contradiction, respect de la dignité). Son \textbf{importance} est de fonder la \cle{légitimité} de la décision, condition de son acceptation durable (paix sociale). Elle ne remplace ni le géomètre (faits) ni le juge (application), elle \cle{éclaire} le juge sur le \cle{fondement éthique} de sa sentence.
    \item \textbf{Argumentation} : Kant (Doc 2) montre que la question « Que dois-je faire ? » relève de la \cle{morale} (philosophie pratique), non de la science. La science (bornage) est \cle{nécessaire} (sans faits, pas de justice éclairée) mais \cle{insuffisante} (le fait ne vaut pas norme — \emph{factum non facit jus}). Les données (Doc 3) prouvent que la coutume domine (53 \%) mais que la médiation (15 \%, taux exécution 85 \%) — espace de \cle{parole raisonnée} — est la plus efficace. La philosophie \cle{institutionalise} cette médiation : elle exige que la décision soit \textbf{argumentée} (non imposée), \textbf{universalisable} (valable pour tout conflit similaire), \textbf{respectueuse des personnes} (fin en soi, non moyen). \textbf{Thèse :} Seule la philosophie peut articuler le fait scientifique (bornage) et la norme coutumière (histoire) dans un \cle{principe de justice} (ex. : sécurité juridique + protection du faible + réparation équitable) acceptable par tous.
    \item \textbf{Synthèse} : Au quartier Ngodi, philosopher n'est pas un luxe intellectuel, c'est une \cle{nécessité civique}. La philosophie ne « tranche » pas comme le chef ni ne « mesure » comme le géomètre : elle \cle{problématise}, elle \cle{critique}, elle \cle{fonde}. Elle permet de passer de la \cle{force} (coutume) ou de la \cle{technique} (science) au \cle{droit} (justice rationnelle). C'est en cela que la philosophie, par sa méthode critique et son objet universel, est la condition de possibilité d'une \cle{paix juste} et non d'une simple \cle{trêve} imposée.
\end{enumerate}

\textbf{Consigne 5 : Restitution et correction collective (Pistes pour le professeur)}
\begin{itemize}
    \item \textbf{Diagramme de Venn :} Vérifier que l'intersection n'est pas vide (rationnalité commune) et que « Prescription normative » est bien côté Philosophie (Science = descriptif).
    \item \textbf{Conclusions élèves :} Repérer les 5 étapes. Sanctionner l'absence de référence explicite à Kant ou aux données (Doc 2/3). Valoriser l'usage des concepts : \cle{problématisation}, \cle{fait/valeur}, \cle{universalisation}, \cle{légitimité}.
    \item \textbf{Débat :} « La philosophie peut-elle trancher ? » $\to$ Non, elle \cle{éclaire} le tranchant (juge, médiateur). Elle donne les \cle{critères} du juste.
\end{itemize}
\end{exoresolu}

\courssub{Bilan de l'activité}
\begin{aretenir}[Ce que l'activité a mis en évidence]
\begin{itemize}
    \item \textbf{La philosophie n'est pas une « opinion de plus » :} Elle se distingue de la coutume (autorité/tradition) et de la science (fait/mesure) par sa \cle{méthode} (rationnelle, critique, universalisante) et son \cle{objet} (le sens, la valeur, le fondement).
    \item \textbf{Le passage du fait au droit (problématisation) :} Le bornage (science) donne le fait ; la coutume donne une norme historique ; la philosophie \cle{critique} les deux au nom de la \cle{Justice} (universel) pour fonder une décision \cle{légitime}.
    \item \textbf{L'articulation Science / Philosophie (Diagramme de Venn) :} Elles partagent la raison (intersection) mais divergent sur la finalité : la science vise la \cle{vérité théorique} (prévoir, maîtriser), la philosophie vise la \cle{sagesse pratique} (orienter, juger, justifier). La science est \cle{nécessaire} à la philosophie (données réelles) ; la philosophie est \cle{nécessaire} à la science (éthique, épistémologie, sens).
    \item \textbf{Compétence évaluée (Savoir-faire) :} \emph{« Appliquer les notions à une situation concrète »} et \emph{« Rédiger et justifier une conclusion »}. La structure en 5 étapes (Questionnement $\to$ Synthèse) est le squelette de la \cle{dissertation philosophique} attendue au Probatoire.
    \item \textbf{Ancrage camerounais :} Le conflit foncier (Douala, Yaoundé, zones rurales) est le \cle{lieu privilégié} où s'affrontent droit coutumier, droit moderne, spéculation et justice. La philosophie n'y est pas abstraite : elle est l'outil pour \cle{penser le vivre-ensemble} dans le pluralisme juridique camerounais.
\end{itemize}
\end{aretenir}

\begin{methode}[Pour réussir l'activité d'intégration (Mémo élève)]
\begin{enumerate}
    \item \textbf{Lis le dossier 2 fois :} 1. Comprendre l'histoire. 2. Repérer les \emph{registres} (Scientifique / Coutumier / Philosophique / Juridique).
    \item \textbf{Identifie la question philosophique :} Elle commence souvent par « Qu'est-ce que... ? », « Faut-il... ? », « Est-il juste que... ? ». Elle ne se résout pas par une mesure ni par un ordre.
    \item \textbf{Remplis le Venn rigoureusement :} Science = Fait / Mesure / Loi naturelle. Philosophie = Valeur / Sens / Norme morale / Critique. Intersection = Raison / Argument / Universalité.
    \item \textbf{Rédige les 5 étapes sans les mélanger :} Étape 1 = Question (pas de réponse). Étape 2 = Doute (critique des réponses faciles). Étape 3 = Analyse (apport cours). Étape 4 = Argumentation (thèse + auteurs + données). Étape 5 = Synthèse (réponse finale + portée).
    \item \textbf{Relis-toi :} Ai-je cité Kant ? Ai-je utilisé les chiffres (Doc 3) ? Ai-je défini « philosophie », « méthode », « objet » ? La conclusion répond-elle à la question posée ?
\end{enumerate}
\end{methode}

\newpage

\coursec{Méthodes et exercices résolus}

\courssub{Méthodes et exercices résolus}

\begin{methode}[Analyser et définir un concept philosophique]
\textbf{Objectif :} Maîtriser la définition philosophique (genre + différence spécifique) et distinguer le sens commun du sens savant.
\begin{enumerate}
    \item \textbf{Identifier le concept} : Souligner le terme à définir dans le sujet.
    \item \textbf{Rechercher le genre} : À quelle grande catégorie appartient-il ? (ex. : activité, discipline, attitude, savoir).
    \item \textbf{Déterminer la différence spécifique} : Qu'est-ce qui le distingue des autres éléments du même genre ? (ex. : \emph{par la raison}, \emph{par la radicalité}, \emph{par l'universalité}).
    \item \textbf{Formuler la définition} : \texttt{Concept = Genre + Différence spécifique}.
    \item \textbf{Illustrer par un contre-exemple} : Montrer ce que le concept n'est pas (opinion, science particulière, croyance).
\end{enumerate}
\cle{Genre} \cle{Différence spécifique} \cle{Contre-exemple}
\end{methode}

\begin{exoresolu}[Définir la philosophie par rapport à l'opinion et au savoir scientifique]
\textbf{Énoncé :} « La philosophie est-elle un savoir comme les autres ? » En vous appuyant sur la distinction entre \emph{opinion} (\emph{doxa}), \emph{science} (\emph{epistémè}) et \emph{philosophie}, rédigez une réponse structurée en trois paragraphes.

\tcblower
\textbf{Solution détaillée :}

\textbf{Paragraphe 1 : L'opposition à l'opinion (Doxa).}
La philosophie se définit d'abord \emph{négativement} par sa rupture avec l'opinion commune (\emph{doxa}). L'opinion est un jugement subjectif, non fondé, variable selon les individus, les époques ou les cultures (ex. : « Au Cameroun, on pense que... »). Elle repose sur le préjugé, la tradition ou l'émotion. La philosophie, née en Grèce (passage du \emph{mythos} au \emph{logos}), exige la \textbf{rationnalité} : elle substitue l'argumentation au récit, la preuve à l'affirmation gratuite. \emph{Donc, la philosophie n'est pas une opinion : elle en est l'examen critique.}

\textbf{Paragraphe 2 : La distinction d'avec les sciences particulières (Epistémè).}
Les sciences (physique, biologie, sociologie) sont des \emph{savoirs positifs} : elles étudient un \textbf{objet déterminé} (la matière, le vivant, le social) par une \textbf{méthode spécifique} (expérience, calcul, observation) et visent des \textbf{résultats cumulatifs et vérifiables}. La philosophie, elle, n'a pas d'objet particulier : son objet est \textbf{universel} (l'Être, le Vrai, le Bien, l'Homme, le Sens). Sa méthode n'est pas l'expérience contrôlée mais \textbf{la réflexion conceptuelle, l'analyse logique et la critique des présupposés}. Elle ne produit pas de « résultats » techniques mais de la \textbf{lucidité}. \emph{Donc, la philosophie n'est pas une science particulière : elle est la réflexion sur les conditions de possibilité et le sens des sciences.}

\textbf{Paragraphe 3 : La spécificité du savoir philosophique.}
La philosophie est un \textbf{savoir de second degré} (métasavoir) : elle interroge les fondements, les méthodes et les limites des autres savoirs. Elle est \textbf{radicale} (va à la racine), \textbf{universelle} (concerne tout homme) et \textbf{critique} (ne prend rien pour acquis). Comme le dit Kant, « apprendre à philosopher, c'est apprendre à penser par soi-même » (\emph{Sapere aude}). C'est un \textbf{savoir pratique} qui vise la sagesse (philo-sophia = amour de la sagesse), c'est-à-dire une orientation de l'existence.

\textbf{Conclusion :} La philosophie n'est pas un savoir « comme les autres » (ni opinion, ni science positive). C'est une \textbf{activité critique et rationnelle} qui interroge le sens global de l'existence et les fondements des connaissances particulières.
\end{exoresolu}

\begin{methode}[Expliquer l'origine et l'objet de la philosophie (Étonnement \textrightarrow{} Problématisation)]
\textbf{Objectif :} Retracer le cheminement de l'attitude philosophique : de l'expérience vécue à la question conceptuelle.
\begin{enumerate}
    \item \textbf{L'origine psychologique : L'étonnement (\emph{Thaumazein}).} Citer Platon (\emph{Théétète}) et Aristote (\emph{Métaphysique}) : « C'est par l'étonnement que les hommes commencent à philosopher. » Préciser : ce n'est pas la simple surprise (éphémère), mais un \textbf{étonnement métaphysique} face à l'évidence du monde (Pourquoi y a-t-il quelque chose plutôt que rien ? La mort, le mal, la justice).
    \item \textbf{Le moteur : Le doute et la recherche de sens.} L'étonnement suscite l'inquiétude (l'ignorance prise conscience). La philosophie naît du désir de \textbf{comprendre} (intelligere = lire au-dedans) et non seulement de connaître des faits.
    \item \textbf{L'objet formel : L'universel et le premier.} La philosophie n'étudie pas \emph{ceci} ou \emph{cela} (objets matériels), mais \textbf{l'Être en tant qu'Être} (Aristote), les \textbf{premiers principes} et les \textbf{premières causes}. Elle vise le \textbf{Sens} (direction, signification, valeur).
    \item \textbf{La formulation du problème :} Transformer l'étonnement vague en \textbf{question philosophique} (ouverte, universelle, débattue, sans réponse technique unique). Ex. : « La mort » (fait) \textrightarrow{} « La vie a-t-elle un sens face à la mort ? » (problème).
\end{enumerate}
\cle{Étonnement métaphysique} \cle{Problématisation} \cle{Objet formel}
\end{methode}

\begin{exoresolu}[De l'étonnement quotidien à la question philosophique]
\textbf{Énoncé :} Un élève de Première Technique observe : « Dans mon quartier à Douala, les jeunes passent le Bac, mais beaucoup restent au chômage. C'est injuste. » Transformez cette constatation en une démarche philosophique en 3 étapes : 1. Identifier l'étonnement initial. 2. Formuler la question philosophique. 3. Identifier l'objet philosophique visé.

\tcblower
\textbf{Solution détaillée :}

\textbf{Étape 1 : L'étonnement initial (Le fait brut).}
L'élève constate un \textbf{fait social} : inadéquation entre formation diplômante et insertion professionnelle. Il ressent un \textbf{sentiment d'injustice} (réaction affective morale). C'est le point de départ : l'évidence sociale (« ça ne va pas ») devient interrogative (« pourquoi ça ne va pas ? »).

\textbf{Étape 2 : La problématisation (La question philosophique).}
Il faut dépasser l'explication sociologique (taux de croissance, adéquation formation-emploi) ou la plainte morale pour atteindre l'universel.
\begin{itemize}
    \item \emph{Mauvaise question :} « Comment créer des emplois au Cameroun ? » (Question technique/économique).
    \item \emph{Bonne question philosophique :} \textbf{« La justice sociale se réduit-elle à l'égalité des chances scolaires ? »} ou \textbf{« Le mérite diplômique fonde-t-il un droit à l'emploi ? »}
\end{itemize}
Ces questions sont \textbf{universelles} (valent pour toute société), \textbf{normatives} (interrogent le \emph{Devoir-être} par rapport à l'\emph{Être}), et \textbf{ouvertes} (plusieurs thèses possibles : méritocratie, droit à l'existence, justice comme équité rawlsienne, etc.).

\textbf{Étape 3 : L'objet philosophique visé.}
L'objet n'est pas « le chômage au Cameroun » (objet de la sociologie/économie), mais \textbf{la Justice} (valeur universelle), plus précisément \textbf{la Justice distributive} (critères de répartition des biens sociaux : mérite, besoin, égalité). La philosophie ici éclaire le concept pour guider l'action politique future.

\textbf{Conclusion :} La démarche philosophique transforme un \textbf{ressenti singulier} en \textbf{interrogation universelle} sur les fondements normatifs de la vie collective.
\end{exoresolu}

\begin{methode}[Maîtriser la méthode philosophique : Problématiser, Argumenter, Conceptualiser]
\textbf{Objectif :} Appliquer les trois piliers de la méthode au Bac (Dissertation / Commentaire).
\begin{enumerate}
    \item \textbf{Problématiser (Le cœur de la méthode).}
    \begin{itemize}
        \item Refuser l'évidence / le lieu commun.
        \item Dégager la \textbf{tension} (antinomie) : deux thèses opposées plausibles (ex. : Liberté vs Déterminisme ; Devoir vs Inclination).
        \item Formuler le \textbf{problème} sous forme de question directe : « Peut-on... ? », « Faut-il... ? », « Est-ce que... ? ».
    \end{itemize}
    \item \textbf{Argumenter (La logique du discours).}
    \begin{itemize}
        \item \textbf{Thèse} : Affirmation répondant au problème.
        \item \textbf{Argument} : Raison logique reliant la thèse au principe (autorité, analogie, causalité, \emph{a fortiori}, \emph{reductio ad absurdum}).
        \item \textbf{Exemple/Contre-exemple} : Ancrage concret (historique, littéraire, vécu camerounais) pour tester la thèse.
        \item \textbf{Objection / Antithèse} : Point de vue adverse pris au sérieux.
        \item \textbf{Synthèse / Dépassement} : Intégration partielle ou dépassement dialectique (concept plus large).
    \end{itemize}
    \item \textbf{Conceptualiser (L'outillage).}
    \begin{itemize}
        \item Définir rigoureusement les concepts-clés du sujet (Genre + Différence spécifique).
        \item Utiliser les \textbf{distinctions conceptuelles} (ex. : Liberté \emph{de} / Liberté \emph{pour} ; Loi \emph{naturelle} / Loi \emph{positive} ; Justice / Légalité).
        \item Mobiliser les \textbf{références autorisées} (Auteurs au programme : Platon, Aristote, Descartes, Kant, Bergson, Sartre, auteurs africains : Mudimbe, Wiredu, Mbembe, etc.) pour \emph{appuyer}, non pour \emph{remplacer} sa propre pensée.
    \end{itemize}
\end{enumerate}
\cle{Problématisation} \cle{Argumentation} \cle{Conceptualisation} \cle{Dialectique}
\end{methode}

\begin{exoresolu}[Structurer une dissertation : « La technique libère-t-elle l'homme ? »]
\textbf{Énoncé :} Sujet type Probatoire. Proposez un plan dialectique détaillé (Introduction, 3 parties avec sous-parties argumentées, Conclusion) en mobilisant la distinction \emph{moyen/fin} et l'aliénation technique.

\tcblower
\textbf{Solution détaillée :}

\textbf{INTRODUCTION}
\begin{itemize}
    \item \textbf{Amorce :} De l'outil préhistorique à l'IA, la technique façonne l'humanité.
    \item \textbf{Définitions :} Technique = ensemble de procédés rationnels pour transformer la nature selon nos fins (moyen). Libérer = affranchir de la contrainte, de la nécessité, donner de l'autonomie.
    \item \textbf{Problématisation :} La technique apparaît d'abord comme \textbf{libération} (maîtrise de la nature, gain de temps, confort). Mais ne risque-t-elle pas de devenir une \textbf{nouvelle servitude} (aliénation, dépendance, inversion moyen/fin) ?
    \item \textbf{Annonce du plan :} 1. La technique comme puissance d'agir (libération). 2. La technique comme système d'asservissement (aliénation). 3. Vers une maîtrise critique : la technique comme médiation éthique.
\end{itemize}

\textbf{PREMIÈRE PARTIE : La technique, puissance de libération (Maîtrise de la nécessité)}
\begin{itemize}
    \item \textbf{Argument 1 : Libération des contraintes naturelles.} L'outil prolonge le corps (Bergson : \emph{L'Évolution créatrice}, l'homme \emph{homo faber}). Le feu, la charrue, le vaccin libèrent de la faim, du froid, de la maladie. Ex. : L'irrigation au Nord-Cameroun libère de la dépendance pluviométrique.
    \item \textbf{Argument 2 : Gain de temps et accès à la culture.} Automatisation des tâches répétitives (Aristote : l'esclave est un « outil animé » ; la machine libère l'homme pour la \emph{scholè} / loisir studieux).
    \item \textbf{Argument 3 : Extension de la liberté politique.} Imprimerie, internet $\rightarrow$ circulation des idées, résistance aux tyrannies (Printemps arabes, lanceurs d'alerte).
\end{itemize}

\textbf{DEUXIÈME PARTIE : La technique, risque d'aliénation (L'inversion moyen/fin)}
\begin{itemize}
    \item \textbf{Argument 1 : L'aliénation marxiste.} Le prolétaire ne possède pas les moyens de production ; la machine dicte son rythme (Taylorisme). L'homme devient « appendice de la machine » (\emph{Le Capital}). Ex. : Cadences infernales dans les usines de transformation du cacao/bois.
    \item \textbf{Argument 2 : L'enchaînement technique (Ellul / Heidegger).} La technique devient \emph{La Technique} : système autonome, efficacité comme seule valeur (\emph{Gestell} / Engeignement). On ne choisit plus les fins, on s'adapte aux moyens. Ex. : « Il faut numériser l'administration » (moyen) devient une fin en soi, oubliant le service public.
    \item \textbf{Argument 3 : Perte de savoir-faire et dépendance.} « Black box » : l'utilisateur ignore le fonctionnement (smartphone, algorithmes de notation bancaire). Perte d'autonomie pratique (Illich : \emph{La Convivialité}).
\end{itemize}

\textbf{TROISIÈME PARTIE : Dépasser l'alternative : La technique comme médiation éthique}
\begin{itemize}
    \item \textbf{Argument 1 : La neutralité instrumentale relativisée.} La technique n'est pas \emph{neutre} (portée de valeurs : efficacité, rentabilité), mais elle n'est pas \emph{fatale}. L'homme reste \textbf{sujet moral} (Sartre : « Nous sommes nos choix »).
    \item \textbf{Argument 2 : L'éthique de la responsabilité (Jonas).} Principe responsabilité : « Agis de sorte que les effets de ton action soient compatibles avec la permanence d'une vie authentiquement humaine. » Évaluation \emph{ex ante} des impacts (écologie, IA).
    \item \textbf{Argument 3 : L'appropriation critique (Technologies appropriées / Low-tech).} Choisir des techniques \emph{conviviales} (Illich), réparables, locales. Ex. : Fours solaires, biogaz, logiciels libres au Cameroun : la technique redevient moyen au service d'une fin humaine choisie.
\end{itemize}

\textbf{CONCLUSION}
\begin{itemize}
    \item \textbf{Bilan :} La technique libère de la \emph{nécessité naturelle} mais peut asservir par la \emph{nécessité technique}.
    \item \textbf{Réponse :} Elle libère \textbf{si et seulement si} l'homme reste le \textbf{législateur des fins} (Kant : l'humanité comme fin, jamais comme moyen). La libération technique appelle une \textbf{libération politique et éthique}.
    \item \textbf{Ouverture :} Le défi de l'IA générative : délégation cognitive vs augmentation humaine.
\end{itemize}
\end{exoresolu}

\begin{methode}[Démontrer l'importance de la philosophie (Formation de l'esprit critique \textrightarrow{} Citoyenneté)]
\textbf{Objectif :} Justifier l'utilité de la philosophie non par son « utilité marchande » mais par sa \emph{finalité intrinsèque} (formation du sujet).
\begin{enumerate}
    \item \textbf{Plan intellectuel : L'esprit critique.}
    \begin{itemize}
        \item Lutte contre les \emph{idoles} (Bacon) / \emph{préjugés} (Descartes) / \emph{idéologies} (Marx) / \emph{fake news} actuelles.
        \item Exercice du \textbf{doute méthodique} : ne rien accepter sans raison valide.
        \item Maîtrise de la \textbf{logique} : repérer les sophismes (homme de paille, \emph{ad hominem}, \emph{post hoc ergo propter hoc}).
    \end{itemize}
    \item \textbf{Plan moral et politique : L'autonomie (Kant).}
    \begin{itemize}
        \item \emph{Sapere aude} : « Aie le courage de te servir de ton propre entendement. »
        \item Sortie de la \emph{minorité} (tutelle : tradition, autorité, mode, algorithmes).
        \item Fondement de la \textbf{citoyenneté active} : juger les lois, voter en conscience, résister à l'injustice (Antigone, Thoreau, Mandela).
    \end{itemize}
    \item \textbf{Plan existentiel : Le sens de la vie (Sagesse).}
    \begin{itemize}
        \item « La vie non examinée ne vaut pas d'être vécue » (Socrate, \emph{Apologie}).
        \item Gestion de l'angoisse, de la mort, de l'échec (Stoïcisme, Existentialisme).
        \item Construction d'un \textbf{projet de vie} cohérent (valeurs hiérarchisées).
    \end{itemize}
    \item \textbf{Plan épistémologique : Régulation des sciences.}
    \begin{itemize}
        \item Épistémologie : valider les modèles, clarifier les concepts (espace, temps, vie, information).
        \item Éthique des sciences : limites de la manipulation du vivant, IA, climat (Jonas).
    \end{itemize}
\end{enumerate}
\cle{Esprit critique} \cle{Autonomie} \cle{Sagesse} \cle{Régulation}
\end{methode}

\begin{exoresolu}[Justifier la place de la philosophie dans un cursus technique]
\textbf{Énoncé :} Un élève en Première F4 (Génie Civil) demande : « À quoi me sert la philosophie pour construire des ponts ? » Répondez en distinguant \emph{compétence technique} et \emph{compétence éthique/épistémique}, avec un exemple concret de risque d'effondrement.

\tcblower
\textbf{Solution détaillée :}

\textbf{1. La distinction fondamentale : Savoir-faire technique vs Savoir-être ingénieur.}
\begin{itemize}
    \item \textbf{Compétence technique (Le « Comment »)} : Calcul de résistance des matériaux (RDM), normes Eurocodes/BAEL, logiciels (Robot, ArchiCAD), gestion de chantier. \emph{C'est le cœur du métier, nécessaire mais insuffisant.}
    \item \textbf{Compétence philosophique (Le « Pourquoi » et le « Jusqu'où »)} :
    \begin{itemize}
        \item \textbf{Épistémologie :} Comprendre les \textbf{limites des modèles}. Un calcul RDM suppose des hypothèses (matériau homogène, sollicitations connues). La philosophie des sciences apprend l'\emph{humilité épistémique} : la carte n'est pas le territoire. \textbf{Exemple :} L'effondrement du viaduc de Polcevera à Gênes (Italie, 2018) ou d'immeubles à Lagos (Nigeria, 2019) : souvent dû à une \emph{confiance aveugle dans le modèle} ou l'ignorance de la corrosion (réalité physique) vs le calcul (modèle idéal).
        \item \textbf{Éthique de la responsabilité (Hans Jonas) :} L'ingénieur engage la \textbf{vie d'autrui}. « Agis de sorte que les effets de ton action ne détruisent pas la possibilité de la vie humaine. » Le choix d'un acier moins cher (corruption/pression économique) vs la sécurité : c'est un \textbf{choix de valeur}, pas un calcul.
        \item \textbf{Critique de l'idéologie technicienne :} Résister à la pression du « toujours plus vite, moins cher » (efficacité) au profit du « durable, sûr, juste » (valeurs humaines).
    \end{itemize}
\end{itemize}

\textbf{2. Exemple concret : Le choix des fondations dans un sol argileux (Yaoundé/Nkolbisson).}
\begin{itemize}
    \item \textbf{Approche purement technique (à court terme) :} Semelles superficielles (moins chères, plus rapides). Le logiciel valide selon les normes actuelles.
    \item \textbf{Approche philosophique (réflexivité) :} L'ingénieur se demande : « Ce modèle intègre-t-il le changement climatique (sécheresses intenses $\rightarrow$ retrait-gonflement des argiles) ? Suis-je responsable si la maison fissure dans 10 ans ? » $\rightarrow$ Décision : Pieux profonds (surcoût) mais sécurité long terme.
\end{itemize}

\textbf{3. Conclusion pour l'élève.}
La philosophie ne te donne pas la \textbf{formule du béton}, elle te donne la \textbf{rigueur conceptuelle} pour ne pas confondre le modèle et la réalité, et le \textbf{courage moral} pour refuser un chantier dangereux. C'est la différence entre un \emph{exécutant} et un \emph{ingénieur responsable}.
\end{exoresolu}

\begin{methode}[Différencier Philosophie et Science (Critères : Objet, Méthode, Vérité, Finalité)]
\textbf{Objectif :} Construire un tableau comparatif rigoureux pour ne pas confondre les deux démarches (écueil classique au Bac).
\begin{enumerate}
    \item \textbf{L'Objet :}
    \begin{itemize}
        \item Science : \textbf{Particulier / Régional} (un domaine délimité : la cellule, l'atome, le marché). \emph{Fait}.
        \item Philosophie : \textbf{Universel / Totalisant} (l'Être, la Vérité, le Bien, le Sujet, le Sens). \emph{Valideur de sens}.
    \end{itemize}
    \item \textbf{La Méthode :}
    \begin{itemize}
        \item Science : \textbf{Hypothético-déductive / Expérimentale}. Observation $\rightarrow$ Hypothèse $\rightarrow$ Expérience/Calcul $\rightarrow$ Validation/Falsification (Popper). \emph{Inter-subjectivité} (reproductibilité).
        \item Philosophie : \textbf{Conceptuelle / Dialectique / Phénoménologique}. Analyse de concepts, expérience de pensée, herméneutique, argumentation rationnelle. \emph{Universalisation} (valable pour tout sujet rationnel).
    \end{itemize}
    \item \textbf{La Vérité / Validité :}
    \begin{itemize}
        \item Science : \textbf{Vérité provisoire, cumulative, falsifiable}. Progrès = accumulation/correction (Newton $\rightarrow$ Einstein).
        \item Philosophie : \textbf{Validité rationnelle, non cumulative}. On ne « résout » pas un problème philosophique comme une équation ; on l'\emph{éclaircit}. Platon/Aristote restent pertinents. Progrès = approfondissement de la lucidité.
    \end{itemize}
    \item \textbf{La Finalité :}
    \begin{itemize}
        \item Science : \textbf{Pouvoir / Prédiction / Maîtrise technique} (Bacon : « Savoir, c'est pouvoir »).
        \item Philosophie : \textbf{Sagesse / Autonomie / Sens} (Kant : « Que puis-je savoir ? Que dois-je faire ? Que m'est-il permis d'espérer ? »).
    \end{itemize}
    \item \textbf{Rapport :} \textbf{Non-hiérarchique mais distinct.} La science a besoin de la philosophie (épistémologie, éthique) pour se situer ; la philosophie a besoin des sciences (données) pour ne pas spéculer dans le vide.
\end{enumerate}
\cle{Objet régional vs Universel} \cle{Falsifiabilité} \cle{Non-cumulativité} \cle{Épistémologie}
\end{methode}

\begin{exoresolu}[Analyser un texte : « La science ne nous apprend pas ce que nous devons faire » (Max Weber, \emph{Le Savant et le Politique})]
\textbf{Énoncé :} Commentaire de texte type Bac. Identifiez la thèse, l'argumentation, les concepts clés (fait/valeur, rationalisation, désenchantement) et discutez la portée.

\tcblower
\textbf{Solution détaillée :}

\textbf{1. Présentation / Thèse.}
\textbf{Auteur/Contexte :} Max Weber (1864-1920), sociologue allemand, conférence de Munich 1917. \textbf{Thèse centrale :} \textbf{La science est incompétente en matière de valeurs.} Elle dit ce \emph{est} (fait), non ce qui \emph{doit être} (valeur). « La science est sans valeur » (\emph{Wertfreiheit}).

\textbf{2. Analyse structurée de l'argumentation.}
\begin{itemize}
    \item \textbf{Argument 1 : Spécificité de la rationalité scientifique.} La science procède par \textbf{réduction causale} et \textbf{calcul des moyens}. Elle répond : « Par quels moyens atteindre une fin donnée ? » Elle ne peut pas dire : « Quelle fin vaut la peine d'être poursuivie ? »
    \item \textbf{Argument 2 : Le fait ne contient pas la valeur (Hume/Guillaumin).} Aucune déduction logique ne permet de passer de « \emph{est} » à « \emph{doit} ». Ex. : La biologie décrit la fécondation ; elle ne dit pas si l'avortement est moral ou immoral.
    \item \textbf{Argument 3 : Le polythéisme des valeurs.} Le monde moderne est « désenchanté » : les valeurs suprêmes (Dieu, Patrie, Progrès, Art, Amour) sont en \textbf{conflit irréductible}. La science ne peut pas trancher entre elles (ex. : Efficacité économique vs Justice sociale).
    \item \textbf{Argument 4 : La responsabilité du politique.} Le savant fournit les \textbf{moyens techniques} (l'arme nucléaire, l'algorithme de notation). Le politique assume le \textbf{choix des fins} (la paix, la justice, le contrôle social). Confondre les deux = \emph{technocratie} (le savant décide) ou \emph{populisme} (on nie les faits scientifiques).
\end{itemize}

\textbf{3. Concepts clés à définir.}
\begin{itemize}
    \item \cle{Séparation fait/valeur} (Naturalistic fallacy).
    \item \cle{Rationalisation} : emprise du calcul moyen/fin sur tous les domaines.
    \item \cle{Désenchantement du monde} : perte de sens unifié, pluralisme des valeurs.
    \item \cle{Éthique de conviction vs Éthique de responsabilité} (politique).
\end{itemize}

\textbf{4. Discussion critique (Portée et limites).}
\begin{itemize}
    \item \textbf{Portée :} Indispensable garde-fou contre le \emph{scientisme} (croire que la science dicte la morale) et la \emph{technocratie}. Protège la démocratie : les choix de société appartiennent aux citoyens, pas aux experts.
    \item \textbf{Limite 1 (Épistémologique) :} La science n'est pas si « neutre ». Le choix des objets d'étude, des indicateurs (PIB vs IDH), des catégories (race, genre) est \textbf{chargé de valeurs} (Bourdieu, Foucault, épistémologie féministe).
    \item \textbf{Limite 2 (Éthique) :} Le savant n'est pas un automate. Il a une \textbf{responsabilité éthique} sur l'usage de ses découvertes (physiciens de Manhattan, biologistes CRISPR). La clause de conscience existe.
    \item \textbf{Limite 3 (Pratique) :} Dans l'urgence (crise climatique, pandémie), la distinction s'effrite : la science \emph{impose} des contraintes de réalité qui \emph{forcent} la main au politique (ex. : confinement). La frontière « fait/valeur » devient poreuse.
\end{itemize}

\textbf{Conclusion :} Weber a raison sur le \textbf{plan logique} (séparation des registres), mais il faut nuancer sur le \textbf{plan pratique} : science et politique s'interpénètrent dans l'action publique. La philosophie est l'instance de médiation.
\end{exoresolu}

\begin{methode}[Synthèse : Résoudre un sujet type Bac (Dissertation ou Commentaire) en 4 temps]
\textbf{Objectif :} Méthode unique pour l'épreuve écrite (4h).
\begin{enumerate}
    \item \textbf{Temps 1 : Analyse du sujet (30-45 min) - \emph{Brouillon}.}
    \begin{itemize}
        \item \textbf{Dissertation :} Définir \textbf{tous} les termes (dictionnaire + philosophie). Identifier le \textbf{type de sujet} (Paradoxe, Question fermée, Question ouverte, Citation). \textbf{Problématiser} : trouver la tension (Thèse plausible vs Antithèse plausible). Rédiger la \textbf{problématique} (une question). Faire le \textbf{plan détaillé} (3 parties, arguments, exemples, auteurs).
        \item \textbf{Commentaire :} Identifier \textbf{thèse}, \textbf{structure} (étapes), \textbf{concepts clés}, \textbf{arguments}, \textbf{exemples}, \textbf{références implicites}. \textbf{Problématiser} le texte (quel problème l'auteur résout-il ?). Faire le \textbf{plan d'explication} (suivre le fil du texte) + \textbf{plan de discussion} (critique/approfondissement).
    \end{itemize}
    \item \textbf{Temps 2 : Rédaction de l'Introduction (20 min) - \emph{Copie}.}
    \begin{itemize}
        \item Accroche (optionnelle, brève). Définitions des termes. \textbf{Problématique claire}. \textbf{Annonce du plan} (verbes d'action : Examiner, Analyser, Montrer, Dépasser).
    \end{itemize}
    \item \textbf{Temps 3 : Rédaction du Développement (2h30) - \emph{Copie}.}
    \begin{itemize}
        \item \textbf{Une idée = Un paragraphe.} Structure : \textbf{Thèse (phrase affirmative)} $\rightarrow$ \textbf{Argumentation} (logique, concepts, auteurs) $\rightarrow$ \textbf{Exemple concret} (Cameroun, histoire, actualité) $\rightarrow$ \textbf{Transition partielle} (limite $\rightarrow$ partie suivante).
        \item \textbf{Transition majeure} entre grandes parties : Résumer l'apport + Montrer l'insuffisance $\rightarrow$ Lancer la partie suivante.
    \end{itemize}
    \item \textbf{Temps 4 : Conclusion + Relecture (15-20 min).}
    \begin{itemize}
        \item \textbf{Bilan synthétique} (réponse nuancée à la problématique).
        \item \textbf{Ouverture} : Autre domaine, auteur non cité, actualité brûlante, question nouvelle. \textbf{Pas de nouvelle thèse.}
        \item \textbf{Relecture :} Orthographe, syntaxe, cohérence, \textbf{noms d'auteurs/titres en italique}, concepts en \textbf{gras} (si autorisé) ou soulignés.
    \end{itemize}
\end{enumerate}
\cle{Problématisation} \cle{Plan dialectique} \cle{Paragraphe type} \cle{Ouverture}
\end{methode}

\begin{exoresolu}[Sujet complet : « L'homme est-il un animal comme les autres ? » (Plan et intro rédigée)]
\textbf{Énoncé :} Traitez ce sujet de dissertation en appliquant la méthode des 4 temps. Fournissez : 1. L'analyse des termes et la problématique. 2. Le plan détaillé (3 parties, 2-3 args/partie). 3. L'introduction rédigée.

\tcblower
\textbf{Solution détaillée :}

\textbf{1. Analyse et Problématisation.}
\begin{itemize}
    \item \textbf{Homme :} Espèce biologique \emph{Homo sapiens} (animalité : naissance, mort, besoins, sexualité, génome 98,7\% partagé avec le chimpanzé).
    \item \textbf{Animal :} Être vivant, sensible, mobile, hétrotrophe, gouverné par l'instinct.
    \item \textbf{Comme les autres :} Similitude de nature, continuité, absence de rupture ontologique.
    \item \textbf{Tension :} \textbf{Continuité biologique} (Darwin, éthologie : culture, outils, deuil chez primates) \textbf{VS Rupture anthropologique} (Langage symbolique, Technique, Droit, Morale, Conscience de soi, Mortalité assumée, Histoire).
    \item \textbf{Problématique :} \textbf{« La continuité biologique suffit-elle à définir l'homme, ou la rupture symbolique et culturelle constitue-t-elle une différence de nature, et non pas seulement de degré ? »}
\end{itemize}

\textbf{2. Plan détaillé.}

\textbf{PARTIE I : L'homme \emph{est} un animal (Continuité naturaliste).}
\begin{itemize}
    \item A. Un corps animal : évolution darwinienne, génétique partagée, besoins vitaux, vulnérabilité, mortalité. (Darwin, \emph{La Filiation de l'homme} ; Biologie).
    \item B. Des comportements « animaux » partagés : socialité, usage d'outils, transmission culturelle, émotions, politique (chimpanzés : De Waal). Différence de \textbf{degré} (complexité), pas de nature. (Éthologie).
    \item C. Le naturalisme philosophique : L'homme n'est pas un « empire dans un empire » (Spinoza, \emph{Éthique}). Pas d'âme substantielle distincte. Monisme : l'esprit émerge de la matière cérébrale.
\end{itemize}

\textbf{PARTIE II : L'homme n'est \emph{pourtant} pas un animal \emph{comme les autres} (Rupture anthropologique).}
\begin{itemize}
    \item A. Le langage \textbf{symbolique} (pas seulement communication) : syntaxe récursive, référence absente, méta-langage, écriture, accumulation transgénérationnelle. (Benveniste, Chomsky, Derrida : « L'animal ne répond pas »).
    \item B. La \textbf{technique} comme extériorisation de la mémoire et projection vers l'avenir (Leroi-Gourhan, Simondon, Stiegler). L'animal s'adapte au milieu ; l'homme \textbf{transforme} le milieu (niche écologique $\rightarrow$ niche technologique).
    \item C. La \textbf{liberté} et la \textbf{responsabilité} : L'homme se \textbf{déprend} de l'instinct (Sartre : « L'existence précède l'essence »). Il se choisit, il est \textbf{projet}. Il a une \textbf{histoire}, pas seulement une évolution.
\end{itemize}

\textbf{PARTIE III : L'homme, un animal \textbf{singulier} : la « nature dénaturée » (Plessner / Gehlen / Merleau-Ponty).}
\begin{itemize}
    \item A. La \textbf{positionnalité excentrique} (Plessner) : L'animal vit \emph{dedans} son milieu (centrique). L'homme vit \emph{devant} lui (excentrique) : il a un corps et il \emph{est} un corps. Double appartenance : nature et culture.
    \item B. L'\textbf{instinct de dénégation} (Freud) / \textbf{Animal malade} (Nietzsche) / \textbf{Être de manque} (Gehlen) : L'homme n'est pas « programmé ». Son instinct est \emph{indéterminé} (plasticité cérébrale). Il \textbf{doit} se créer (culture = second nature).
    \item C. \textbf{Éthique :} Cette singularité fonde la \textbf{dignité} (Kant : fin en soi) et les \textbf{Droits de l'Homme}, distincts du bien-être animal. Responsabilité \emph{sur} la nature (écologie).
\end{itemize}

\textbf{3. Introduction rédigée.}
« L'homme est un animal », affirme la biologie moderne depuis Darwin. Le séquençage du génome confirme notre parenté étroite avec le chimpanzé (98,7\% d'ADN partagé). L'éthologie (Jane Goodall, Frans de Waal) brouille chaque jour un peu plus la frontière : outils, culture, politique, empathie existent chez nos cousins primates. Pourtant, nul chimpanzé n'a écrit de traité de morale, construit de cathédrale, ou jugé un criminel pour crime contre l'humanité. L'homme se distingue par une \textbf{rupture} : le langage symbolique, la technique cumulative, le droit, la conscience de la mort et la liberté de se choisir. \textbf{Dès lors, la continuité biologique suffit-elle à définir l'homme, ou la rupture symbolique et culturelle constitue-t-elle une différence de nature, et non pas seulement de degré ?} Nous examinerons d'abord les fondements de l'animalité humaine (I), avant d'analyser les structures de la rupture anthropologique (II), pour enfin comprendre comment l'homme assume sa singularité comme une \textbf{tâche éthique} : celle de s'humaniser en se dénaturant (III).»
\end{exoresolu}

\begin{attention}[Les 5 erreurs qui coûtent des points au Bac (Probatoire / Baccalauréat Technique)]
\begin{enumerate}
    \item \textbf{Confondre « Exemple » et « Argument ».} \emph{Erreur :} « Par exemple, au Cameroun, il y a la corruption. » \emph{Correction :} L'exemple \textbf{illustre} un argument logique. Structure : « L'argument d'autorité montre que... (Logique). \textbf{Par exemple}, l'affaire [X] au Cameroun illustre ce mécanisme... » \textbf{Point clé :} Pas d'exemple « nu » sans lien explicite avec la thèse défendue.
    \item \textbf{Le « Name-dropping » (Citation décorative).} \emph{Erreur :} « Comme dit Kant, "Sapere aude". » (Puis on passe à autre chose). \emph{Correction :} \textbf{Expliquer} la citation : « Kant, dans \emph{Qu'est-ce que les Lumières ?}, définit les Lumières comme la sortie de la minorité : "Sapere aude". Cela signifie que l'autonomie philosophique exige le courage de penser par soi-même, sans tuteur. » \textbf{Point clé :} L'auteur \emph{appuie} votre raisonnement, il ne le \emph{remplace} pas.
    \item \textbf{Ignorer la problématisation (Plan « Catalogue » ou « Pour/Contre » binaire).} \emph{Erreur :} Partie I : « Oui, la technique libère. » Partie II : « Non, elle aliène. » Partie III : « Elle fait les deux. » \emph{Correction :} Chaque partie doit faire \textbf{avancer la pensée}. I. Libération par la maîtrise (niveau technique). II. Aliénation par le système technique (niveau structurel/ontologique). III. Dépassement par l'éthique de la responsabilité (niveau politique/existentiel). \textbf{Point clé :} Dialectique = Dépassement (Aufhebung), pas simple alternance.
    \item \textbf{Définitions floues ou absentes.} \emph{Erreur :} Traiter « La technique libère-t-elle l'homme ? » sans définir « technique » (moyen/fin, système), « libérer » (contrainte/autonomie), « homme » (sujet/espèce). \emph{Correction :} \textbf{Introduction = Contrat de lecture.} Définitions rigoureuses (Genre + Différence spécifique) des termes du sujet. Si le sens change dans le développement, le signaler explicitement (« Au sens fort de... »).
    \item \textbf{Oublier le contexte camerounais / l'ancrage concret.} \emph{Erreur :} Dissertation 100\% abstraite (Platon, Descartes, Kant uniquement). \emph{Correction :} Le programme exige : « Appliquer les notions à des situations concrètes de la vie courante ». \textbf{Exemples attendus :} Justice coutumière vs Justice étatique (Droit), Gestion de l'eau au Sahel (Éthique/Technique), Place des chefferies (Politique), Identité culturelle vs Modernité (Anthropologie), Corruption / « Arrangement » (Morale). \textbf{Point clé :} L'exemple local prouve que vous \emph{pensez} la philosophie, vous ne la \emph{récitez} pas.
\end{enumerate}
\end{attention}

\newpage

\coursec{Exercices}

\courssub{Je vérifie mes connaissances}

\begin{exercice}[QCM : les fondamentaux]
Pour chaque question, entoure la (ou les) bonne(s) réponse(s).
\begin{enumerate}
\item Le mot \emph{philosophie} vient du grec \emph{philein} (aimer) et \emph{sophia} (sagesse). Il signifie donc : \\
a) l'amour du pouvoir \quad b) l'amour de la sagesse \quad c) la science des dieux \quad d) l'art de convaincre
\item La philosophie serait née, selon la tradition, au VI\textsuperscript{e} siècle avant J.-C. : \\
a) à Rome \quad b) en Égypte \quad c) en Grèce (Milet) \quad d) en Chine
\item Parmi ces auteurs, lequel a défini la philosophie comme \og l'étonnement \fg{} qui pousse l'homme à chercher ? \\
a) Platon \quad b) Newton \quad c) Pasteur \quad d) Einstein
\item Selon Kant, la philosophie se résume en quatre questions. Laquelle n'en fait PAS partie ? \\
a) Que puis-je savoir ? \quad b) Que dois-je faire ? \quad c) Que m'est-il permis d'espérer ? \quad d) Combien puis-je gagner ?
\end{enumerate}
\end{exercice}

\begin{exercice}[Vrai ou faux ? Justifie ta réponse]
Réponds par \textbf{Vrai} ou \textbf{Faux}, puis justifie en une ou deux phrases.
\begin{enumerate}
\item La philosophie se contente d'opinions personnelles sans argumentation.
\item La science et la philosophie poursuivent exactement le même objet.
\item La réflexion philosophique utilise une méthode rigoureuse : problématisation, analyse, discussion, synthèse.
\item La philosophie n'a aucune utilité pratique dans la vie quotidienne.
\end{enumerate}
\end{exercice}

\begin{exercice}[Définitions à compléter]
Recopie et complète les définitions suivantes avec les mots de la liste : \emph{sagesse ; rationnelle ; universel ; méthode ; critique}.
\begin{enumerate}
\item Étymologiquement, la philosophie est l'amour de la \ldots\ldots\ldots
\item La philosophie est une réflexion \ldots\ldots\ldots et \ldots\ldots\ldots qui interroge le sens de l'existence, du monde et de l'action humaine.
\item La démarche philosophique exige une \ldots\ldots\ldots rigoureuse et un examen \ldots\ldots\ldots de toutes les opinions reçues.
\end{enumerate}
\end{exercice}

\begin{exercice}[Correspondance auteur--idée]
Relie chaque auteur à l'idée qui lui correspond.
\begin{center}
\begin{tabularx}{\linewidth}{|l|X|}
\hline
\textbf{Auteur} & \textbf{Idée} \\
\hline
1. Socrate & a) La philosophie répond à quatre questions fondamentales \\
\hline
2. Aristote & b) \og Une vie sans examen ne vaut pas la peine d'être vécue \fg{} \\
\hline
3. Kant & c) La philosophie naît de l'étonnement \\
\hline
4. Épicure & d) La philosophie guérit les troubles de l'âme \\
\hline
\end{tabularx}
\end{center}
\end{exercice}

\courssub{Je m'entraîne}

\begin{exercice}[Situations de la vie courante]
Voici trois situations observées au Cameroun :
\begin{enumerate}
\item[\textbf{A.}] À Yaoundé, une étudiante se demande si tricher à l'examen du Probatoire peut être justifié lorsque \og tout le monde le fait \fg{}.
\item[\textbf{B.}] Un chauffeur de taxi à Douala répare le moteur de sa voiture en suivant le manuel technique.
\item[\textbf{C.}] Un commerçant au marché Mokolo calcule son bénéfice de la journée.
\end{enumerate}
\begin{enumerate}
\item Identifie la situation qui relève d'une démarche philosophique.
\item Justifie ton choix en montrant pourquoi les deux autres situations n'en relèvent pas (utilise les critères : universalité de la question, réflexion sur les valeurs, argumentation).
\end{enumerate}
\end{exercice}

\begin{exercice}[Remettre la méthode dans l'ordre]
On a appliqué la démarche philosophique à la question \og Faut-il toujours dire la vérité ? \fg{}, mais les cinq étapes ci-dessous ont été mélangées. Remets-les dans l'ordre logique, puis rédige en cinq lignes la synthèse finale.
\begin{enumerate}
\item[\textbf{a.}] Confronter les deux positions : la vérité peut blesser ; le mensonge détruit la confiance.
\item[\textbf{b.}] Formuler clairement le problème : la vérité est-elle un devoir absolu ou dépend-elle des circonstances ?
\item[\textbf{c.}] Définir les notions clés : vérité, mensonge, devoir, conséquence.
\item[\textbf{d.}] Proposer une réponse nuancée et justifiée : dire la vérité est la règle, mais la manière et le moment comptent.
\item[\textbf{e.}] Examiner les arguments pour la position \og il faut toujours dire la vérité \fg{} (exemple : la confiance dans la famille ou au travail).
\end{enumerate}
\end{exercice}

\begin{exercice}[Philosophie ou science ?]
Pour chacune des questions suivantes, indique si elle relève principalement de la \textbf{science} ou de la \textbf{philosophie}, et justifie brièvement.
\begin{enumerate}
\item Quelle est la composition chimique de l'eau de source de la région de l'Ouest ?
\item L'homme est-il libre de ses choix ?
\item Quelle distance sépare Yaoundé de Bafoussam ?
\item Qu'est-ce qu'une société juste ?
\item Le progrès technique rend-il l'homme plus heureux ?
\end{enumerate}
\end{exercice}

\begin{exercice}[L'utilité de la philosophie au quotidien]
Un élève de Première technique affirme : \og La philosophie ne sert à rien, moi je veux devenir technicien. \fg{}
\begin{enumerate}
\item Donne deux arguments montrant que la philosophie est utile à un technicien (pense à l'éthique professionnelle, à l'esprit critique, à la prise de décision).
\item Illustre chaque argument par un exemple concret tiré de la vie professionnelle ou scolaire camerounaise.
\end{enumerate}
\end{exercice}

\courssub{Je me prépare au Bac}

\begin{exercice}[Dissertation type Baccalauréat technique]
\textbf{Sujet :} \emph{La philosophie est-elle utile à la vie de l'homme ?}

Certains affirment que la philosophie est un luxe inutile : elle ne produit ni biens, ni techniques, ni profits. D'autres soutiennent qu'elle est indispensable : elle apprend à penser par soi-même, à distinguer le vrai du faux, le juste de l'injuste, et donne un sens à l'existence.

\textbf{Travail à faire :}
\begin{enumerate}
\item Dégage la problématique du sujet et formule la question principale qu'il pose.
\item Construis un plan dialectique en trois parties : thèse (la philosophie est inutile), antithèse (la philosophie est indispensable), synthèse (ta position argumentée).
\item Rédige l'introduction complète de la dissertation (accroche, définition des termes, problématique, annonce du plan).
\item Rédige le développement de la première partie en mobilisant au moins un exemple concret de la vie courante camerounaise.
\end{enumerate}
\end{exercice}

\begin{exercice}[Commentaire de texte type Probatoire]
\textbf{Texte :} \og Tout l'intérêt de ma raison se réunit dans les trois questions suivantes : Que puis-je savoir ? Que dois-je faire ? Que m'est-il permis d'espérer ? [\ldots] La première question est purement spéculative. La seconde est purement pratique. La troisième [\ldots] conduit à cette conclusion que toute espérance concerne la félicité. \fg{} (Kant, \emph{Logique}, adapté)

\textbf{Travail à faire :}
\begin{enumerate}
\item Présente brièvement l'auteur et situe le texte dans le thème de la leçon.
\item Dégage la thèse du texte : que nous apprend Kant sur l'objet de la philosophie ?
\item Explique chacune des trois questions en donnant pour chacune un exemple concret tiré de ta vie d'élève ou de la société camerounaise.
\item Montre en quoi ces questions recouvrent les trois grandes dimensions de la philosophie : la connaissance, la morale, l'espérance.
\item Discussion : ces trois questions suffisent-elles à définir toute la philosophie ? Justifie ta réponse en cinq à dix lignes.
\end{enumerate}
\end{exercice}

\begin{exercice}[Question de réflexion argumentée]
\textbf{Sujet :} \emph{La science peut-elle répondre à toutes les questions de l'homme ?}

Dans un quartier de Douala, deux amis discutent. Le premier dit : \og Aujourd'hui, la science explique tout : les maladies, la météo, l'univers. Nous n'avons plus besoin de la philosophie. \fg{} Le second répond : \og La science dit \emph{comment} fonctionne le monde, mais elle ne dit pas \emph{pourquoi} il faut vivre ni \emph{comment} bien agir. \fg{}

\textbf{Travail à faire :}
\begin{enumerate}
\item Rappelle les différences entre la démarche scientifique et la démarche philosophique (objet, méthode, type de réponse).
\item Montre, à l'aide de deux exemples précis, ce que la science permet d'expliquer.
\item Montre, à l'aide de deux exemples précis, les questions qui échappent à la science et relèvent de la philosophie (valeurs, sens, liberté, justice).
\item Prends position dans le débat entre les deux amis et défends ton point de vue en un texte argumenté de quinze à vingt lignes, avec une introduction, deux arguments et une conclusion.
\end{enumerate}
\end{exercice}

\newpage

\coursec{Corrigés des exercices}

\begin{corrige}[Exercice 1 — QCM : les fondamentaux]
\begin{enumerate}
\item \textbf{Réponse b)} : l'amour de la sagesse. \emph{Philein} signifie \og aimer \fg{} et \emph{sophia} signifie \og sagesse \fg{} ou \og savoir \fg{}. Le philosophe n'est donc pas celui qui possède la sagesse, mais celui qui la recherche.
\item \textbf{Réponse c)} : en Grèce, à Milet (colonie grecque d'Asie Mineure), au VI\textsuperscript{e} siècle avant J.-C., avec Thalès, considéré comme le premier philosophe.
\item \textbf{Réponse a)} : Platon. Dans le \emph{Théétète}, il affirme que l'étonnement (\emph{thaumazein}) est l'origine de la philosophie ; Aristote reprendra cette idée dans la \emph{Métaphysique}.
\item \textbf{Réponse d)} : \og Combien puis-je gagner ? \fg{} ne fait pas partie des quatre questions kantiennes, qui sont : Que puis-je savoir ? Que dois-je faire ? Que m'est-il permis d'espérer ? Qu'est-ce que l'homme ?
\end{enumerate}
\textbf{Points de vigilance :} ne pas confondre \emph{sophia} (sagesse) avec \emph{sophiste} (maître de rhétorique) ; bien retenir que la quatrième question de Kant, \og Qu'est-ce que l'homme ? \fg{}, résume les trois autres.
\end{corrige}

\begin{corrige}[Exercice 2 — Vrai ou faux ?]
\begin{enumerate}
\item \textbf{Faux.} La philosophie se distingue précisément de l'opinion (\emph{doxa}) : elle exige une argumentation rationnelle, des preuves et un examen critique. Une opinion non justifiée n'est pas de la philosophie.
\item \textbf{Faux.} La science étudie des objets délimités et mesurables (la matière, le vivant) par l'expérience et le calcul ; la philosophie interroge le sens, les valeurs et les fondements (la justice, la liberté, la vérité) par la réflexion conceptuelle.
\item \textbf{Vrai.} La démarche philosophique suit des étapes rigoureuses : problématisation, définition des notions, analyse des arguments, discussion des thèses, synthèse argumentée.
\item \textbf{Faux.} La philosophie a une utilité pratique : elle forme l'esprit critique, éclaire les décisions morales et professionnelles (éthique), et aide à vivre en donnant un sens à l'existence, comme le montre Épicure pour qui elle \og guérit les troubles de l'âme \fg{}.
\end{enumerate}
\textbf{Points de vigilance :} une justification est exigée pour chaque réponse ; une réponse \og Vrai/Faux \fg{} seule, sans justification, ne vaut que la moitié des points.
\end{corrige}

\begin{corrige}[Exercice 3 — Définitions à compléter]
\begin{enumerate}
\item Étymologiquement, la philosophie est l'amour de la \textbf{sagesse}.
\item La philosophie est une réflexion \textbf{rationnelle} et \textbf{universelle} (critère : le mot \emph{universel} s'accorde ici au féminin : universelle) qui interroge le sens de l'existence, du monde et de l'action humaine.
\item La démarche philosophique exige une \textbf{méthode} rigoureuse et un examen \textbf{critique} de toutes les opinions reçues.
\end{enumerate}
\textbf{Points de vigilance :} chaque mot de la liste n'est utilisé qu'une seule fois ; attention à l'accord \emph{universelle} au féminin.
\end{corrige}

\begin{corrige}[Exercice 4 — Correspondance auteur--idée]
\begin{center}
\begin{tabularx}{\linewidth}{|l|X|}
\hline
\rowcolor{popblueL} \textbf{Auteur} & \textbf{Idée correspondante} \\
\hline
1. Socrate & \textbf{b)} \og Une vie sans examen ne vaut pas la peine d'être vécue \fg{} \\
\hline
2. Aristote & \textbf{c)} La philosophie naît de l'étonnement \\
\hline
3. Kant & \textbf{a)} La philosophie répond à quatre questions fondamentales \\
\hline
4. Épicure & \textbf{d)} La philosophie guérit les troubles de l'âme \\
\hline
\end{tabularx}
\end{center}
\textbf{Justifications :} Socrate, dans l'\emph{Apologie}, fait de l'examen de soi la condition d'une vie digne ; Aristote ouvre la \emph{Métaphysique} par l'étonnement comme origine du philosopher ; Kant formule ses quatre questions dans la \emph{Logique} ; Épicure compare la philosophie à une médecine de l'âme.
\textbf{Point de vigilance :} Platon et Aristote associent tous deux la philosophie à l'étonnement ; ici, c'est Aristote qui est proposé.
\end{corrige}

\begin{corrige}[Exercice 5 — Situations de la vie courante]
\begin{enumerate}
\item \textbf{La situation A} relève d'une démarche philosophique.
\item \textbf{Justification :}
\begin{itemize}
\item \textbf{Situation A (philosophique) :} la question \og la tricherie peut-elle être justifiée ? \fg{} porte sur une \textbf{valeur} (l'honnêteté, la justice) et possède une portée \textbf{universelle} : elle concerne tout homme, pas seulement cette étudiante. Elle exige une \textbf{argumentation} rationnelle et non un simple calcul d'intérêt.
\item \textbf{Situation B (technique) :} réparer un moteur en suivant un manuel relève du \emph{savoir-faire} : il s'agit d'appliquer des procédures efficaces pour obtenir un résultat matériel. La question posée est \og comment faire ? \fg{}, non \og que dois-je faire ? \fg{}.
\item \textbf{Situation C (calcul utilitaire) :} calculer un bénéfice relève de l'arithmétique commerciale : la réponse est chiffrée, certaine et particulière. Il n'y a ni interrogation sur les valeurs, ni universalité de la question.
\end{itemize}
\end{enumerate}
\textbf{Points de vigilance :} les trois critères (universalité, réflexion sur les valeurs, argumentation) doivent être explicitement mobilisés ; il ne suffit pas de dire \og A est philosophique \fg{} sans montrer pourquoi B et C ne le sont pas.
\end{corrige}

\begin{corrige}[Exercice 6 — Remettre la méthode dans l'ordre]
\textbf{Ordre logique : b $\rightarrow$ c $\rightarrow$ e $\rightarrow$ a $\rightarrow$ d.}
\begin{enumerate}
\item \textbf{b.} Formuler le problème (problématisation).
\item \textbf{c.} Définir les notions clés (analyse des concepts).
\item \textbf{e.} Examiner les arguments de la première position (thèse).
\item \textbf{a.} Confronter les deux positions (discussion dialectique).
\item \textbf{d.} Proposer une réponse nuancée et justifiée (synthèse).
\end{enumerate}
\textbf{Synthèse rédigée (cinq lignes) :} Dire la vérité est la règle fondamentale de la vie en société, car sans elle la confiance — au sein de la famille, à l'école ou au travail — devient impossible. Cependant, la vérité n'est pas un devoir absolu dans sa forme : une vérité brutale peut blesser inutilement. Il faut donc distinguer le \emph{fond} (ne jamais mentir sur l'essentiel) et la \emph{forme} (choisir le moment et les mots). La sagesse consiste à allier l'honnêteté au respect d'autrui.
\textbf{Points de vigilance :} la problématisation vient toujours en premier ; la synthèse ne doit pas être un simple compromis mou, mais une position \emph{argumentée}.
\end{corrige}

\begin{corrige}[Exercice 7 — Philosophie ou science ?]
\begin{enumerate}
\item \textbf{Science.} La composition chimique de l'eau se détermine par l'expérience et l'analyse en laboratoire ; la réponse est mesurable et vérifiable.
\item \textbf{Philosophie.} La liberté n'est pas un objet mesurable : c'est une question de sens et de valeur qui se discute par l'argumentation (déterminisme vs libre arbitre).
\item \textbf{Science.} La distance Yaoundé--Bafoussam (environ \SI{290}{km} par la route) se mesure ; la réponse est un fait chiffré, indiscutable.
\item \textbf{Philosophie.} La justice est une valeur : aucune expérience ne peut mesurer ce qui est \og juste \fg{} ; la question relève de la réflexion morale et politique.
\item \textbf{Philosophie.} Le bonheur et le sens du progrès sont des notions évaluatives : la technique peut mesurer ses effets matériels, mais non décider si ces effets rendent \og plus heureux \fg{}.
\end{enumerate}
\textbf{Critère général à retenir :} la science répond à des questions de \emph{fait} (comment ? combien ?) par l'expérience ; la philosophie répond à des questions de \emph{sens} et de \emph{valeur} (pourquoi ? faut-il ?) par la réflexion.
\end{corrige}

\begin{corrige}[Exercice 8 — L'utilité de la philosophie au quotidien]
\begin{enumerate}
\item \textbf{Argument 1 — L'éthique professionnelle :} le technicien ne fait pas qu'exécuter ; il doit décider \emph{bien}. La philosophie l'aide à distinguer le juste de l'injuste dans son travail. \emph{Exemple :} un technicien du bâtiment à Bafoussam refuse de signer le contrôle d'un immeuble construit avec du ciment de mauvaise qualité, malgré la pression de l'entrepreneur : il sait que sa responsabilité engage la sécurité des habitants.

\textbf{Argument 2 — L'esprit critique et la prise de décision :} la philosophie apprend à ne pas croire sans examiner, à peser les arguments avant de trancher. \emph{Exemple :} face à une rumeur circulant sur WhatsApp à propos d'un vaccin, un élève technicien vérifie les sources avant de la partager ; de même, en classe, il sait justifier sa réponse au Probatoire au lieu de réciter sans comprendre.
\end{enumerate}
\textbf{Points de vigilance :} chaque argument doit être accompagné d'un exemple \emph{concret et camerounais} ; un argument abstrait sans illustration perd la moitié de sa valeur.
\end{corrige}

\begin{corrige}[Exercice 9 — Dissertation : La philosophie est-elle utile à la vie de l'homme ?]
\textbf{1. Problématique.} Le sujet oppose l'apparente inutilité matérielle de la philosophie à sa nécessité pour la pensée et l'action. Question principale : \emph{la philosophie, qui ne produit rien de matériel, apporte-t-elle néanmoins à l'homme un bien indispensable ?}

\textbf{2. Plan dialectique.}
\begin{itemize}
\item \textbf{I. Thèse :} la philosophie semble inutile. Elle ne produit ni biens, ni techniques, ni profits ; elle tourne en rond sur des questions sans réponse définitive ; le temps qu'elle prend serait mieux employé aux études techniques ou au travail.
\item \textbf{II. Antithèse :} la philosophie est indispensable. Elle apprend à penser par soi-même (Socrate : l'examen de soi) ; elle distingue le vrai du faux, le juste de l'injuste ; elle donne un sens à l'existence et apaise l'âme (Épicure).
\item \textbf{III. Synthèse :} l'utilité de la philosophie n'est pas matérielle mais humaine : elle ne fabrique pas des objets, elle forme des hommes libres et responsables. C'est une utilité supérieure, car aucune technique ne dit \emph{comment} il faut s'en servir.
\end{itemize}

\textbf{3. Introduction rédigée.} \og À quoi ça sert, la philosophie ? \fg{} : telle est la question que se posent beaucoup d'élèves, pressés d'apprendre un métier utile. Or, par \emph{philosophie}, il faut entendre une réflexion rationnelle et universelle qui interroge le sens de l'existence, du monde et de l'action humaine ; par \emph{utile}, ce qui sert effectivement la vie de l'homme. Le problème est donc le suivant : une discipline qui ne produit ni biens ni techniques peut-elle être tenue pour utile, ou son utilité est-elle d'un autre ordre ? Nous examinerons d'abord les raisons de soupçonner la philosophie d'inutilité, puis nous montrerons qu'elle est indispensable à la pensée et à l'action, avant de soutenir que son utilité, non matérielle mais humaine, est la plus haute de toutes.

\textbf{4. Développement de la première partie (Thèse).} La philosophie paraît d'abord inutile, car elle ne produit rien de tangible. L'agronome augmente les récoltes de maïs dans la région de l'Ouest, l'informaticien crée des applications, le mécanicien répare les véhicules : leurs savoirs se mesurent en résultats. La philosophie, elle, discute sans fin : Platon, Kant ou Épicure se contredisent, et aucune de leurs questions ne reçoit de réponse définitive et vérifiable. Dans un pays où les besoins matériels sont urgents, un jeune Camerounais peut juger plus rentable d'apprendre la soudure ou la comptabilité que de s'interroger sur la justice. Exemple concret : au marché Mokolo, le commerçant qui calcule ses bénéfices obtient un gain immédiat, tandis que celui qui médite sur le sens de la richesse ne gagne, semble-t-il, rien. L'utilité se mesurerait donc au profit, et la philosophie échouerait à cette épreuve. Mais faut-il réduire l'utile au profitable ? C'est ce qu'il faut examiner.

\textbf{Barème indicatif (sur 20) :} problématique claire (2 pts) ; plan dialectique cohérent (4 pts) ; introduction complète en quatre moments (6 pts) ; développement de la thèse avec exemple camerounais (6 pts) ; correction de la langue (2 pts).
\textbf{Points de vigilance :} l'introduction doit contenir les quatre moments (accroche, définitions, problématique, annonce du plan) ; la thèse doit être défendue \emph{sincèrement}, même si on la réfutera ensuite ; interdiction d'annoncer sa position dès l'introduction.
\end{corrige}

\begin{corrige}[Exercice 10 — Commentaire de texte : Kant, \emph{Logique}]
\textbf{1. Présentation.} Emmanuel Kant (1724--1804) est un philosophe allemand, figure des Lumières. Ce texte, extrait de sa \emph{Logique}, s'inscrit dans le thème de la leçon : la définition de la philosophie par son objet.

\textbf{2. Thèse du texte.} Kant affirme que toute la philosophie se ramène à trois questions fondamentales — savoir, agir, espérer — qui définissent son objet : la philosophie n'est pas un savoir parmi d'autres, mais l'interrogation totale de la raison humaine sur elle-même.

\textbf{3. Explication des trois questions avec exemples.}
\begin{itemize}
\item \og \textbf{Que puis-je savoir ?} \fg{} interroge les limites et la validité de la connaissance. \emph{Exemple :} un élève de Première technique se demande si les informations trouvées sur Internet sont vraies et vérifiables.
\item \og \textbf{Que dois-je faire ?} \fg{} concerne l'action morale : quels sont mes devoirs ? \emph{Exemple :} dois-je dénoncer un camarade qui triche au Probatoire, ou me taire par solidarité ?
\item \og \textbf{Que m'est-il permis d'espérer ?} \fg{} porte sur la félicité et le sens de la vie. \emph{Exemple :} un jeune diplômé de Douala espère que son travail honnête conduira à une vie heureuse et digne.
\end{itemize}

\textbf{4. Les trois dimensions de la philosophie.} La première question recouvre la \textbf{théorie de la connaissance} (le vrai) ; la deuxième, la \textbf{morale} (le bien, le devoir) ; la troisième, la dimension de l'\textbf{espérance} (le bonheur, le sens ultime). Ainsi la philosophie embrasse tout l'homme : sa pensée, son action et son destin.

\textbf{5. Discussion (cinq à dix lignes).} Ces trois questions définissent bien le cœur de la philosophie, car elles couvrent le savoir, l'action et le sens. Kant lui-même ajoute d'ailleurs une quatrième question qui les résume : \og Qu'est-ce que l'homme ? \fg{}. Cependant, on peut objecter que certaines questions philosophiques — Qu'est-ce que la beauté ? Qu'est-ce qu'une société juste ? — débordent ce cadre. En réalité, elles s'y rattachent : la justice relève du \og Que dois-je faire ? \fg{} collectif. Les questions de Kant ne ferment donc pas la philosophie : elles en tracent les axes essentiels, que chaque génération remplit à sa manière.

\textbf{Barème indicatif (sur 20) :} présentation de l'auteur et du contexte (2 pts) ; thèse dégagée (4 pts) ; explication des trois questions avec exemples (6 pts) ; mise en relation avec les trois dimensions (4 pts) ; discussion argumentée (4 pts).
\textbf{Points de vigilance :} ne pas paraphraser le texte, mais l'\emph{expliquer} ; chaque exemple doit être précis et personnel ; la discussion exige une position justifiée, pas un simple oui ou non.
\end{corrige}

\begin{corrige}[Exercice 11 — Question de réflexion : La science peut-elle répondre à toutes les questions de l'homme ?]
\textbf{1. Différences entre science et philosophie.}
\begin{center}
\begin{tabularx}{\linewidth}{|l|X|X|}
\hline
\rowcolor{popblueL} \textbf{Critère} & \textbf{Science} & \textbf{Philosophie} \\
\hline
Objet & Faits mesurables et délimités (matière, vivant, nombres) & Sens, valeurs, fondements (liberté, justice, vérité) \\
\hline
Méthode & Observation, expérience, calcul, vérification & Réflexion, analyse des concepts, argumentation, discussion \\
\hline
Type de réponse & Certaine, provisoire, vérifiable par tous & Discutée, justifiée rationnellement, jamais définitivement close \\
\hline
\end{tabularx}
\end{center}

\textbf{2. Ce que la science explique (deux exemples).} La science explique les mécanismes du monde : elle identifie le virus responsable d'une épidémie et met au point un vaccin (médecine) ; elle prévoit les pluies et guide les semis des agriculteurs de l'Adamaoua (météorologie). Ses réponses sont mesurables et efficaces.

\textbf{3. Ce qui échappe à la science (deux exemples).} Aucune expérience ne peut dire si l'euthanasie est moralement acceptable : c'est une question de \emph{valeur}, non de fait. De même, la science peut décrire les inégalités économiques au Cameroun, mais elle ne peut pas décider ce qu'est une société \emph{juste} : cela relève de la réflexion philosophique sur le sens et les valeurs.

\textbf{4. Texte argumenté (quinze à vingt lignes).}

\emph{Introduction.} Le premier ami affirme que la science, parce qu'elle explique de plus en plus de phénomènes, rend la philosophie superflue. Mais expliquer le monde suffit-il à répondre à toutes les questions de l'homme ? Nous soutiendrons, avec le second ami, que la science et la philosophie sont complémentaires et non rivales.

\emph{Argument 1.} La science répond au \emph{comment}, non au \emph{pourquoi}. Elle dit comment fonctionne le cerveau, mais non pourquoi il faut vivre ; elle fabrique des techniques puissantes, mais ne dit pas si leur usage est bon. La bombe ou l'intelligence artificielle sont des produits de la science : seule la réflexion morale peut en juger l'emploi.

\emph{Argument 2.} Les questions les plus humaines — la liberté, la justice, le bonheur — ne sont pas mesurables. On ne pèse pas la dignité d'un homme comme on pèse un sac de cacao. Or ce sont ces questions qui orientent nos choix quotidiens : tricher ou non, respecter ou non autrui, servir ou non la communauté.

\emph{Conclusion.} Le second ami a raison : la science donne des moyens, la philosophie interroge les fins. Loin de s'exclure, elles se complètent : une société qui aurait la science sans la philosophie serait puissante mais aveugle. L'homme a besoin des deux ailes pour voler : le savoir et la sagesse.

\textbf{Barème indicatif (sur 20) :} tableau des différences exact (4 pts) ; exemples scientifiques précis (4 pts) ; exemples philosophiques précis (4 pts) ; texte structuré (introduction, deux arguments, conclusion) (6 pts) ; clarté de l'expression (2 pts).
\textbf{Points de vigilance :} ne pas opposer science et philosophie de façon caricaturale (la science n'est pas \og mauvaise \fg{}) ; chaque exemple doit illustrer nettement le critère (fait vs valeur) ; le texte final doit respecter la structure annoncée et la longueur demandée.
\end{corrige}

\newpage

\coursec{Fiche bilan}

\begin{popfigure}
\begin{center}\tcbox[colback=popblue,colframe=popblue,coltext=white,arc=9pt,boxsep=4pt,left=6pt,right=6pt]{\bfseries\small La philosophie}\end{center}
\vspace{-2pt}
\begin{tcbraster}[raster columns=3,raster equal height=rows,raster column skip=6pt,raster row skip=6pt,enhanced,blankest]
\begin{tcolorbox}[enhanced,colback=white,colframe=poporange,boxrule=0.9pt,arc=3pt,title={Définitions},fonttitle=\bfseries\scriptsize,coltitle=white,colbacktitle=poporange,left=4pt,right=4pt,top=2pt,bottom=2pt,boxsep=1pt,toptitle=1.5pt,bottomtitle=1.5pt,halign=left]
\begin{itemize}[nosep,leftmargin=*,itemsep=1pt,font=\scriptsize]
\item Philo-sophia : amour de la sagesse
\item Réflexion critique et rationnelle
\end{itemize}
\end{tcolorbox}
\begin{tcolorbox}[enhanced,colback=white,colframe=popgreen,boxrule=0.9pt,arc=3pt,title={Origine et objet},fonttitle=\bfseries\scriptsize,coltitle=white,colbacktitle=popgreen,left=4pt,right=4pt,top=2pt,bottom=2pt,boxsep=1pt,toptitle=1.5pt,bottomtitle=1.5pt,halign=left]
\begin{itemize}[nosep,leftmargin=*,itemsep=1pt,font=\scriptsize]
\item Grèce, VI\textsuperscript{e} s. av. J.-C.
\item Étonnement, doute
\item Objet : le tout, l'homme, les valeurs
\end{itemize}
\end{tcolorbox}
\begin{tcolorbox}[enhanced,colback=white,colframe=popgold,boxrule=0.9pt,arc=3pt,title={Méthode},fonttitle=\bfseries\scriptsize,coltitle=white,colbacktitle=popgold,left=4pt,right=4pt,top=2pt,bottom=2pt,boxsep=1pt,toptitle=1.5pt,bottomtitle=1.5pt,halign=left]
\begin{itemize}[nosep,leftmargin=*,itemsep=1pt,font=\scriptsize]
\item Doute, questionnement
\item Analyse, argumentation
\item Dialogue, rigueur
\end{itemize}
\end{tcolorbox}
\begin{tcolorbox}[enhanced,colback=white,colframe=poppink,boxrule=0.9pt,arc=3pt,title={Importance},fonttitle=\bfseries\scriptsize,coltitle=white,colbacktitle=poppink,left=4pt,right=4pt,top=2pt,bottom=2pt,boxsep=1pt,toptitle=1.5pt,bottomtitle=1.5pt,halign=left]
\begin{itemize}[nosep,leftmargin=*,itemsep=1pt,font=\scriptsize]
\item Esprit critique
\item Liberté, autonomie
\item Vie morale et citoyenne
\end{itemize}
\end{tcolorbox}
\begin{tcolorbox}[enhanced,colback=white,colframe=poppurple,boxrule=0.9pt,arc=3pt,title={Philosophie et science},fonttitle=\bfseries\scriptsize,coltitle=white,colbacktitle=poppurple,left=4pt,right=4pt,top=2pt,bottom=2pt,boxsep=1pt,toptitle=1.5pt,bottomtitle=1.5pt,halign=left]
\begin{itemize}[nosep,leftmargin=*,itemsep=1pt,font=\scriptsize]
\item Points communs : raison, vérité
\item Différences : expérience vs réflexion
\end{itemize}
\end{tcolorbox}
\end{tcbraster}

\legende{Carte mentale de la leçon : les cinq savoirs essentiels autour de la question \og Qu'est-ce que la philosophie ? \fg{}}
\end{popfigure}

\begin{bilan}
\textbf{Définitions clés à connaître par c\oe ur}
\begin{itemize}
\item \cle{Philosophie} : du grec \emph{philein} (aimer) et \emph{sophia} (sagesse, savoir) : littéralement \og amour de la sagesse \fg{}. Terme attribué à \textbf{Pythagore}.
\item \cle{Philosophie} (sens actuel) : réflexion rationnelle, critique et universelle qui interroge le sens de l'existence, la connaissance et les valeurs (le vrai, le bien, le beau).
\item \cle{Étonnement} (\emph{thaumazein}) : selon \textbf{Platon} et \textbf{Aristote}, sentiment de surprise devant le monde, origine du philosopher.
\item \cle{Doute} : attitude méthodique qui consiste à suspendre son jugement pour examiner les opinions (Socrate, Descartes).
\item \cle{Argumentation} : ensemble ordonné de raisons visant à justifier ou à réfuter une thèse.
\item \cle{Sagesse} : savoir parfait sur les choses humaines et divines, que le philosophe recherche sans prétendre le posséder.
\end{itemize}

\textbf{Repères historiques et auteurs}

\begin{center}
\begin{tabularx}{\linewidth}{|l|X|}\hline
\rowcolor{popblueL} \textbf{Repère} & \textbf{À retenir} \\ \hline
VI\textsuperscript{e} s. av. J.-C. & Naissance de la philosophie en Grèce (Milet) : passage du \emph{mythos} (récit) au \emph{logos} (raison). \\ \hline
Socrate & Maïeutique : faire accoucher les esprits par le dialogue ; \og Connais-toi toi-même \fg{}. \\ \hline
Platon / Aristote & L'étonnement est le commencement de la philosophie. \\ \hline
Descartes & Le doute méthodique comme première règle de la méthode. \\ \hline
Kant & La philosophie répond à : que puis-je savoir ? que dois-je faire ? que m'est-il permis d'espérer ? \\ \hline
\end{tabularx}
\end{center}

\textbf{Méthode express de la philosophie}
\begin{enumerate}
\item S'étonner et \textbf{problématiser} : transformer une opinion en question.
\item \textbf{Douter} : examiner et critiquer les préjugés et les évidences.
\item \textbf{Analyser} : définir les mots, distinguer les notions.
\item \textbf{Argumenter} : avancer des raisons, des exemples, des objections.
\item \textbf{Conclure} : formuler une réponse justifiée et nuancée.
\end{enumerate}

\textbf{Philosophie et science : l'essentiel}
\begin{itemize}
\item \textbf{Points communs} : recherche de la vérité, usage de la raison, refus de l'opinion et de la superstition.
\item \textbf{Différences} : la science procède par \emph{expérience} et \emph{vérification} sur des objets particuliers ; la philosophie procède par \emph{réflexion} et \emph{questionnement} sur le tout (sens, valeurs, fondements).
\item \textbf{Complémentarité} : la philosophie interroge les fondements et les conséquences morales des sciences ; les sciences nourrissent la réflexion philosophique.
\end{itemize}

\textbf{Importance de la philosophie}
\begin{itemize}
\item Formation de l'\cle{esprit critique} et de l'autonomie de jugement.
\item Préparation à la vie \cle{citoyenne} et morale (débat, tolérance, responsabilité).
\item Outil de réussite scolaire et professionnelle : rigueur, clarté, argumentation.
\end{itemize}
\end{bilan}

\begin{aretenir}[Checklist avant le Bac]
\begin{itemize}
\item[$\square$] Je sais définir la philosophie par son étymologie (\emph{philein} + \emph{sophia}) et citer Pythagore.
\item[$\square$] Je sais donner au moins deux définitions actuelles de la philosophie (réflexion critique, rationnelle, universelle).
\item[$\square$] Je sais situer l'origine de la philosophie (Grèce, VI\textsuperscript{e} s. av. J.-C.) et expliquer le passage du mythe au logos.
\item[$\square$] Je sais expliquer le rôle de l'étonnement (Platon, Aristote) et du doute (Socrate, Descartes).
\item[$\square$] Je sais nommer l'objet de la philosophie : le monde, l'homme, la connaissance, les valeurs.
\item[$\square$] Je sais décrire les étapes de la méthode philosophique : problématiser, douter, analyser, argumenter, conclure.
\item[$\square$] Je sais définir la maïeutique de Socrate et donner un exemple de dialogue.
\item[$\square$] Je sais citer au moins trois importances de la philosophie (esprit critique, liberté, vie citoyenne).
\item[$\square$] Je sais comparer philosophie et science : un point commun, une différence, une complémentarité.
\item[$\square$] Je sais illustrer chaque notion par un exemple concret de la vie courante camerounaise.
\item[$\square$] Je sais rédiger une introduction de dissertation : accroche, problématique, annonce du plan.
\item[$\square$] Je sais justifier une conclusion en m'appuyant sur des arguments et des auteurs cités correctement.
\end{itemize}
\end{aretenir}

\findecours

\end{document}
